% Panafricanisme et développement de l'Afrique — Philosophie Tle Tech — Cours signé OnBuch+
% Programme officiel MINESEC. Compiler avec : tectonic N09-panafricanisme-developpement-afrique.tex
\newcommand{\DOCMATIERE}{Philosophie}
\newcommand{\DOCNIVEAU}{Tle Tech}
\newcommand{\DOCMODULE}{Philosophie – classe de Terminale technique}
\newcommand{\DOCLECON}{N09}
\newcommand{\DOCTITRE}{Panafricanisme et développement de l'Afrique}
% OnBuch+ — préambule « cours signés » ENRICHI (compilé avec tectonic / XeLaTeX)
% Système visuel « Pop » d'OnBuch+ : fond crème (jamais blanc), bordures encre
% épaisses, ombres « dures » décalées, titres Archivo Black en capitales.
% Version MATHÉMATIQUES : graphes (pgfplots/tikz), boîtes pédagogiques
% (définition, propriété, méthode, attention, le savais-tu, exercice résolu,
% exercice, TP, bilan), page de garde et signature OnBuch+ ancrée en bas.
%
% Macros attendues dans main.tex AVANT \input{preamble} :
%   \DOCMATIERE  \DOCNIVEAU  \DOCMODULE  \DOCLECON (numéro)  \DOCTITRE
\documentclass[11pt]{article}

\usepackage{fontspec}
\usepackage[french]{babel}
\usepackage[a4paper,margin=2cm,top=3.4cm,bottom=2.8cm,headheight=1.4cm,footskip=1.3cm]{geometry}
\usepackage{amsmath,amssymb,amsfonts}
\usepackage{mathtools} % \xmapsto, \coloneqq... souvent utilisés par les modèles
\usepackage{cancel} % simplifications barrées (bilans de forces)
% \abs{z} (module, valeur absolue) et \norm{u} (norme), utilisés par les modèles
\providecommand{\abs}{}\let\abs\relax\DeclarePairedDelimiter{\abs}{\lvert}{\rvert}
\providecommand{\norm}{}\let\norm\relax\DeclarePairedDelimiter{\norm}{\lVert}{\rVert}
\usepackage[table]{xcolor} % [table] : fournit aussi \rowcolors (alternance de lignes)
\usepackage{pagecolor}
\usepackage{tikz}
\usetikzlibrary{arrows.meta,positioning,calc,shapes.geometric,shapes.misc,shapes.arrows,decorations.pathmorphing,decorations.pathreplacing,decorations.markings,patterns,shadows,mindmap,trees,fit,backgrounds,matrix,intersections,angles,quotes}
\usepackage{pgfplots}
\usepackage[version=4]{mhchem} % équations nucléaires \ce{^{235}_{92}U}
\usepackage[european,siunitx]{circuitikz} % schémas électriques
\pgfplotsset{compat=1.18}
\usepgfplotslibrary{fillbetween,groupplots} % aires entre courbes (calcul intégral)
\usepackage{enumitem}
\usepackage{fancyhdr}
% [most] charge toutes les bibliothèques tcolorbox (listings, minted, raster,
% documentation...) et gonfle inutilement la mémoire TeX sur un document long
% (~300 boîtes) : on ne charge que ce qui est réellement utilisé.
\usepackage[skins,breakable]{tcolorbox}
\usepackage{graphicx}
\usepackage{colortbl}
\usepackage{booktabs}
\usepackage{tabularx}
\usepackage{multirow}
\usepackage{diagbox} % case d'en-tête diagonale des tableaux à double entrée
% Colonnes extensibles centrée / gauche / droite (C, L, R), souvent utilisées par les modèles
\newcolumntype{C}{>{\centering\arraybackslash}X}
\newcolumntype{L}{>{\raggedright\arraybackslash}X}
\newcolumntype{R}{>{\raggedleft\arraybackslash}X}
\usepackage{array}
\usepackage{multicol}
\usepackage{siunitx}
% parse-numbers=false : accepte aussi \SI{2,5 \times 10^{-3}}{...}
\sisetup{locale=FR,output-decimal-marker={,},group-minimum-digits=4,parse-numbers=false}
\DeclareSIUnit\ohm{\ensuremath{\symup{\Omega}}} % U+2126 absent de la police
\DeclareSIUnit\division{div} % réglages d'oscilloscope (ms/div, V/div)
\DeclareSIUnit\tr{tr} % tours (vitesse de rotation, ex. \SI{1500}{\tr\per\minute})
\DeclareSIUnit\molar{M} % concentration molaire (ex. \SI{3}{\molar}, \SI{1}{\micro\molar})
\DeclareSIUnit\an{an} % année (le modèle confond parfois avec \year, un compteur TeX) — ex. \SI{5}{\kilo\meter\per\an}
\DeclareSIUnit\FCFA{FCFA} % franc CFA (ex. \SI{500}{\FCFA\per\kilo\gram})
\DeclareSIUnit\degreeBrix{\textdegree Brix} % teneur en sucre d'un sirop/jus (réfractométrie)
\DeclareSIUnit\month{mois} % le modèle utilise parfois \month, un compteur TeX, comme unité de durée
\DeclareSIUnit\microliter{\micro\liter} % le modèle écrit parfois \microliter en un seul token
\DeclareSIUnit\million{millions} % dénombrement (ex. \SI{15}{\million\per\mL} spermatozoïdes)
\usepackage{qrcode}
% \degree : commande fréquemment utilisée par les modèles pour un angle ou une
% température en dehors de \SI{}{} (non fournie par les paquets chargés).
\providecommand{\degree}{^{\circ}}
% \qed (fin de démonstration), utilisé par les modèles sans amsthm
\providecommand{\qed}{\hfill$\square$}
% \barre : double barre des valeurs interdites dans un tableau de variations (tkz-tab)
\providecommand{\barre}{\ensuremath{\|}}
% environnement proof (amsthm non chargé) utilisé par les modèles
\ifdefined\proof\else\newenvironment{proof}[1][Démonstration]{\par\noindent\textit{#1.}\ }{\qed\par}\fi
\usepackage{lastpage}
\usepackage{hyperref}
\hypersetup{hidelinks}

% ============================================================
% Palette « Pop » OnBuch+ — mêmes hex que la classe P (plus_ui.dart)
% ============================================================
\definecolor{poporange}{HTML}{F59321}
\definecolor{popdark}{HTML}{E07A0C}
\definecolor{popcream}{HTML}{FFF6E8}
\definecolor{popink}{HTML}{1C1714}
\definecolor{poppurple}{HTML}{7A5AE0}
\definecolor{popgreen}{HTML}{2E9E6A}
\definecolor{poppink}{HTML}{E0457B}
\definecolor{popblue}{HTML}{2F7DE1}
\definecolor{popgold}{HTML}{C98A12}
\definecolor{popmuted}{HTML}{8A7F76}
\definecolor{popline}{HTML}{EADFD2}
\definecolor{popgray}{HTML}{8A7F76}
\colorlet{popgrey}{popgray}
% Alias tolérants (le modèle nomme parfois les couleurs différemment)
\colorlet{popwhite}{white}
\colorlet{popblack}{popink}
\colorlet{popred}{poppink}
\colorlet{popyellow}{popgold}
% Teintes claires (fonds de tableaux/encadrés) — le modèle nomme parfois
% les couleurs avec un suffixe L (ex. popblueL) sans les avoir définies.
\colorlet{poporangeL}{poporange!20}
\colorlet{popdarkL}{popdark!20}
\colorlet{poppurpleL}{poppurple!20}
\colorlet{popgreenL}{popgreen!20}
\colorlet{poppinkL}{poppink!20}
\colorlet{popblueL}{popblue!20}
\colorlet{popgoldL}{popgold!20}
\colorlet{popredL}{popred!20}
\colorlet{popgrayL}{popgray!20}
% Teintes claires pour fonds de tableaux / zones
\colorlet{popblueL}{popblue!10!white}
\colorlet{poporangeL}{poporange!14!white}
\colorlet{popgreenL}{popgreen!12!white}
\colorlet{poppurpleL}{poppurple!12!white}
\colorlet{poppinkL}{poppink!12!white}
\colorlet{popgoldL}{popgold!14!white}

\pagecolor{popcream}

% ============================================================
% Typographie
% ============================================================
\defaultfontfeatures{Ligatures=TeX}
\setmainfont{PlusJakartaSans-Medium.ttf}[Path=../../../fonts/,BoldFont=PlusJakartaSans-Bold.ttf]
% Maths en sans-serif (Fira Math) avec chiffres et lettres droites de Plus
% Jakarta, pour que les formules se fondent dans le texte.
\usepackage[math-style=ISO,bold-style=ISO]{unicode-math}
\setmathfont{FiraMath-Regular.otf}[Path=../../../fonts/]
\setmathfont{PlusJakartaSans-Medium.ttf}[Path=../../../fonts/,range={up/num,up/{Latin,latin},\mathup}]
\usepackage{icomma} % virgule décimale sans espace en mode math
% Caractères mathématiques Unicode absents des polices de texte (Plus Jakarta,
% Archivo Black) : rendus par leur équivalent mathématique, en texte comme en
% formule — sinon ils s'affichent en carré vide (ex. « Arithmétique dans ℤ »).
\usepackage{newunicodechar}
\newunicodechar{ℝ}{\ensuremath{\mathbb{R}}}
\newunicodechar{ℤ}{\ensuremath{\mathbb{Z}}}
\newunicodechar{ℂ}{\ensuremath{\mathbb{C}}}
\newunicodechar{ℕ}{\ensuremath{\mathbb{N}}}
\newunicodechar{ℚ}{\ensuremath{\mathbb{Q}}}
\newunicodechar{𝔻}{\ensuremath{\mathbb{D}}}
\newunicodechar{α}{\ensuremath{\alpha}}
\newunicodechar{β}{\ensuremath{\beta}}
\newunicodechar{γ}{\ensuremath{\gamma}}
\newunicodechar{δ}{\ensuremath{\delta}}
\newunicodechar{ε}{\ensuremath{\varepsilon}}
\newunicodechar{θ}{\ensuremath{\theta}}
\newunicodechar{λ}{\ensuremath{\lambda}}
\newunicodechar{μ}{\ensuremath{\mu}}
\newunicodechar{σ}{\ensuremath{\sigma}}
\newunicodechar{τ}{\ensuremath{\tau}}
\newunicodechar{φ}{\ensuremath{\varphi}}
\newunicodechar{ψ}{\ensuremath{\psi}}
\newunicodechar{ω}{\ensuremath{\omega}}
\newunicodechar{Σ}{\ensuremath{\Sigma}}
% majuscules produites par \MakeUppercase dans les titres (α-aminés -> Α)
\newunicodechar{Α}{\ensuremath{\alpha}}
\newunicodechar{Β}{\ensuremath{\beta}}
\newunicodechar{Γ}{\ensuremath{\Gamma}}
\newunicodechar{Δ}{\ensuremath{\Delta}}
\newunicodechar{Ε}{\ensuremath{\varepsilon}}
\newunicodechar{Θ}{\ensuremath{\Theta}}
\newunicodechar{Λ}{\ensuremath{\Lambda}}
\newunicodechar{Μ}{\ensuremath{\mu}}
\newunicodechar{Τ}{\ensuremath{\tau}}
\newunicodechar{Φ}{\ensuremath{\Phi}}
\newunicodechar{Ψ}{\ensuremath{\Psi}}
\newunicodechar{Ω}{\ensuremath{\Omega}}
\newunicodechar{↦}{\ensuremath{\mapsto}}
\newunicodechar{∈}{\ensuremath{\in}}
\newunicodechar{∉}{\ensuremath{\notin}}
\newunicodechar{∧}{\ensuremath{\wedge}}
\newunicodechar{⇔}{\ensuremath{\Leftrightarrow}}
\newunicodechar{⟺}{\ensuremath{\Longleftrightarrow}}
\newunicodechar{⇒}{\ensuremath{\Rightarrow}}
\newunicodechar{⟹}{\ensuremath{\Longrightarrow}}
\newunicodechar{∘}{\ensuremath{\circ}}
\newunicodechar{∪}{\ensuremath{\cup}}
\newunicodechar{∩}{\ensuremath{\cap}}
\newunicodechar{≡}{\ensuremath{\equiv}}
\newunicodechar{⊂}{\ensuremath{\subset}}
\newunicodechar{⊥}{\ensuremath{\perp}}
\newunicodechar{∃}{\ensuremath{\exists}}
\newunicodechar{∀}{\ensuremath{\forall}}
\newunicodechar{∛}{\ensuremath{\sqrt[3]{\,}}}
\newunicodechar{∜}{\ensuremath{\sqrt[4]{\,}}}
\newunicodechar{⃗}{\ensuremath{{}^{\to}}}
\newunicodechar{ⁿ}{\ifmmode^{n}\else\textsuperscript{\ensuremath{n}}\fi}
\newunicodechar{ˣ}{\ifmmode^{x}\else\textsuperscript{\ensuremath{x}}\fi}
\newunicodechar{ᵇ}{\ifmmode^{b}\else\textsuperscript{\ensuremath{b}}\fi}
\newunicodechar{ʳ}{\ifmmode^{r}\else\textsuperscript{\ensuremath{r}}\fi}
\newunicodechar{ᵘ}{\ifmmode^{u}\else\textsuperscript{\ensuremath{u}}\fi}
\newunicodechar{ʸ}{\ifmmode^{y}\else\textsuperscript{\ensuremath{y}}\fi}
\newunicodechar{ᵃ}{\ifmmode^{a}\else\textsuperscript{\ensuremath{a}}\fi}
\newunicodechar{ᵉ}{\ifmmode^{e}\else\textsuperscript{\ensuremath{e}}\fi}
\newunicodechar{ᵏ}{\ifmmode^{k}\else\textsuperscript{\ensuremath{k}}\fi}
\newunicodechar{ᵖ}{\ifmmode^{p}\else\textsuperscript{\ensuremath{p}}\fi}
\newunicodechar{ᴺ}{\ifmmode^{N}\else\textsuperscript{\ensuremath{N}}\fi}
\newunicodechar{ᶜ}{\ifmmode^{c}\else\textsuperscript{\ensuremath{c}}\fi}
\newunicodechar{ʰ}{\ifmmode^{h}\else\textsuperscript{\ensuremath{h}}\fi}
\newunicodechar{ᵠ}{\ifmmode^{q}\else\textsuperscript{\ensuremath{q}}\fi}
\newunicodechar{ₙ}{\ifmmode_{n}\else\textsubscript{\ensuremath{n}}\fi}
\newunicodechar{ₐ}{\ifmmode_{a}\else\textsubscript{\ensuremath{a}}\fi}
\newunicodechar{ᵢ}{\ifmmode_{i}\else\textsubscript{\ensuremath{i}}\fi}
\newunicodechar{ᵣ}{\ifmmode_{r}\else\textsubscript{\ensuremath{r}}\fi}
\newunicodechar{ₖ}{\ifmmode_{k}\else\textsubscript{\ensuremath{k}}\fi}
\newunicodechar{ᵦ}{\ifmmode_{\beta}\else\textsubscript{\ensuremath{\beta}}\fi}
\newunicodechar{ₓ}{\ifmmode_{x}\else\textsubscript{\ensuremath{x}}\fi}
\newfontfamily\popdisplay{ArchivoBlack-Regular.ttf}[Path=../../../fonts/]
\newfontfamily\poplogo{SpaceGrotesk-Bold.ttf}[Path=../../../fonts/]
\newfontfamily\popsemi{PlusJakartaSans-SemiBold.ttf}[Path=../../../fonts/]
\newfontfamily\popxbold{PlusJakartaSans-ExtraBold.ttf}[Path=../../../fonts/]
\newfontfamily\popxboldls{PlusJakartaSans-ExtraBold.ttf}[Path=../../../fonts/,LetterSpace=3.0]

\setlist[enumerate]{leftmargin=1.7em,itemsep=5pt,topsep=4pt}
\setlist[itemize]{leftmargin=1.7em,itemsep=4pt,topsep=4pt,label={\color{poporange}\textbullet}}
\setlength{\parindent}{0pt}
\setlength{\parskip}{5pt}
\renewcommand{\arraystretch}{1.25}

% pgfplots : style Pop par défaut
\pgfplotsset{
  every axis/.append style={axis line style={popink,line width=1pt},tick style={popink},
    grid style={popline,line width=0.5pt},label style={font=\small},tick label style={font=\footnotesize}},
  popcurve/.style={poporange,line width=1.8pt,no markers,smooth},
}
% Même style disponible dans un \draw TikZ simple (hors pgfplots)
\tikzset{popcurve/.style={poporange,line width=1.8pt}}
\tikzset{popgrid/.style={popline,very thin}}
\colorlet{popcurve}{poporange} % le nom du style est parfois utilisé comme couleur

% ============================================================
% Pill / squiggle
% ============================================================
\newcommand{\poppill}[3]{%
  \tikz[baseline=-0.55ex]\node[draw=popink,line width=1.2pt,rounded corners=2.6pt,%
    fill=#2,inner xsep=6.5pt,inner ysep=3pt]{%
    \popxboldls\fontsize{7}{7}\selectfont\color{#3}\MakeUppercase{#1}};%
}
\newcommand{\popsquiggle}[1]{%
  \tikz[baseline=-0.4ex]{
    \draw[#1,line width=2.6pt,line cap=round]
      plot [smooth] coordinates {(0,0.09) (0.34,0.01) (0.68,0.21) (1.02,0.01) (1.36,0.10)};
  }%
}
% Mot-clé surligné (marqueur orange sous le texte)
\newcommand{\cle}[1]{{\popxbold\color{popdark}#1}}

% ============================================================
% En-tête / pied
% ============================================================
\pagestyle{fancy}
\fancyhf{}
\renewcommand{\headrulewidth}{0pt}
\renewcommand{\footrulewidth}{0pt}
\lhead{\raisebox{-0.15ex}{\poplogo\fontsize{13}{13}\selectfont\color{popink}OnBuch\color{poporange}+}}
\rhead{\poppill{\DOCMATIERE\ \textbullet\ \DOCNIVEAU\ \textbullet\ Leçon \DOCLECON}{white}{popink}}
\lfoot{\popxbold\fontsize{7.6}{7.6}\selectfont\color{popmuted}\MakeUppercase{Cours signé OnBuch+}\ \textbullet\ {\popsemi\DOCTITRE}}
\rfoot{\tikz[baseline=-0.6ex]\node[rounded corners=8.5pt,fill=popink,minimum height=17pt,minimum width=17pt,inner xsep=5pt,inner sep=0]{\popxbold\fontsize{8}{8}\selectfont\color{popcream}\thepage\,/\,\pageref*{LastPage}};}

% ============================================================
% Titres
% ============================================================
\newcommand{\chaptitre}[2]{%
  \begin{tcolorbox}[enhanced,colback=poporange,colframe=popink,boxrule=2.6pt,arc=11pt,
      drop shadow=popink,left=15pt,right=15pt,top=12pt,bottom=11pt,before skip=3pt,after skip=18pt]
    \poppill{#2}{popink}{popcream}\par\vspace{9pt}
    {\raggedright\hyphenpenalty=10000\exhyphenpenalty=10000\popdisplay\fontsize{22}{24}\selectfont\color{popcream}\MakeUppercase{#1}\par}
  \end{tcolorbox}
}
% Section numérotée : pastille orange avec le numéro + titre Archivo
\newcounter{popsec}
\newcommand{\coursec}[1]{%
  \stepcounter{popsec}\setcounter{popsub}{0}%
  \par\vspace{16pt}\noindent
  \begin{minipage}{\linewidth}
  \tikz[baseline=-0.7ex]\node[draw=popink,line width=1.6pt,fill=poporange,rounded corners=4pt,minimum size=22pt,inner sep=0]{\popdisplay\fontsize{12}{12}\selectfont\color{popcream}\thepopsec};%
  \hspace{8pt}{\popdisplay\fontsize{15}{17}\selectfont\color{popink}\MakeUppercase{#1}}\par
  \vspace{4pt}\noindent{\color{poporange}\rule{36pt}{3.6pt}}
  \end{minipage}\par\nopagebreak\vspace{8pt}
}
\newcounter{popsub}
\newcommand{\courssub}[1]{%
  \stepcounter{popsub}%
  \par\vspace{9pt}\noindent{\popxbold\fontsize{11.5}{13}\selectfont\color{popink}{\color{poporange}\thepopsec.\thepopsub}\hspace{6pt}#1}\par\nopagebreak\vspace{4pt}
}

% ============================================================
% Boîtes pédagogiques — même recette : fond blanc, bordure épaisse colorée,
% ombre dure encre, badge plein. Titre optionnel [#1].
% ============================================================
\tcbset{popbase/.style={breakable,enhanced,colback=white,boxrule=2.3pt,arc=9pt,
  shadow={3pt}{-3pt}{0pt}{popink},left=14pt,right=14pt,top=11pt,bottom=11pt,before skip=11pt,after skip=15pt,
  fontupper=\popsemi\fontsize{10.6}{14.5}\selectfont\color{popink},
  fontlower=\popsemi\fontsize{10.6}{14.5}\selectfont\color{popink}}}
% Coche dessinée (le glyphe ✓ n'existe pas dans les polices du document)
\newcommand{\popcheck}[1]{\tikz[baseline=-0.5ex]\draw[#1,line width=1.6pt,line cap=round,line join=round](-0.13,0.0)--(-0.04,-0.09)--(0.13,0.1);}
\providecommand{\coche}{\popcheck{popgreen}}
\providecommand{\resultat}[1]{\par\smallskip\noindent\fcolorbox{popgreen}{popgreenL}{\parbox{\dimexpr\linewidth-2\fboxsep-2\fboxrule\relax}{\textbf{Résultat.} #1}}\par\smallskip}
\newcommand{\popboxhead}[3]{\poppill{#1}{#2}{white}\ifx\relax#3\relax\else\hspace{6pt}{\popxbold\fontsize{10.6}{12}\selectfont\color{popink}#3}\fi\par\vspace{7pt}}

\newtcolorbox{aretenir}[1][]{popbase,colframe=popgreen,before upper={\popboxhead{À retenir}{popgreen}{#1}}}
\newtcolorbox{exemplebox}[1][]{popbase,colframe=poporange,before upper={\popboxhead{Exemple}{poporange}{#1}}}
\newtcolorbox{definition}[1][]{popbase,colframe=popblue,before upper={\popboxhead{Définition}{popblue}{#1}}}
\newtcolorbox{propriete}[1][]{popbase,colframe=poppurple,before upper={\popboxhead{Propriété}{poppurple}{#1}}}
\newtcolorbox{methode}[1][]{popbase,colframe=popgold,before upper={\popboxhead{Méthode}{popgold}{#1}}}
\newtcolorbox{attention}[1][]{popbase,colframe=poppink,before upper={\popboxhead{Attention}{poppink}{#1}}}
\newtcolorbox{experience}[1][]{popbase,colframe=popblue,colback=popblueL,before upper={\popboxhead{Expérience}{popblue}{#1}}}
% « Le savais-tu ? » : carte violette pleine (contexte, culture, Cameroun)
\newtcolorbox{savaistu}[1][]{popbase,colback=poppurple,colframe=popink,
  fontupper=\popsemi\fontsize{10.4}{14.2}\selectfont\color{white},
  before upper={\poppill{Le savais-tu\,?}{white}{poppurple}\ifx\relax#1\relax\else\hspace{6pt}{\popxbold\color{white}#1}\fi\par\vspace{7pt}}}
% Exercice résolu : énoncé en haut, \tcblower puis la solution sur fond crème
\newcounter{popexo}\newcounter{popexr}
\newtcolorbox{exoresolu}[1][]{popbase,colframe=poporange,colbacklower=popcream,
  segmentation style={popink,line width=1.2pt,dashed},
  before upper={\stepcounter{popexr}\popboxhead{Exercice résolu \thepopexr}{poporange}{#1}},
  before lower={\poppill{Solution}{popink}{popcream}\par\vspace{6pt}}}
% Exercice (non résolu en place) — la correction est dans la partie Corrigés
\newtcolorbox{exercice}[1][]{popbase,colframe=popink,
  before upper={\stepcounter{popexo}\popboxhead{Exercice \thepopexo}{popink}{#1}}}
% Corrigé
\newtcolorbox{corrige}[1][]{popbase,colframe=popgreen,colback=popgreenL,
  before upper={\popboxhead{Corrigé}{popgreen}{#1}}}
% Objectifs / prérequis / bilan
\newtcolorbox{objectifs}{popbase,colframe=popink,colback=poporangeL,
  before upper={\popboxhead{Objectifs de la leçon}{popink}{}}}
\newtcolorbox{prerequis}{popbase,colframe=popink,colback=popblueL,
  before upper={\popboxhead{Prérequis}{popblue}{}}}
\newtcolorbox{bilan}{popbase,colframe=popink,colback=white,boxrule=2.8pt,
  before upper={\popboxhead{Fiche bilan}{poporange}{L'essentiel à connaître pour l'examen}}}
% Figure encadrée avec légende
\newtcolorbox{popfigure}[1][]{enhanced,breakable=false,colback=white,colframe=popline,boxrule=1.4pt,arc=8pt,
  left=10pt,right=10pt,top=10pt,bottom=8pt,before skip=10pt,after skip=12pt,halign=center,
  lower separated=false,halign lower=center,
  fontlower=\popsemi\fontsize{9}{11.5}\selectfont\color{popmuted},
  #1}
\newcommand{\legende}[1]{\par\vspace{6pt}{\popsemi\fontsize{9}{11.5}\selectfont\color{popmuted}#1}}

% ============================================================
% Signature OnBuch+
% ============================================================
\newcommand{\poplogoinline}[1]{{\poplogo\fontsize{#1}{#1}\selectfont\color{popink}OnBuch\color{poporange}+}}
\newcommand{\signatureonbuch}{%
  \begin{tcolorbox}[enhanced,colback=popink,colframe=popink,boxrule=2.2pt,arc=10pt,
      drop shadow=poporange,left=14pt,right=14pt,top=10pt,bottom=10pt,before skip=0pt,after skip=0pt]
    \tikz[baseline=-0.7ex]\node[circle,fill=popgreen,draw=popcream,line width=1.3pt,minimum size=22pt,inner sep=0]{\popcheck{white}};%
    \hspace{9pt}\begin{minipage}[c]{0.62\linewidth}
      {\popxbold\fontsize{12}{13}\selectfont\color{popcream}Signé\ \textbullet\ {\poplogo OnBuch\color{poporange}+}}\par\vspace{2pt}
      {\raggedright\popsemi\fontsize{8}{10.5}\selectfont\color{popline}\MakeUppercase{Équipe pédagogique OnBuch+}\\\MakeUppercase{Conforme au programme officiel MINESEC}\par}
    \end{minipage}\hfill
    {\poplogo\fontsize{22}{22}\selectfont\color{popcream}OnBuch\color{poporange}+}
  \end{tcolorbox}%
}

% ============================================================
% Page de garde
% #1 titre · #2 matière · #3 niveau · #4 module · #5 n° leçon · #6 durée ·
% #7 description / objectifs (texte ou itemize)
% ============================================================
\newcommand{\pagegarde}[7]{%
  \thispagestyle{empty}
  \newgeometry{a4paper,margin=1.7cm,top=1.6cm,bottom=1.4cm}%
  \begin{tikzpicture}[remember picture,overlay]
    \fill[poporange!18] ([xshift=-3.2cm,yshift=-3.6cm]current page.north east) circle (4.6cm);
    \fill[poppurple!14] ([xshift=2.4cm,yshift=7.6cm]current page.south west) circle (3.8cm);
  \end{tikzpicture}%
  {\poplogo\fontsize{30}{30}\selectfont\color{popink}OnBuch\color{poporange}+}\hfill
  \raisebox{0.6ex}{\poppill{Cours signé \textbullet\ Premium}{popink}{popcream}}\par
  \vspace{1pt}\popsquiggle{poppurple}\par
  \vspace{8pt}
  \begin{tcolorbox}[enhanced,colback=poporange,colframe=popink,boxrule=2.8pt,arc=13pt,
      drop shadow=popink,left=18pt,right=18pt,top=12pt,bottom=13pt,after skip=10pt]
    \poppill{#2 \textbullet\ #3}{popink}{popcream}\hspace{5pt}\poppill{#4}{white}{popink}\par\vspace{8pt}
    {\popdisplay\fontsize{14}{14}\selectfont\color{popink}LEÇON #5}\par\vspace{4pt}
    {\raggedright\hyphenpenalty=10000\exhyphenpenalty=10000\popdisplay\fontsize{25}{27}\selectfont\color{popcream}\MakeUppercase{#1}\par}
    \vspace{8pt}
    {\popxbold\fontsize{9.5}{9.5}\selectfont\color{popink}\MakeUppercase{Durée conseillée : #6}}
  \end{tcolorbox}
  \begin{tcolorbox}[enhanced,colback=white,colframe=popink,boxrule=2.2pt,arc=10pt,
      drop shadow=popink,left=16pt,right=16pt,top=10pt,bottom=10pt,after skip=8pt]
    \poppill{Dans cette leçon}{poppurple}{white}\par\vspace{5pt}
    \popsemi\fontsize{9.3}{12.6}\selectfont\color{popink}#7
  \end{tcolorbox}
  \noindent\begin{minipage}[t]{0.487\linewidth}
    \begin{tcolorbox}[enhanced,colback=popgreen,colframe=popink,boxrule=2.2pt,arc=10pt,
        drop shadow=popink,left=11pt,right=11pt,top=10pt,bottom=10pt,height=3.1cm,valign=center]
      \tcbox[colback=white,colframe=popink,boxrule=1.3pt,arc=5pt,boxsep=0pt,nobeforeafter,
        left=3pt,right=3pt,top=3pt,bottom=3pt,valign=center]{\qrcode[height=1.9cm]{https://play.google.com/store/apps/details?id=com.luvvix.onbuch&hl=fr}}%
      \hspace{8pt}\begin{minipage}[c]{0.5\linewidth}
        {\popdisplay\fontsize{11}{12.5}\selectfont\color{white}\MakeUppercase{Télécharge OnBuch}}\par\vspace{4pt}
        {\raggedright\popsemi\fontsize{8.4}{11}\selectfont\color{white}Gratuit \textbullet\ Google Play\\Tuteur IA, annales, concours, résultats\par}
      \end{minipage}
    \end{tcolorbox}
  \end{minipage}\hfill
  \begin{minipage}[t]{0.487\linewidth}
    \begin{tcolorbox}[enhanced,colback=poppurple,colframe=popink,boxrule=2.2pt,arc=10pt,
        drop shadow=popink,left=11pt,right=11pt,top=10pt,bottom=10pt,height=3.1cm,valign=center]
      \tikz[baseline=-0.7ex]\node[circle,fill=white,draw=popink,line width=1.3pt,minimum size=1.6cm,inner sep=0]{\popdisplay\fontsize{24}{24}\selectfont\color{poppurple}+};%
      \hspace{8pt}\begin{minipage}[c]{0.6\linewidth}
        {\popdisplay\fontsize{11}{12.5}\selectfont\color{white}\MakeUppercase{Passe à OnBuch+}}\par\vspace{4pt}
        {\raggedright\popsemi\fontsize{8.4}{11}\selectfont\color{white}Tous les cours signés, Tuteur IA illimité, fiches PDF — onglet OnBuch+ de l'app\par}
      \end{minipage}
    \end{tcolorbox}
  \end{minipage}
  \vspace{8pt}
  \popcontenu
  \vfill
  \signatureonbuch
  \restoregeometry
  \newpage
}

% Rangée « Ce cours contient » (page de garde)
\newcommand{\poptile}[3]{% #1 couleur #2 gros texte #3 légende
  \begin{tcolorbox}[enhanced,colback=#1,colframe=popink,boxrule=2pt,arc=9pt,shadow={2.5pt}{-2.5pt}{0pt}{popink},
      left=6pt,right=6pt,top=9pt,bottom=9pt,halign=center,valign=center,height=1.75cm,width=\linewidth,nobeforeafter]
    {\popdisplay\fontsize{12.5}{13}\selectfont\color{white}\MakeUppercase{#2}}\par\vspace{3pt}
    {\centering\popsemi\fontsize{7.8}{9.5}\selectfont\color{white}#3\par}
  \end{tcolorbox}}
\newcommand{\popcontenu}{%
  \par\noindent{\popxboldls\fontsize{7.5}{7.5}\selectfont\color{popmuted}CE COURS SIGNÉ CONTIENT}\par\vspace{7pt}
  \noindent\begin{minipage}[t]{0.235\linewidth}\poptile{popblue}{Cours}{complet, illustré, pas à pas}\end{minipage}\hfill
  \begin{minipage}[t]{0.235\linewidth}\poptile{poporange}{Méthodes}{et exercices résolus}\end{minipage}\hfill
  \begin{minipage}[t]{0.235\linewidth}\poptile{poppink}{Exercices}{type examen + corrigés}\end{minipage}\hfill
  \begin{minipage}[t]{0.235\linewidth}\poptile{popgreen}{Fiche bilan}{l'essentiel à retenir}\end{minipage}\par}

% Sommaire
\newtcolorbox{sommaire}{enhanced,colback=white,colframe=popink,boxrule=2.2pt,arc=10pt,
  drop shadow=popink,left=18pt,right=18pt,top=13pt,bottom=13pt,before skip=6pt,after skip=20pt,
  before upper={\poppill{Sommaire}{poppurple}{white}\par\vspace{9pt}\popsemi\fontsize{10.6}{14.5}\selectfont\color{popink}}}

% Signature de fin de cours — poussée tout en bas de la dernière page
\newcommand{\findecours}{%
  \par\vfill\nopagebreak
  \begin{center}\popsquiggle{poporange}\end{center}\vspace{4pt}
  \signatureonbuch
}
\AtBeginDocument{\def\llbracket{\mathopen{[\mkern-3.5mu[}}\def\rrbracket{\mathclose{]\mkern-3.5mu]}}} % intervalles d'entiers (glyphe ⟦ absent de la police)
% Glyphes absents de Fira Math : redéfinis à partir de glyphes disponibles
\AtBeginDocument{%
  \def\setminus{\mathbin{\backslash}}\let\smallsetminus\setminus
  \def\vdots{\mathord{\vbox{\baselineskip3.2pt\lineskiplimit0pt\kern4pt\hbox{.}\hbox{.}\hbox{.}}}}%
  \def\ddots{\mathinner{\mkern1mu\raise6.5pt\hbox{.}\mkern2mu\raise3.6pt\hbox{.}\mkern2mu\raise0.7pt\hbox{.}\mkern1mu}}%
  \def\triangle{\mathord{\scalebox{0.9}{$\symup{\Delta}$}}}%
  \def\nabla{\mathord{\raisebox{\depth}{\scalebox{1}[-1]{$\symup{\Delta}$}}}}%
  \def\checkmark{\popcheck{popdark}}%
  \def\star{\mathbin{\ast}}\def\bigstar{\mathord{\ast}}%
  \def\diamond{\mathbin{\scalebox{0.6}{$\Diamond$}}}%
  \def\bigcup{\mathop{\mathchoice{\raisebox{-0.3em}{\scalebox{1.7}{$\cup$}}}{\scalebox{1.25}{$\cup$}}{\cup}{\cup}}\displaylimits}%
  \def\bigcap{\mathop{\mathchoice{\raisebox{-0.3em}{\scalebox{1.7}{$\cap$}}}{\scalebox{1.25}{$\cap$}}{\cap}{\cap}}\displaylimits}%
  \def\lesseqqgtr{\mathrel{\lessgtr}}\def\bigoplus{\mathop{\oplus}\displaylimits}%
}
\providecommand{\ssi}{si et seulement si} % abréviation fréquente
\AtBeginDocument{\providecommand{\wideparen}{\overparen}} % arc de cercle

%% FIGFIX-BEGIN v5 — correctifs globaux des figures (docs/onbuch-generation/tools/figfix)
\usepackage{adjustbox}\usepackage{etoolbox}\tcbuselibrary{raster}
% toute figure TikZ/pgfplots plus large que la ligne est réduite à la largeur disponible
\BeforeBeginEnvironment{tikzpicture}{\begin{adjustbox}{max width=\linewidth}}
\AfterEndEnvironment{tikzpicture}{\end{adjustbox}}
% légendes pgfplots lisibles par-dessus les courbes
\pgfplotsset{every axis/.append style={legend style={fill=white,fill opacity=0.92,text opacity=1,draw=black!15,inner sep=3pt}}}
% étiquettes d'arêtes (syntaxe edge["..."]) avec fond blanc
\tikzset{every edge quotes/.append style={fill=white,inner sep=1.2pt}}
%% FIGFIX-END
\begin{document}

\pagegarde{Panafricanisme et développement de l'Afrique}{Philosophie}{Tle Tech}{Philosophie – classe de Terminale technique}{N09}{2 séances (fiche de progression harmonisée)}{Cette leçon interroge les notions de développement et de sous-développement, puis présente le panafricanisme — son histoire, ses objectifs et ses difficultés — comme philosophie du développement de l'Afrique. Elle invite l'élève à appliquer ces notions à des situations concrètes camerounaises et africaines et à construire une argumentation philosophique rigoureuse.
\vspace{4pt}
\begin{multicols}{2}\small
\begin{itemize}
\item À la fin de cette leçon, je sais définir les concepts de développement et de sous-développement
\item Je sais distinguer croissance économique et développement humain à partir d'indicateurs concrets (PIB, IDH)
\item Je sais retracer les grandes étapes historiques du panafricanisme (congrès, OUA, UA)
\item Je sais identifier les objectifs politiques, économiques et culturels du panafricanisme
\item Je sais analyser les difficultés internes et externes qui freinent l'unité africaine
\item Je sais appliquer les notions de développement et d'unité africaine à des situations concrètes du Cameroun et de la vie courante
\end{itemize}\end{multicols}}

\begin{sommaire}
\begin{multicols}{2}
\begin{enumerate}
\item Le développement : définition et indicateurs
\item Le sous-développement : réalité et causes
\item Historique du panafricanisme
\item Objectifs du panafricanisme comme philosophie du développement
\item Difficultés et limites du panafricanisme
\item Synthèse : l'unité africaine est-elle la voie du développement ?
\item Activité d'intégration
\item Méthodes et exercices résolus
\item Exercices
\item Corrigés des exercices
\item Fiche bilan
\end{enumerate}
\end{multicols}
\end{sommaire}

\begin{objectifs}
\begin{itemize}
\item À la fin de cette leçon, je sais définir les concepts de développement et de sous-développement
\item Je sais distinguer croissance économique et développement humain à partir d'indicateurs concrets (PIB, IDH)
\item Je sais retracer les grandes étapes historiques du panafricanisme (congrès, OUA, UA)
\item Je sais identifier les objectifs politiques, économiques et culturels du panafricanisme
\item Je sais analyser les difficultés internes et externes qui freinent l'unité africaine
\item Je sais appliquer les notions de développement et d'unité africaine à des situations concrètes du Cameroun et de la vie courante
\item Je sais interpréter des documents (cartes, données chiffrées, photographies) relatifs au développement de l'Afrique
\item Je sais rédiger et justifier une démarche argumentative, une réponse ou une conclusion en dissertation philosophique
\end{itemize}
\end{objectifs}

\begin{prerequis}
\begin{itemize}
\item Notion de colonisation et de décolonisation de l'Afrique (histoire de 1ère)
\item Notion de liberté et de responsabilité (leçons de philosophie de Terminale)
\item Notion d'État et de souveraineté (leçon sur l'État)
\item Repères géographiques de base : grands ensembles régionaux africains (CEDEAO, CEMAC)
\item Notion de justice sociale et d'inégalités (début d'année de Terminale)
\end{itemize}
\end{prerequis}

\newpage

\coursec{Le développement : définition et indicateurs}

\courssub{Situation d'accroche et problématique}

En 2023, un jeune diplômé de Douala, assis dans un cybercafé du quartier Deïdo, fait un constat troublant en manipulant son smartphone. Ce téléphone, assemblé en Chine, contient du cobalt extrait du Congo-Kinshasa. Le cacao récolté dans la région du Centre au Cameroun est exporté brut vers l'Europe, puis réimporté sous forme de tablettes de chocolat vendues dix fois plus cher dans les supermarchés de Yaoundé. Et lorsqu'il veut se rendre à Lagos, capitale économique du Nigeria voisin, il découvre que les vols les plus accessibles passent par Paris ou Addis-Abeba : pour relier deux métropoles africaines, il faut sortir de l'Afrique. Comment un continent si riche de ses ressources (pétrole, cobalt, cacao, bois, uranium, terres arables) peut-il rester si pauvre de ses revenus ? Pourquoi, soixante ans après les indépendances, ce paradoxe persiste-t-il ?

Certains intellectuels et hommes d'État africains, de Kwame Nkrumah à Thabo Mbeki, répondent : parce que le continent est \cle{divisé}, fragmenté en plus de cinquante États souvent trop petits, trop faibles et tournés vers leurs anciennes métropoles coloniales plutôt que les uns vers les autres. Dès lors, l'\cle{unité africaine} prônée par le \cle{panafricanisme} serait-elle la condition du développement de l'Afrique ? Mais avant de répondre, il faut d'abord clarifier les mots : que signifie exactement \og développement \fg{} ? Comment le mesure-t-on ? Et qu'est-ce qui distingue un pays riche d'un pays développé ? Telle est la problématique de cette première section : \emph{le développement est-il une simple accumulation de richesses, ou un processus plus profond de transformation de la vie individuelle et collective ?}

\courssub{Étymologie et sens commun du mot « développement »}

Commençons, comme toujours en philosophie, par interroger le langage ordinaire. Le mot \og développement \fg{} vient du verbe \emph{développer}, formé du préfixe \emph{dé-} (qui marque l'ouverture, le défaire) et de \emph{velopper}, d'origine ancienne, qui signifie \emph{envelopper}, \emph{enrouler}. \cle{Développer}, c'est donc littéralement \og dérouler ce qui était enroulé \fg{}, \og ouvrir ce qui était fermé \fg{}. On parle ainsi du développement d'une feuille qui se déplie, du développement d'un bourgeon qui devient fleur, du développement de l'enfant qui grandit, ou encore, en photographie ancienne, du développement d'un cliché qui fait apparaître l'image cachée sur la pellicule.

Le sens commun dégage donc une intuition fondamentale : le développement est un \cle{déploiement}, un \cle{épanouissement}, le passage d'un état où des possibilités sont contenues, repliées, virtuelles, à un état où elles s'actualisent et deviennent visibles. Le germe contient déjà la plante, mais en puissance ; le développement est le processus par lequel cette puissance devient acte. On retrouve ici la grande distinction aristotélicienne entre la \cle{puissance} et l'\cle{acte} : développer, c'est faire passer ce qui existe en germe à son accomplissement.

Appliquée à une société, cette intuition signifie qu'un pays \og en voie de développement \fg{} n'est pas un pays sans richesses ni sans potentialités : c'est un pays dont les potentialités (ressources naturelles, jeunesse, savoir-faire, cultures) ne sont pas encore pleinement déployées au bénéfice de tous. Le jeune diplômé de Douala, le cacao du Centre, le cobalt congolais : autant de puissances qui attendent leur acte.

\courssub{Définition philosophique du développement}

\begin{definition}[Développement]
Le \cle{développement} est le processus par lequel une société accroît ses richesses matérielles tout en améliorant durablement les conditions de vie, les libertés et l'épanouissement de tous ses membres. Philosophiquement, il désigne le passage d'un état moins accompli à un état plus accompli de la vie individuelle et collective.
\end{definition}

Analysons cette définition. Elle comporte deux dimensions inséparables. D'une part, une dimension \emph{quantitative} et matérielle : produire davantage de biens et de services, construire des routes, des écoles, des hôpitaux, maîtriser les techniques. D'autre part, une dimension \emph{qualitative} et humaine : vivre plus longtemps et en meilleure santé, être instruit, participer aux décisions, jouir de libertés réelles, être traité avec dignité. Une société développée n'est pas seulement une société qui \emph{a} davantage ; c'est une société dont les membres \emph{sont} davantage : plus capables, plus libres, plus épanouis.

C'est pourquoi le développement est un concept \cle{global} et \cle{normatif} : il porte un jugement de valeur. Dire qu'une société se développe, ce n'est pas seulement constater qu'elle change, c'est affirmer qu'elle progresse vers un état jugé meilleur, plus accompli, plus conforme à la dignité humaine. Toute la difficulté — et tout l'enjeu philosophique de cette leçon — tient à cette question : qui définit ce qui est \og meilleur \fg{}, et selon quels critères ?

\courssub{Le développement économique : croissance, industrialisation, infrastructures}

La dimension la plus visible et la plus mesurable du développement est économique. On l'évalue traditionnellement par plusieurs indicateurs.

\begin{itemize}
\item Le \cle{Produit Intérieur Brut} (PIB) : valeur totale des biens et services produits par un pays en une année. Rapporté à la population, le \cle{PIB par habitant} donne une première idée du niveau de richesse moyen.
\item Le taux de \cle{croissance économique} : variation annuelle du PIB, qui indique le dynamisme de la production.
\item Le degré d'\cle{industrialisation} : part de l'industrie dans l'économie, capacité à transformer sur place les matières premières au lieu de les exporter brutes.
\item Les \cle{infrastructures} : routes, ports, chemins de fer, réseaux d'électricité et d'eau potable, télécommunications, qui sont les artères de la vie économique.
\end{itemize}

L'exemple du cacao camerounais illustre parfaitement l'enjeu de l'industrialisation : exporter la fève brute rapporte peu ; transformer la fève en chocolat sur place créerait des emplois, des compétences et une valeur ajoutée beaucoup plus grande. Le développement économique suppose donc le passage d'une économie de simple extraction et exportation à une économie de \cle{transformation}.

\courssub{Le développement humain : santé, éducation, libertés, dignité}

Mais la richesse produite profite-t-elle à tous ? Un pays peut afficher un PIB élevé alors que la majorité de sa population vit dans la pauvreté, sans accès à l'eau potable ni aux soins. C'est pour dépasser la seule mesure économique que le Programme des Nations unies pour le développement (PNUD) a créé en 1990 l'\cle{Indice de Développement Humain} (IDH), qui combine trois dimensions :

\begin{enumerate}
\item la \cle{santé}, mesurée par l'espérance de vie à la naissance ;
\item l'\cle{éducation}, mesurée par la durée moyenne et la durée attendue de scolarisation ;
\item le \cle{niveau de vie}, mesuré par le revenu national brut par habitant.
\end{enumerate}

L'IDH varie entre 0 et 1 : plus il est proche de 1, plus le développement humain est élevé. Cet indice repose sur une philosophie précise, celle de l'économiste indien Amartya Sen : le développement doit être pensé en termes de \cle{capabilités} (\emph{capabilities}), c'est-à-dire de capacités réelles des personnes à mener la vie qu'elles ont des raisons de valoriser. Être développé, c'est pouvoir effectivement se nourrir, se soigner, s'instruire, circuler, participer à la vie publique. La richesse n'est qu'un moyen ; la fin, ce sont les libertés réelles des personnes.

\begin{savaistu}[Amartya Sen : le développement comme liberté]
Amartya Sen, économiste et philosophe indien, prix Nobel d'économie en 1998, définit le développement comme \og un processus d'élargissement des libertés réelles dont jouissent les gens \fg{} (\emph{Development as Freedom}, 1999). Dans cette perspective, la famine, l'analphabétisme ou l'absence de soins ne sont pas seulement des malheurs : ce sont des \emph{privations de liberté}. Un pays se développe donc lorsque ses citoyens deviennent capables de choisir et de réaliser leur propre vie. Cette approche a inspiré la création de l'IDH.
\end{savaistu}

\courssub{Distinction cruciale : croissance n'est pas développement}

Nous touchons ici au cœur conceptuel de la leçon. La \cle{croissance} est un phénomène \emph{quantitatif} : l'augmentation de la richesse produite. Le \cle{développement} est un phénomène \emph{qualitatif} et \emph{global} : la transformation des structures économiques, sociales, politiques et culturelles au bénéfice de tous. Il peut y avoir croissance sans développement : c'est le cas des pays pétroliers où l'exploitation du pétrole gonfle le PIB, enrichit une minorité, mais laisse la majorité sans écoles, sans hôpitaux et sans emplois, tout en rendant l'économie dépendante d'une seule ressource.

\begin{attention}[Erreur fréquente au Probatoire / Baccalauréat technique]
Ne confondez jamais \cle{croissance économique} et \cle{développement} dans une dissertation. La croissance est quantitative (combien produit-on ?) ; le développement est qualitatif et global (qui en bénéficie, et comment vit-on ?). Un sujet qui vous demande si la richesse suffit au développement attend précisément que vous mobilisiez cette distinction, avec l'exemple des pays pétroliers à fort PIB mais à fortes inégalités.
\end{attention}

\begin{exemplebox}[Le Gabon et l'île Maurice : la richesse ne suffit pas]
Le Gabon, riche de son pétrole, affiche un PIB par habitant nettement supérieur à celui de l'île Maurice (environ \num{8800} dollars contre environ \num{11500} dollars en 2022, l'île Maurice ayant dépassé le Gabon ces dernières années grâce à sa diversification). Pourtant, l'IDH de Maurice (environ \num{0,80}, le plus élevé d'Afrique subsaharienne avec les Seychelles) dépasse celui du Gabon (environ \num{0,71}). Pourquoi ? Parce qu'à Maurice, la richesse est mieux répartie, l'éducation et la santé sont accessibles à tous, et l'économie est diversifiée (textile, tourisme, services financiers, TIC), tandis qu'au Gabon la rente pétrolière profite inégalement à la population. Conclusion philosophique : le développement ne se mesure pas à la quantité de richesses, mais à leur conversion en capacités réelles pour tous.
\end{exemplebox}

\begin{popfigure}
\begin{center}
\begin{tikzpicture}
\begin{axis}[
ybar, bar width=0.45cm,
width=0.95\linewidth, height=6.5cm,
symbolic x coords={Cameroun,Gabon,Maurice,France,Rwanda},
xtick=data,
ymin=0, ymax=45000,
ylabel={PIB par habitant (dollars, environ 2022)},
legend style={at={(0.5,-0.22)},anchor=north,legend columns=2},
nodes near coords, every node near coord/.append style={font=\tiny,rotate=90,anchor=west},
]
\addplot[fill=popblue] coordinates {(Cameroun,1600) (Gabon,8800) (Maurice,11500) (France,41000) (Rwanda,970)};
\addplot[fill=poporange] coordinates {(Cameroun,5880) (Gabon,7100) (Maurice,8000) (France,9000) (Rwanda,5480)};
\legend{PIB/habitant, IDH $\times$ \num{10000}}
\end{axis}
\end{tikzpicture}
\end{center}
\legende{Comparaison du PIB par habitant et de l'IDH (multiplié par \num{10000} pour être lisible sur le même graphique) de cinq pays. On observe que le Gabon, malgré un PIB par habitant élevé, a un IDH inférieur à celui de Maurice : la richesse produite ne se convertit pas automatiquement en développement humain. (Valeurs approximatives 2022, sources : Banque mondiale et PNUD.)}
\end{popfigure}

\courssub{Les deux dimensions du développement intégral}

On peut désormais schématiser la structure complète du concept : le développement véritable, ou \cle{développement intégral}, résulte de la convergence de la croissance économique et du développement humain. L'une sans l'autre reste boiteuse : richesse sans partage (Gabon) ou idéal sans moyens.

\begin{popfigure}
\begin{center}
\begin{tikzpicture}[
node distance=1.6cm,
branche/.style={rectangle, rounded corners, draw=popblue, fill=popblueL, text width=4.6cm, align=center, font=\small},
sous/.style={rectangle, draw=popmuted, fill=popcream, text width=4.2cm, align=center, font=\footnotesize},
final/.style={rectangle, rounded corners, draw=poporange, fill=poporangeL, text width=5.5cm, align=center, font=\small\bfseries},
fleche/.style={-{Stealth[length=3mm]}, thick, popdark}
]
\node[branche, align=center] (croissance){\textbf{Croissance économique}\\ (quantitatif)};
\node[sous, below=0.5cm of croissance] (c1) {PIB, industrialisation, infrastructures, transformation des matières premières};
\node[branche, right=3.2cm of croissance, align=center] (humain){\textbf{Développement humain}\\ (qualitatif)};
\node[sous, below=0.5cm of humain] (h1) {Santé, éducation, libertés réelles, dignité (IDH, capabilités de Sen)};
\node[final, below=2.2cm of $(c1)!0.5!(h1)$, align=center] (integral){Développement intégral\\ épanouissement de tous les membres de la société};
\draw[fleche] (c1.south) -- (integral.north west);
\draw[fleche] (h1.south) -- (integral.north east);
\end{tikzpicture}
\end{center}
\legende{Schéma fonctionnel du développement intégral : la croissance économique (dimension quantitative) et le développement humain (dimension qualitative) doivent converger. Ni la richesse sans partage, ni les droits sans moyens ne suffisent.}
\end{popfigure}

\courssub{Le développement endogène : se développer par ses propres forces}

Reste une dernière question, décisive pour la suite de la leçon : le développement peut-il être importé ? L'histoire des politiques de développement en Afrique montre que les modèles copiés de l'extérieur (plans d'ajustement structurel des années 1980-1990 imposés par les institutions financières internationales, par exemple) ont souvent échoué, parce qu'ils ignoraient les réalités locales. D'où l'idée de \cle{développement endogène} : un développement \og venu de l'intérieur \fg{} (du grec \emph{endon}, dedans, et \emph{genos}, origine), qui s'appuie sur les propres forces, ressources, savoirs et valeurs de la société concernée.

Se développer de façon endogène, pour un pays africain, c'est par exemple valoriser les savoir-faire locaux (agriculture vivrière, artisanat, médecine traditionnelle à côté de la médecine moderne), transformer sur place les matières premières, financer le développement par l'épargne et la fiscalité nationales plutôt que par la seule aide extérieure, et définir ses propres priorités au lieu de recevoir des modèles tout faits. Le développement endogène ne signifie pas le repli sur soi : il signifie que le modèle de développement doit être \emph{approprié} par ceux qu'il concerne. C'est précisément l'une des ambitions du panafricanisme, que nous étudierons dans les sections suivantes : faire de l'Afrique unie l'acteur de son propre développement.

\courssub{Application au Cameroun : quelques données récentes}

Appliquons maintenant ces notions à une situation concrète, comme le demande le programme. Où en est le Cameroun ?

\begin{center}
\begin{tabularx}{\linewidth}{|l|X|}
\hline
\rowcolor{popblueL} \textbf{Indicateur} & \textbf{Valeur approximative (2021-2023)} \\
\hline
PIB par habitant & environ \num{1600} dollars \\
\hline
IDH (rang mondial) & environ \num{0,58} (environ 150ème rang sur 191 pays) \\
\hline
Espérance de vie & environ 60 ans \\
\hline
Taux brut de scolarisation au secondaire & environ 55 à 60 \% \\
\hline
Accès à l'électricité & environ 65 \% de la population (fort contraste ville/village) \\
\hline
Accès à l'eau potable & environ 70 à 75 \% de la population \\
\hline
\end{tabularx}
\end{center}

Que conclure de ces chiffres ? Le Cameroun connaît une croissance économique réelle (environ 3 à 4 \% par an), dispose de ressources considérables (pétrole, gaz, bois, cacao, café, aluminium, agriculture) et d'atouts humains précieux (une population jeune, un taux d'alphabétisation parmi les meilleurs d'Afrique centrale). Mais son IDH reste moyen : les richesses ne se convertissent pas encore pleinement en santé, en éducation et en emplois pour tous. Le Cameroun illustre donc à la fois la distinction croissance/développement et le défi du développement endogène : transformer sur place, investir dans le capital humain, réduire les inégalités entre régions et entre villes et campagnes.

\begin{exoresolu}[Croissance et développement au Cameroun]
\textbf{Énoncé.} Un élève affirme : \og Le PIB du Cameroun augmente chaque année, donc le Cameroun se développe. \fg{} Ce raisonnement est-il valide ? Justifiez votre réponse en mobilisant les notions du cours et un exemple.
\tcblower
\textbf{Solution détaillée.} Le raisonnement est incomplet, car il confond \cle{croissance} et \cle{développement}. Premièrement, la croissance est quantitative : elle mesure l'augmentation de la richesse produite. Or le développement est qualitatif et global : il suppose que cette richesse améliore effectivement la santé, l'éducation, les libertés et la dignité de tous. Deuxièmement, l'expérience le confirme : le Gabon a longtemps eu un PIB par habitant supérieur à celui de Maurice, tout en ayant un IDH inférieur, parce que sa rente pétrolière profitait inégalement à la population. Troisièmement, au Cameroun même, le PIB progresse alors qu'environ un habitant sur trois n'a pas accès à l'électricité et que l'IDH reste moyen. Conclusion : la croissance est une \emph{condition} du développement, mais non une condition \emph{suffisante}. Pour que le raisonnement de l'élève soit valide, il faudrait vérifier que la richesse produite se convertit en capabilités réelles pour l'ensemble de la population.
\end{exoresolu}

\begin{methode}[Définir un concept en dissertation]
Pour définir rigoureusement une notion philosophique en introduction ou au début d'un développement, procédez en quatre étapes :
\begin{enumerate}
\item \textbf{Sens commun et étymologie} : partez de l'usage ordinaire du mot (développer = dérouler, épanouir).
\item \textbf{Définition technique} : formulez la définition du cours (processus d'accroissement des richesses et d'amélioration des conditions de vie et des libertés de tous).
\item \textbf{Distinction avec les notions voisines} : opposez le concept à ce avec quoi on le confond (développement $\neq$ croissance ; développement exogène $\neq$ endogène).
\item \textbf{Exemple} : illustrez par un cas précis et chiffré (Gabon/Maurice, Cameroun).
\end{enumerate}
Cette démarche montre au correcteur que vous maîtrisez le concept au lieu de le réciter.
\end{methode}

\begin{aretenir}[L'essentiel de la section]
\begin{itemize}
\item Le \cle{développement} est le passage d'un état moins accompli à un état plus accompli de la vie individuelle et collective : accroissement des richesses \emph{et} amélioration des conditions de vie, des libertés et de la dignité de tous.
\item Il a deux dimensions : économique (PIB, industrialisation, infrastructures) et humaine (santé, éducation, libertés réelles — IDH, capabilités d'Amartya Sen).
\item \cle{Croissance} $\neq$ \cle{développement} : la première est quantitative, le second qualitatif et global (exemple du Gabon et de Maurice).
\item Le \cle{développement endogène} s'appuie sur les forces et les valeurs propres d'une société : c'est l'une des voies que le panafricanisme propose à l'Afrique.
\end{itemize}
\end{aretenir}

\begin{exercice}[Pour s'entraîner]
1. Définissez le développement et distinguez-le de la croissance en vous appuyant sur un exemple africain. 2. Pourquoi Amartya Sen définit-il le développement comme un élargissement des libertés ? 3. Le développement du Cameroun peut-il être endogène sans être fermé sur lui-même ? Rédigez un paragraphe argumenté d'une quinzaine de lignes.
\end{exercice}

\begin{corrige}[Exercice 1]
1. Le développement est le processus par lequel une société accroît ses richesses tout en améliorant les conditions de vie, les libertés et l'épanouissement de tous ; la croissance n'est que l'augmentation quantitative de la production. Exemple : le Gabon (fort PIB/habitant, IDH moyen) contre Maurice (PIB comparable, IDH élevé grâce à l'éducation, la santé et la diversification). 2. Pour Sen, la richesse n'est qu'un moyen : ce qui compte, ce sont les capabilités, c'est-à-dire les capacités réelles des personnes de vivre la vie qu'elles valorisent (se nourrir, se soigner, s'instruire, participer). La pauvreté est donc une privation de liberté. 3. Oui : le développement endogène signifie s'appuyer d'abord sur ses propres ressources, savoirs et priorités (transformation locale du cacao, valorisation des savoir-faire, financement national), non se couper du monde. Le Cameroun peut commercer et coopérer avec l'extérieur tout en définissant lui-même son modèle, notamment dans le cadre d'une intégration africaine. Le paragraphe doit articuler thèse, argument, exemple et conclusion.
\end{corrige}

\coursec{Le sous-développement : réalité et causes}

Dans la section précédente, nous avons défini le \cle{développement} comme un processus global de transformation économique, sociale et culturelle qui améliore durablement les conditions de vie des populations, et nous avons examiné ses indicateurs (PIB, IDH, espérance de vie, scolarisation). Mais ce constat s'impose alors : si le développement est possible, pourquoi une grande partie de l'Afrique demeure-t-elle pauvre ? Faut-il croire que l'Afrique est simplement « en retard » sur le chemin parcouru par l'Europe ? Ou bien sa pauvreté a-t-elle des causes précises, historiques et politiques, qu'il faut analyser ? C'est toute la question du \cle{sous-développement}, que nous abordons maintenant en trois temps : sa définition, ses manifestations, puis ses causes et le débat philosophique qu'elles soulèvent.

\courssub{Définition du sous-développement}

Commençons par une intuition simple. Imaginez un élève très intelligent, capable de devenir ingénieur, mais qu'on empêche d'aller à l'école et qu'on oblige à travailler gratuitement pour enrichir quelqu'un d'autre. Cet élève n'est pas « naturellement incapable » : ses potentialités sont \emph{bloquées} et \emph{exploitées au profit d'un autre}. C'est exactement l'idée contenue dans le concept de sous-développement appliqué à une société entière.

\begin{definition}[Sous-développement]
Le \cle{sous-développement} désigne la situation d'une société dont les richesses et les potentialités humaines sont bloquées ou exploitées au profit de l'extérieur, maintenant la population dans la \textbf{pauvreté}, la \textbf{dépendance} et l'impossibilité de transformer ses propres ressources. Il ne s'agit pas d'une absence de richesses, mais d'un \emph{blocage} : le pays possède des ressources (cacao, pétrole, minerais, terres fertiles, jeunesse nombreuse) sans pouvoir en tirer profit pour son propre peuple.
\end{definition}

Trois traits essentiels composent donc le sous-développement :
\begin{itemize}
\item la \textbf{pauvreté} : une grande partie de la population vit dans la précarité matérielle ;
\item la \textbf{dépendance} : l'économie du pays dépend des décisions prises ailleurs (cours mondiaux des matières premières, créanciers, firmes étrangères) ;
\item le \textbf{blocage des potentialités} : le pays ne transforme pas ses richesses, n'industrialise pas, n'utilise pas pleinement ses compétences humaines.
\end{itemize}

\begin{attention}[Erreur fréquente à éviter]
Ne présentez jamais le sous-développement comme un simple « \emph{retard} » que l'Afrique rattraperait en copiant l'Occident. Cette vision, dite \emph{évolutionniste}, suppose que tous les pays suivent le même chemin et que l'Afrique serait seulement « en arrière ». Or la thèse des théoriciens de la dépendance (dont André Gunder Frank) montre le contraire : le sous-développement de l'Afrique et le développement de l'Occident sont \textbf{deux faces d'un même processus historique}. L'Afrique n'est pas « en retard » : elle a été \emph{sous-développée} par des siècles d'exploitation. Dire « retard » au lieu de « processus produit » est une faute d'analyse fréquente au Baccalauréat technique.
\end{attention}

\begin{savaistu}[D'où vient l'expression « pays en voie de développement » ?]
Après la Seconde Guerre mondiale, on parlait ouvertement de « pays sous-développés ». Mais ce terme, jugé péjoratif et blessant, a progressivement été remplacé dans le langage diplomatique par « pays en voie de développement » (PVD), puis par « pays du Sud » ou « pays émergents » pour les plus dynamiques. Ce changement de vocabulaire n'est pas innocent : il traduit l'idée optimiste que tout pays est \emph{en marche} vers le développement. Mais les philosophes nous invitent à nous méfier des mots : changer le nom ne change pas la réalité. Un pays peut être « en voie de développement » depuis soixante ans sans jamais arriver !
\end{savaistu}

\courssub{Les manifestations du sous-développement}

Comment reconnaît-on concrètement une société sous-développée ? Par des signes observables et mesurables, que tout élève peut vérifier autour de lui :

\begin{itemize}
\item la \textbf{pauvreté monétaire} : une part importante de la population vit avec moins de quelques francs CFA par jour, sans accès stable au logement décent, à l'eau potable ou à l'électricité ;
\item la \textbf{malnutrition} : malgré des terres fertiles, des enfants souffrent de carences, car les meilleures terres servent parfois aux cultures d'exportation (cacao, café, coton) plutôt qu'à l'alimentation locale ;
\item la \textbf{faible espérance de vie} : les maladies pourtant soignables (paludisme, infections) restent meurtrières faute d'infrastructures sanitaires suffisantes ;
\item l'\textbf{analphabétisme} : une fraction notable de la population adulte ne sait ni lire ni écrire, ce qui limite sa participation à la vie économique et citoyenne ;
\item l'\textbf{exode rural} : les jeunes quittent les villages pour les villes (Douala, Yaoundé), où beaucoup survivent dans l'économie informelle ;
\item la \textbf{fuite des cerveaux} : médecins, ingénieurs et infirmiers formés en Afrique partent exercer en Europe ou en Amérique, privant le continent des compétences qu'il a payées pour former.
\end{itemize}

Ces manifestations forment un cercle vicieux : la pauvreté empêche l'éducation, le manque d'éducation entretient la pauvreté, et la fuite des compétences prive le pays des moyens de sortir de ce cercle.

\courssub{Le sous-développement : un processus historiquement produit}

Venons-en au cœur philosophique de la leçon. Le sous-développement est-il un état naturel, un point de départ que tous les peuples auraient connu ? Le sociologue et économiste \textbf{André Gunder Frank} répond non, dans une formule célèbre : il faut parler du « \emph{développement du sous-développement} ».

\begin{propriete}[Thèse d'André Gunder Frank]
Le sous-développement n'est pas un retard naturel : c'est un \textbf{processus historiquement produit}. Les mêmes mécanismes (traite négrière, colonisation, commerce inégal) qui ont \emph{développé} l'Europe ont simultanément \emph{sous-développé} l'Afrique. Centre et périphérie sont liés : la richesse de l'un s'est construite en partie sur l'appauvrissement de l'autre. Le sous-développement est donc une \emph{création de l'histoire}, non une fatalité de la nature.
\end{propriete}

Cette thèse renverse le regard : l'Afrique n'est pas un continent qui « n'a pas encore » commencé à se développer ; c'est un continent dont le développement a été \emph{orienté vers l'extérieur}. Ses ports, ses chemins de fer, ses mines ont été construits non pour relier les Africains entre eux, mais pour évacuer les richesses vers l'Europe. On comprend alors pourquoi le Cameroun exporte ses matières premières brutes et importe des produits finis : ce circuit n'est pas un accident, c'est un \emph{héritage organisé}.

\begin{popfigure}
\begin{center}
\begin{tikzpicture}[scale=1]
% Contour très simplifié de l'Afrique
\draw[thick, popdark, fill=popcream] (0,3.2) -- (1.2,3.8) -- (2.6,3.6) -- (3.2,2.8) -- (3.6,1.8) -- (3.1,0.8) -- (2.6,-0.4) -- (2.2,-1.6) -- (1.7,-2.2) -- (1.2,-1.4) -- (0.9,-0.2) -- (0.4,0.9) -- (-0.2,1.9) -- (-0.6,2.6) -- cycle;
\node[popdark, font=\bfseries] at (1.5,1.2) {AFRIQUE};
% Matières premières
\node[align=center, font=\small, popgreen!60!black] at (1.5,0.4) {Cacao, pétrole,\\café, bois, minerais\\(bruts)};
% Flèche sortante
\draw[-{Stealth[length=4mm]}, line width=2.5pt, poporange] (3.4,1.6) .. controls (5,2.4) .. (6.6,2.6);
\node[align=left, font=\small, poporange!80!black] at (6.9,3.3) {Exportation des\\matières premières\\brutes (prix imposés)};
% Flèche entrante
\draw[-{Stealth[length=4mm]}, line width=2.5pt, popblue] (6.6,-0.6) .. controls (5,-1.2) .. (3.4,-0.2);
\node[align=left, font=\small, popblue!80!black] at (6.9,-1.4) {Importation de produits\\manufacturés chers\\(chocolat, carburant, machines)};
% Europe schématique
\draw[thick, popmuted, fill=popblueL] (7.6,1.2) rectangle (9.4,2.2);
\node[font=\small\bfseries, popdark] at (8.5,1.7) {Occident};
% Nord
\draw[-{Stealth}, popdark] (-1,3.4) -- (-1,4.1);
\node[font=\small, popdark] at (-1,4.35) {N};
\end{tikzpicture}
\end{center}
\legende{Croquis schématique du « commerce inégal » : l'Afrique exporte ses matières premières brutes à bas prix (flèche orange) et importe des produits manufacturés coûteux (flèche bleue). La valeur ajoutée de la transformation est captée à l'extérieur du continent.}
\end{popfigure}

\courssub{Les causes externes du sous-développement}

Les premières causes sont extérieures au continent : elles relèvent de l'histoire des relations entre l'Afrique et le reste du monde.

\begin{enumerate}
\item \textbf{La traite négrière} (du XVI\textsuperscript{e} au XIX\textsuperscript{e} siècle) : des millions d'Africains jeunes et valides ont été déportés. Le continent a ainsi perdu une part considérable de sa force de travail, tandis que cette main-d'œuvre enrichissait les plantations d'Amérique et les ports d'Europe.
\item \textbf{La colonisation} (fin XIX\textsuperscript{e} -- 1960) : les puissances coloniales ont organisé des économies d'extraction. On produisait ce dont l'Europe avait besoin, pas ce dont les Africains avaient besoin. Les frontières, les infrastructures et les administrations ont été dessinées pour l'exploitation, non pour le développement.
\item \textbf{Le partage inégal du commerce mondial} : les pays africains vendent des matières premières brutes dont les prix sont fixés sur les bourses de Londres, Paris ou New York, et achètent des produits transformés dont les prix sont fixés par les industriels étrangers. C'est ce qu'on appelle la \emph{détérioration des termes de l'échange} : avec un sac de cacao, on achète chaque année moins de machines qu'avant.
\item \textbf{L'endettement} : les États africains empruntent pour financer leurs équipements ; le remboursement de la dette absorbe ensuite une part importante de leurs budgets, au détriment de la santé et de l'éducation.
\item \textbf{Les prix des matières premières imposés} : le producteur de cacao camerounais ne fixe pas le prix de son cacao ; il le subit. Quand le cours mondial chute, c'est tout le budget du pays qui vacille.
\end{enumerate}

\courssub{Les causes internes du sous-développement}

Il serait cependant malhonnête — et philosophiquement faible — de tout attribuer à l'extérieur. Le sous-développement a aussi des causes internes, que les Africains eux-mêmes peuvent combattre :

\begin{itemize}
\item la \textbf{mauvaise gouvernance} : mauvaise gestion des deniers publics, absence de planification à long terme, institutions faibles ;
\item la \textbf{corruption} : le détournement des ressources publiques prive l'État des moyens d'investir dans les routes, les écoles et les hôpitaux ;
\item les \textbf{conflits} : guerres civiles et instabilité politique détruisent les infrastructures et découragent l'investissement ;
\item la \textbf{faible industrialisation} : faute d'usines de transformation, le continent vend brut ce qu'il pourrait vendre transformé, donc beaucoup plus cher ;
\item la \textbf{dépendance alimentaire} : des pays aux terres fertiles importent du riz, du blé ou de l'huile, alors qu'ils pourraient nourrir leur population.
\end{itemize}

\begin{popfigure}
\begin{center}
\begin{tikzpicture}[scale=0.9]
% Barres horizontales
\def\bars{
  30/Commerce inégal (externe)/poporangeL,
  25/Colonisation et héritage (externe)/popblueL,
  20/Mauvaise gouvernance (interne)/poporange!70,
  15/Faible industrialisation (interne)/popblue!70,
  10/Conflits et instabilité (interne)/popgreenL}
\foreach \val/\label/\col [count=\i] in \bars {
  \draw[fill=\col, draw=popdark] (0,-\i*1.1) rectangle (\val*0.075, -\i*1.1+0.7);
  \node[left, align=right, font=\small] at (-0.3, -\i*1.1+0.35) {\label};
  \node[right, font=\small\bfseries] at (\val*0.075+0.15, -\i*1.1+0.35) {\val\%};
}
% Axe
\draw[->, thick] (0,0) -- (2.5,0) node[right, font=\small] {Poids estimé (\%)};
% Graduations
\foreach \x in {0,10,20,30} {
  \draw (\x*0.075,0.1) -- (\x*0.075,-0.1) node[below, font=\tiny] {\x};
}
% Légende interne/externe
\node[font=\small\bfseries, poporange] at (1.5, 0.8) {Causes externes : 55\%};
\node[font=\small\bfseries, popblue] at (1.5, -6.2) {Causes internes : 45\%};
\end{tikzpicture}
\end{center}
\legende{Répartition indicative des causes du sous-développement (en \%), à discuter en classe : faut-il accorder plus de poids aux causes externes (55 \%) ou aux causes internes (45 \%) ? Ce partage est un support de débat, non une vérité statistique.}
\end{popfigure}

\courssub{Situation concrète : le cas du Cameroun}

\begin{exemplebox}[Le cacao camerounais : une richesse qui s'évapore]
Le Cameroun est l'un des grands producteurs mondiaux de cacao. Pourtant, il \textbf{exporte plus de 70 \% de son cacao brut}, en fèves non transformées ; la transformation locale (beurre de cacao, poudre, chocolat) ne concerne qu'une faible part de la production. Conséquence : une tablette de chocolat vendue à Douala peut rapporter à l'industriel étranger plusieurs fois ce qu'a gagné le planteur camerounais. Le même raisonnement vaut pour le pétrole : le Cameroun exporte du pétrole brut et importe du carburant raffiné, plus cher. La \emph{valeur ajoutée} — c'est-à-dire la richesse créée par la transformation — part à l'étranger, avec les emplois qu'elle aurait pu créer.
\end{exemplebox}

\begin{methode}[Analyser un document chiffré en philosophie]
Face à un tableau ou une statistique au Baccalauréat technique, procédez en quatre étapes :
\begin{enumerate}
\item \textbf{Identifier la source} : qui a produit ce chiffre, quand, dans quel but ? (ex. : données sur les exportations de cacao, ministère du Commerce) ;
\item \textbf{Dégager la tendance} : le phénomène augmente-t-il, diminue-t-il, stagne-t-il ? (ex. : la part du cacao transformé localement reste faible) ;
\item \textbf{Comparer} : mettre en rapport deux données (ex. : plus de 70 \% exporté brut contre moins de 30 \% transformé sur place) ;
\item \textbf{Conclure en argumentant} : relier le chiffre à la notion du cours (ex. : cette disproportion illustre la dépendance et le blocage des potentialités, traits du sous-développement).
\end{enumerate}
Un chiffre seul ne prouve rien : c'est son interprétation argumentée qui a une valeur philosophique.
\end{methode}

\courssub{Débat : le sous-développement est-il une fatalité ?}

Si le sous-développement est un processus produit par l'histoire, alors il peut être \emph{défait} par l'action humaine : ce qui a été fait peut être défait. C'est le fondement de l'\textbf{optimisme du développement endogène} : un développement « venu de l'intérieur », appuyé sur les ressources, les savoirs et les initiatives propres aux Africains, plutôt que sur l'imitation servile des modèles étrangers. Des expériences locales — coopératives de planteurs transformant elles-mêmes leur cacao, transformation artisanale du manioc, technologies mobiles adaptées aux besoins — montrent que des voies africaines de développement existent.

Mais l'optimisme ne doit pas devenir naïveté : les causes externes (prix imposés, dette) ne disparaissent pas par la seule volonté. Le débat reste donc ouvert : le sous-développement est-il une fatalité ? La réponse philosophique que nous construirons est nuancée : ce n'est pas une fatalité \emph{naturelle}, puisqu'il est historiquement produit ; mais s'en libérer exige à la fois des réformes internes (gouvernance, industrialisation) et une transformation des rapports internationaux — ce qui annonce la question du panafricanisme, traitée dans les sections suivantes.

\begin{exoresolu}[Analyse d'une situation concrète]
\textbf{Énoncé.} Un élève déclare : « L'Afrique est pauvre parce qu'elle est naturellement en retard ; il suffit qu'elle copie l'Europe pour se développer. » En vous appuyant sur la thèse d'André Gunder Frank et sur l'exemple du cacao camerounais, discutez cette affirmation.
\tcblower
\textbf{Solution détaillée.} L'affirmation repose sur deux erreurs. \emph{Première erreur} : le sous-développement n'est pas un retard naturel mais un processus historiquement produit. Selon André Gunder Frank, la traite négrière, la colonisation et le commerce inégal ont « développé le sous-développement » : l'Afrique n'est pas en arrière sur un chemin commun, elle a été insérée dans un système qui l'appauvrit en enrichissant d'autres. \emph{Seconde erreur} : copier l'Europe ne suffit pas, car le problème n'est pas un modèle à imiter mais une position à transformer. L'exemple du cacao le prouve : le Cameroun exporte plus de 70 \% de sa production brute et importe des produits transformés ; tant que la valeur ajoutée est captée à l'extérieur, imiter les apparences du développement (consommer comme l'Occident) aggrave la dépendance au lieu de la réduire. \emph{Conclusion} : la sortie du sous-développement suppose un développement endogène — transformer sur place, bien gouverner, s'unir — plutôt que la copie d'un modèle extérieur. Le sous-développement n'est donc pas une fatalité, mais il ne se vainc pas par l'imitation.
\end{exoresolu}

\begin{exercice}[À faire en classe]
Relevez dans votre environnement trois exemples concrets de dépendance économique (produits importés alors qu'ils pourraient être fabriqués localement). Pour chacun, indiquez la cause possible (interne ou externe) et proposez une solution réaliste. Rédigez ensuite un paragraphe argumenté répondant à la question : « Le sous-développement du Cameroun est-il une fatalité ? »
\end{exercice}

\begin{corrige}[Exercice]
Exemples attendus (au choix) : importation de riz alors que les plaines du Nord peuvent en produire (cause interne : faible investissement agricole ; solution : soutien aux coopératives) ; importation de chocolat alors que le cacao est local (cause externe et interne : absence d'usines de transformation, prix du brut imposé ; solution : industrialisation locale) ; importation de médicaments (cause externe : brevets ; solution : production générique). Paragraphe type : le sous-développement n'est pas une fatalité car, étant historiquement produit (Frank), il peut être transformé par l'action ; mais cette transformation exige des réformes internes (gouvernance, transformation locale) et une renégociation de la place de l'Afrique dans le commerce mondial — ce qui prépare l'idée panafricaine d'unité pour peser davantage.
\end{corrige}

\begin{aretenir}[L'essentiel de la section]
\begin{itemize}
\item Le \cle{sous-développement} est la situation d'une société dont les richesses et potentialités sont bloquées ou exploitées au profit de l'extérieur : pauvreté, dépendance, blocage.
\item Ses manifestations : pauvreté monétaire, malnutrition, faible espérance de vie, analphabétisme, exode rural, fuite des cerveaux.
\item Thèse d'\textbf{André Gunder Frank} : le sous-développement n'est pas un retard naturel mais un processus historiquement produit (« le développement du sous-développement »).
\item Causes externes : traite négrière, colonisation, commerce inégal, endettement, prix imposés. Causes internes : mauvaise gouvernance, corruption, conflits, faible industrialisation, dépendance alimentaire.
\item Le sous-développement n'est pas une fatalité : l'optimisme du développement endogène ouvre la voie, que le panafricanisme tentera d'organiser.
\end{itemize}
\end{aretenir}

\coursec{Historique du panafricanisme}

Nous venons de voir que le sous-développement de l'Afrique n'est pas une fatalité naturelle : il a des causes historiques (traite négrière, colonisation, partage inégal des richesses) et des causes structurelles (dépendance économique, faiblesse de l'industrie, fragmentation du continent en petits États). Une question surgit alors naturellement : face à ce constat, les Africains eux-mêmes ont-ils élaboré une réponse ? La réponse est oui, et cette réponse porte un nom : le \cle{panafricanisme}. Comprendre son histoire, c'est comprendre comment une idée née parmi les descendants d'esclaves déportés en Amérique est devenue le projet politique de tout un continent.

\courssub{Les origines : un mouvement né dans la diaspora}

Commençons par une intuition simple. Imaginez des millions d'Africains arrachés à leur terre par la traite négrière, réduits en esclavage aux Antilles et en Amérique, puis leurs descendants, libres mais méprisés, victimes du racisme et de la ségrégation. Ces hommes et ces femmes découvrent peu à peu qu'ils partagent un même destin avec les Africains restés au continent, eux-mêmes colonisés à la fin du XIX\textsuperscript{e} siècle. De cette prise de conscience d'une \emph{communauté de souffrance et de destin} naît une idée force : les Noirs du monde entier doivent s'unir pour se libérer et se développer. C'est exactement le panafricanisme.

\begin{definition}[Le panafricanisme]
Le \cle{panafricanisme} est un mouvement de pensée et d'action politique visant l'\textbf{unité}, la \textbf{solidarité} et la \textbf{libération} de tous les Africains du continent et de la diaspora, en vue de leur \textbf{développement}. Le préfixe grec \emph{pan} signifie « tout » : le panafricanisme veut rassembler \emph{tous} les Africains, où qu'ils soient.
\end{definition}

Il faut insister sur un point décisif : le panafricanisme n'est pas né en Afrique, mais dans la \cle{diaspora}, c'est-à-dire parmi les populations noires dispersées par la traite négrière (Antilles, États-Unis, Amérique latine). Pourquoi ? Parce que c'est là, au contact direct du racisme le plus violent, mais aussi des idées modernes de liberté et d'égalité, que la conscience d'une identité africaine commune s'est d'abord forgée. Trois figures fondatrices doivent être retenues :

\begin{itemize}
\item \textbf{Henry Sylvester Williams} (1869-1911), avocat originaire de Trinité, organise à \textbf{Londres en 1900} le \emph{premier congrès panafricain} de l'histoire. Cet événement marque la naissance officielle du mouvement : pour la première fois, des Noirs des deux côtés de l'Atlantique se réunissent pour dénoncer ensemble le colonialisme et le racisme.
\item \textbf{W.E.B. Du Bois} (1868-1963), intellectuel noir américain, sociologue et historien, est le grand organisateur des congrès panafricains. Sa formule célèbre annonce le XX\textsuperscript{e} siècle : « le problème du XX\textsuperscript{e} siècle est le problème de la ligne de partage des couleurs ». Il finira sa vie au Ghana, invité par Nkrumah.
\item \textbf{Marcus Garvey} (1887-1940), leader jamaïcain, fonde l'UNIA (Association universelle pour l'amélioration de la condition noire) et prône le « retour à l'Afrique » (\emph{Back to Africa}) et la fierté noire. Il popularise le panafricanisme auprès des masses.
\end{itemize}

\begin{savaistu}[Le 25 mai, Journée de l'Afrique]
Le 25 mai, date de la signature de la Charte de l'OUA à Addis-Abeba en 1963, est célébré chaque année comme « \textbf{Journée de l'Afrique} » (\emph{Africa Day}). C'est un jour férié dans plusieurs pays africains, et l'occasion, dans les lycées camerounais, de conférences et de débats sur l'unité du continent.
\end{savaistu}

\courssub{Les congrès panafricains (1919-1945) : de la revendication à l'exigence d'indépendance}

Après le congrès de Londres (1900), une série de congrès panafricains jalonne la première moitié du XX\textsuperscript{e} siècle, animée principalement par Du Bois. Le congrès de \textbf{Paris en 1919} se tient en marge de la conférence de paix qui suit la Première Guerre mondiale : les Africains et les Afro-Américains, qui ont combattu pour la France et l'Angleterre, demandent en retour que les peuples colonisés soient associés à leur propre gouvernement. Les revendications restent alors modestes : on réclame des réformes, de meilleures conditions, une représentation — pas encore l'indépendance.

Le tournant décisif est le congrès de \textbf{Manchester en 1945} (5\textsuperscript{e} congrès). Trois changements majeurs le caractérisent :

\begin{enumerate}
\item \textbf{Le centre de gravité se déplace} : pour la première fois, les Africains du continent (Kwame Nkrumah du Ghana, Jomo Kenyatta du Kenya) y jouent un rôle de premier plan, aux côtés de la diaspora. Le panafricanisme revient ainsi vers l'Afrique.
\item \textbf{Les exigences se radicalisent} : on ne demande plus des réformes, on exige l'\textbf{indépendance}. La résolution finale proclame le droit des peuples colonisés à disposer d'eux-mêmes.
\item \textbf{Les méthodes s'affirment} : les congressistes recommandent l'action de masse — grèves, boycotts, désobéissance civile — sur le modèle de Gandhi.
\end{enumerate}

On peut donc établir le raisonnement suivant : entre 1900 et 1945, le panafricanisme est passé d'une \emph{pétition morale} (demander justice aux colonisateurs) à un \emph{programme politique} (conquérir la liberté par la lutte). La Seconde Guerre mondiale a joué ici un rôle d'accélérateur : les puissances coloniales, affaiblies et discréditées (elles prétendaient défendre la liberté tout en niant celle des autres), ne pouvaient plus légitimer leur domination.

\courssub{La Négritude : la dimension culturelle du panafricanisme}

Dans les années 1930, à Paris, de jeunes étudiants noirs — le Sénégalais \textbf{Léopold Sédar Senghor}, le Martiniquais \textbf{Aimé Césaire} et le Guyanais \textbf{Léon-Gontran Damas} — fondent le mouvement de la \cle{Négritude}. Le mot lui-même, forgé par Césaire, est une provocation assumée : il reprend l'injure raciste pour en faire un titre de fierté.

La Négritude est la \emph{réhabilitation des valeurs de la civilisation noire} : contre le mépris colonial qui présentait l'Africain comme un être sans histoire ni culture, Senghor affirme que « l'émotion est nègre comme la raison est hellène », valorisant la sensibilité, le rythme, l'art et la spiritualité africains. Césaire, dans son \emph{Cahier d'un retour au pays natal} (1939), célèbre l'identité noire et dénonce la colonisation comme un crime contre l'humanité.

Quel est le lien avec le panafricanisme ? La Négritude en est la \textbf{dimension culturelle} : là où les congrès panafricains revendiquent la liberté politique, la Négritude restaure la \emph{dignité culturelle} sans laquelle aucune libération n'est complète. Un peuple convaincu de sa propre infériorité ne peut pas se développer.

\begin{attention}[Erreur fréquente : confondre panafricanisme et Négritude]
Il ne faut pas \textbf{réduire le panafricanisme à la Négritude}. La Négritude n'en est que la \emph{dimension culturelle} (fierté de l'identité noire, valorisation des cultures africaines). Le panafricanisme est plus vaste : c'est un projet \emph{politique et économique} d'unité et de libération de tout le continent et de la diaspora. Dans une dissertation, écrire « le panafricanisme, c'est-à-dire la Négritude » est une erreur de définition qui coûte des points.
\end{attention}

\courssub{Les indépendances et Kwame Nkrumah : le panafricanisme au pouvoir}

Les années 1960 voient l'accession massive des colonies africaines à l'indépendance (le Cameroun accède à l'indépendance le 1\textsuperscript{er} janvier 1960). Mais le moment symbolique décisif a lieu quelques années plus tôt : le \textbf{6 mars 1957}, la Côte-de-l'Or devient le \textbf{Ghana}, première colonie d'Afrique noire à accéder à l'indépendance, sous la conduite de \textbf{Kwame Nkrumah}.

Nkrumah est la figure centrale du panafricanisme politique. Formé aux États-Unis, participant du congrès de Manchester, il fait de son pays indépendant la base arrière de la libération de tout le continent. Sa déclaration du 6 mars 1957 est restée célèbre :

\begin{quote}
\emph{« L'indépendance du Ghana est sans signification si elle n'est pas liée à la libération totale de l'Afrique. »}
\end{quote}

Analysons cette phrase, car elle contient toute la philosophie panafricaniste. Pour Nkrumah, l'indépendance d'un seul pays est une victoire incomplète, presque illusoire : un Ghana libre entouré d'une Afrique colonisée resterait fragile, isolé, incapable de se développer véritablement. La liberté est donc \emph{indivisible} : elle n'existe pleinement que si elle est partagée par tous. C'est pourquoi Nkrumah accueille à Accra les nationalistes de tout le continent, finance les mouvements de libération et convoque en 1958 la première Conférence des États africains indépendants. Sa thèse, exposée dans \emph{Africa Must Unite} (1963), est radicale : seuls des « États-Unis d'Afrique » pourront peser face aux anciennes puissances coloniales et éviter le \emph{néocolonialisme}, c'est-à-dire la domination économique qui survit à l'indépendance politique.

\courssub{la création de l'OUA à Addis-Abeba}

L'indépendance acquise, les États africains doivent répondre à une question redoutable : \emph{comment} réaliser l'unité ? Deux stratégies s'affrontent :

\begin{itemize}
\item Le \textbf{groupe de Casablanca} (1961), dit des « \cle{unionistes} » : Ghana de Nkrumah, Guinée de Sékou Touré, Mali, Maroc, Égypte, Algérie. Ils veulent une \emph{fédération immédiate}, un gouvernement continental, une armée commune — bref, les États-Unis d'Afrique, tout de suite.
\item Le \textbf{groupe de Monrovia} (1961), dit des « \cle{gradualistes} » : majorité des États, dont le Nigeria, l'Éthiopie, le Liberia et la plupart des pays francophones. Ils jugent l'union politique immédiate irréaliste et prônent une \emph{coopération progressive} entre États souverains, étape par étape (économie d'abord, politique ensuite).
\end{itemize}

Le \textbf{25 mai 1963}, à \textbf{Addis-Abeba} (Éthiopie), 32 chefs d'État africains signent la Charte fondatrice de l'\textbf{Organisation de l'Unité Africaine (OUA)}. L'OUA est un \emph{compromis} entre les deux groupes : elle affirme l'objectif panafricaniste de l'unité (concession aux unionistes), mais sur la base du respect de la souveraineté de chaque État, de la non-ingérence dans les affaires intérieures et de l'intangibilité des frontières héritées de la colonisation (concession aux gradualistes). Ses objectifs principaux : promouvoir l'unité et la solidarité, coordonner la coopération, défendre la souveraineté des États, et achever la décolonisation du continent (lutte contre l'apartheid en Afrique du Sud et contre le colonialisme portugais).

\begin{exemplebox}[L'OUA, un compromis révélateur]
Le choix d'Addis-Abeba comme siège n'est pas anodin : l'Éthiopie de Hailé Sélassié, jamais durablement colonisée (victoire d'Adoua sur l'Italie en 1896), incarne la continuité de la souveraineté africaine. Le compromis de 1963 montre une leçon philosophique générale : une idée généreuse (l'unité totale) doit souvent composer avec la réalité (la jeunesse et la fragilité des États). Le réalisme gradualiste l'a emporté sur l'utopie fédérale — avec les avantages (pas de conflit entre États) et les inconvénients (une unité faible, peu contraignante) que cela suppose.
\end{exemplebox}

\courssub{et après : de l'OUA à l'Union Africaine et à la ZLECAf}

À la fin du XX\textsuperscript{e} siècle, l'OUA apparaît dépassée : continent décolonisé, elle a rempli sa mission de libération, mais elle est restée impuissante face aux conflits, aux coups d'État et à la pauvreté. Le principe de non-ingérence l'empêchait d'agir. Sous l'impulsion notamment du Libyen Mouammar Kadhafi et du Sud-Africain Thabo Mbeki, l'OUA est transformée en \textbf{Union Africaine (UA)} en \textbf{2002} (sommet de Durban, Afrique du Sud). L'UA, inspirée du modèle de l'Union européenne, se dote d'institutions plus intégrées : Commission, Parlement panafricain, Conseil de paix et de sécurité, et même une Cour africaine. Nouveauté décisive : l'UA se reconnaît un \emph{droit d'intervention} dans un État membre en cas de crimes graves (génocide, crimes de guerre), rompant avec la non-ingérence absolue.

L'aboutissement économique le plus récent du panafricanisme est la \textbf{ZLECAf} (Zone de libre-échange continentale africaine), signée à Kigali en 2018 et entrée en vigueur opérationnelle le 1\textsuperscript{er} janvier 2021. Son principe : supprimer progressivement les droits de douane entre pays africains pour que les Africains commercent d'abord entre eux, au lieu d'exporter leurs matières premières brutes vers l'extérieur.

\begin{exemplebox}[Le panafricanisme en chiffres]
L'Union Africaine compte \textbf{55 États membres} : c'est l'organisation régionale la plus vaste du monde. La ZLECAf constitue un marché de plus de \textbf{1,3 milliard de personnes}, avec un produit intérieur brut cumulé d'environ \num{3400} milliards de dollars. À titre de comparaison, les échanges \emph{intra-africains} ne représentent qu'environ 15 à 18 \% du commerce total du continent, contre près de 60 \% en Europe : c'est précisément ce retard que la ZLECAf veut rattraper. Pour le Cameroun, cela signifie la possibilité de vendre davantage de produits transformés (cacao, bois, aluminium) à ses voisins plutôt que de tout exporter brut vers l'Europe ou l'Asie.
\end{exemplebox}

\courssub{Frise chronologique et carte : visualiser l'histoire du panafricanisme}

\begin{methode}[Construire une frise chronologique en dissertation]
Pour mémoriser et mobiliser efficacement l'histoire du panafricanisme, construisez une frise en suivant quatre étapes :
\begin{enumerate}
\item \textbf{Sélectionner} les dates charnières (celles où quelque chose change réellement : naissance, tournant, institutionnalisation), et non toutes les dates.
\item \textbf{Hiérarchiser} : distinguer les dates majeures (1900, 1945, 1963) des dates secondaires.
\item \textbf{Ordonner} les événements sur un axe gradué de manière régulière.
\item \textbf{Justifier} chaque choix en une phrase : « je retiens Manchester 1945 parce que c'est le tournant vers l'exigence d'indépendance ». C'est cette justification qui transforme une liste de dates en raisonnement historique.
\end{enumerate}
\end{methode}

\begin{popfigure}
\begin{center}
\begin{tikzpicture}[scale=1.0]
% Axe principal
\draw[-{Stealth[length=3mm]},line width=1.2pt,popdark] (0,0) -- (15.2,0);
% Graduations et événements
\foreach \x/\date/\lieu/\desc/\col in {
  0.8/1900/Londres/{1\textsuperscript{er} congrès panafricain (H.S. Williams)}/popblue,
  3.6/1945/Manchester/{Tournant : exigence d'indépendance}/poppurple,
  6.2/1957/Accra/{Indépendance du Ghana (Nkrumah)}/poporange,
  8.8/1963/Addis-Abeba/{Création de l'OUA (25 mai)}/popgreen,
  11.6/2002/Durban/{OUA devient Union Africaine}/poppink,
  14.2/2021/Continent/{Entrée en vigueur de la ZLECAf}/popgold}
{
  \draw[\col,line width=1.5pt] (\x,-0.15) -- (\x,0.15);
  \node[\col,font=\bfseries\small,above] at (\x,0.15) {\date};
  \node[\col,font=\footnotesize,above=0.42cm,align=center] at (\x,0.15) {\lieu};
  \node[popdark,font=\scriptsize,below=0.12cm,align=center,text width=2.2cm] at (\x,0) {\desc};
}
\end{tikzpicture}
\end{center}
\legende{Frise chronologique des grandes étapes du panafricanisme, de 1900 à 2021. Chaque date charnière marque une transformation du mouvement : naissance dans la diaspora (1900), radicalisation (1945), passage au pouvoir (1957), institutionnalisation (1963), approfondissement (2002) et intégration économique (2021).}
\end{popfigure}

\begin{popfigure}
\begin{center}
\begin{tikzpicture}[scale=0.9]
% Contour très simplifié de l'Afrique
\draw[fill=popcream,draw=popdark,line width=1pt] 
  (0.2,4.2) -- (1.2,5.0) -- (2.6,5.3) -- (4.2,5.1) -- (5.4,4.6) -- (6.4,4.4)
  -- (7.2,3.8) -- (7.4,3.0) -- (6.8,2.4) -- (6.6,1.6) -- (6.2,0.8) -- (5.6,0.2)
  -- (4.8,-0.4) -- (4.0,-0.6) -- (3.2,-0.2) -- (2.6,0.6) -- (2.0,1.4) -- (1.2,1.8)
  -- (0.4,2.4) -- (-0.2,3.0) -- (-0.4,3.6) -- cycle;
% Zones régionales (approximatives)
\fill[popblueL,opacity=0.75] (0.0,3.4) -- (1.4,3.8) -- (1.8,2.6) -- (0.6,2.2) -- (-0.2,2.8) -- cycle;
\fill[popgreenL,opacity=0.75] (2.4,2.8) -- (4.0,2.8) -- (4.2,1.6) -- (3.0,1.2) -- (2.2,1.8) -- cycle;
\fill[poporangeL,opacity=0.75] (3.0,1.0) -- (4.4,1.2) -- (4.6,0.0) -- (3.6,-0.3) -- (2.8,0.4) -- cycle;
% Villes
\fill[popink] (5.9,3.1) circle (2.2pt); \node[right,font=\scriptsize\bfseries,popink] at (6.0,3.1) {Addis-Abeba};
\fill[popink] (1.5,2.35) circle (2.2pt); \node[left,font=\scriptsize\bfseries,popink] at (1.4,2.35) {Accra};
\fill[popink] (2.9,2.0) circle (2.2pt); \node[below,font=\scriptsize\bfseries,popink] at (2.9,1.85) {Douala/Yaoundé};
% Nord
\draw[-{Stealth},popdark,line width=1pt] (7.8,4.6) -- (7.8,5.4); \node[above,font=\scriptsize\bfseries,popdark] at (7.8,5.4) {N};
% Légende des zones
\node[font=\scriptsize,anchor=west] at (-1.2,-1.2) {\textcolor{popblue}{\rule{0.35cm}{0.2cm}} CEDEAO (Afrique de l'Ouest)};
\node[font=\scriptsize,anchor=west] at (3.2,-1.2) {\textcolor{popgreen}{\rule{0.35cm}{0.2cm}} CEMAC (Afrique centrale)};
\node[font=\scriptsize,anchor=west] at (7.0,-1.2) {\textcolor{poporange}{\rule{0.35cm}{0.2cm}} SADC (Afrique australe)};
\end{tikzpicture}
\end{center}
\legende{Croquis très simplifié de l'Afrique (échelle indicative, contours approximatifs). 1 : Accra (Ghana), capitale de l'indépendance de 1957 et siège du secrétariat de la ZLECAf ; 2 : Addis-Abeba (Éthiopie), siège de l'OUA puis de l'Union Africaine ; 3 : le Cameroun, membre de la CEMAC. Les ensembles régionaux (CEDEAO, CEMAC, SADC) sont les briques économiques de l'intégration continentale voulue par le panafricanisme.}
\end{popfigure}

\courssub{Exercice résolu}

\begin{exoresolu}[Analyse de document : la déclaration de Nkrumah]
\textbf{Énoncé.} Le 6 mars 1957, Kwame Nkrumah déclare : « L'indépendance du Ghana est sans signification si elle n'est pas liée à la libération totale de l'Afrique. » (1) Identifiez la thèse défendue. (2) Expliquez pourquoi cette déclaration est un texte panafricaniste. (3) Montrez en quoi la création de l'OUA en 1963 répond partiellement à cette exigence.
\tcblower
\textbf{Solution détaillée.} (1) La thèse de Nkrumah est que la liberté est \emph{indivisible} : une indépendance nationale isolée est incomplète et fragile ; seule la libération de tout le continent donne son sens plein à l'indépendance de chaque État. (2) Ce texte est panafricaniste parce qu'il refuse de penser la liberté à l'échelle d'un seul pays et la rapporte à l'unité de \emph{tous} les Africains : il exprime la solidarité entre le continent et affirme que le développement du Ghana dépend de celui de l'Afrique entière, ce qui correspond exactement à la définition du panafricanisme comme mouvement visant l'unité et la libération de tous les Africains. (3) L'OUA, créée à Addis-Abeba le 25 mai 1963, répond partiellement à cette exigence : elle institutionnalise la solidarité africaine et soutient activement la décolonisation (apartheid, colonies portugaises). Mais « partiellement » seulement : le compromis avec les gradualistes de Monrovia a sacrifié le projet fédéral de Nkrumah (les États-Unis d'Afrique) au profit d'une simple coopération entre États souverains. On peut donc conclure que l'histoire du panafricanisme est celle d'une tension permanente entre l'idéal de l'unité totale et le réalisme des étapes graduelles.
\end{exoresolu}

\begin{aretenir}[L'essentiel de l'historique]
\begin{itemize}
\item Le panafricanisme naît dans la \textbf{diaspora} (Antilles, Amérique) à la fin du XIX\textsuperscript{e} siècle ; premier congrès : \textbf{Londres 1900} (Sylvester Williams). Figures : Du Bois, Garvey.
\item Les congrès (1919-1945) conduisent au tournant de \textbf{Manchester 1945} : exigence d'indépendance et retour du mouvement vers le continent.
\item La \textbf{Négritude} (Senghor, Césaire, Damas) est la dimension \emph{culturelle} du panafricanisme, non sa totalité.
\item \textbf{1957} : indépendance du Ghana ; Nkrumah fait du panafricanisme un programme d'État.
\item \textbf{25 mai 1963} : Charte de l'\textbf{OUA} à Addis-Abeba, compromis entre unionistes (Casablanca) et gradualistes (Monrovia).
\item \textbf{2002} : l'OUA devient l'\textbf{Union Africaine} ; \textbf{2021} : lancement opérationnel de la \textbf{ZLECAf}, marché de plus de 1,3 milliard de personnes.
\end{itemize}
\end{aretenir}

Cette histoire montre que le panafricanisme est passé par trois âges : l'âge de l'\emph{idée} (diaspora, congrès), l'âge du \emph{pouvoir} (indépendances, Nkrumah) et l'âge des \emph{institutions} (OUA, UA, ZLECAf). Mais une idée peut-elle survivre à ses institutions ? L'unité proclamée est-elle une unité réelle ? C'est en examinant les objectifs puis les difficultés du panafricanisme que nous pourrons répondre à la question qui clôture cette leçon : l'unité africaine est-elle véritablement la voie du développement ?

\coursec{Objectifs du panafricanisme comme philosophie du développement}

Après avoir retracé l'\emph{historique} du panafricanisme — des premiers congrès de la diaspora (1900, 1919, 1945) à la création de l'OUA en 1963 —, nous comprenons que ce mouvement n'est pas seulement une réaction sentimentale au colonialisme. Il se structure en une \cle{philosophie du développement} dotée d'objectifs précis, articulés autour de trois piliers indissociables : le politique, l'économique et le culturel. Comme l'affirme Kwame Nkrumah : « Cherchez d'abord le royaume politique et tout le reste vous sera donné par surcroît. » Mais ce royaume politique n'est pas une fin en soi ; il est l'instrument indispensable pour conquérir la souveraineté économique et restaurer la dignité culturelle.

\courssub{L'objectif politique : libération totale, unité des États et voix commune}

L'objectif politique premier est la \cle{libération totale} du continent : éradication du colonialisme direct, puis du néocolonialisme (dépendance économique, ingérences militaires, conditionnalités du FMI/Banque mondiale). Cette libération passe par l'\cle{unité des États africains}. Deux conceptions s'affrontent historiquement :
\begin{itemize}
    \item Le \cle{groupe de Casablanca} (Ghana, Guinée, Mali, Égypte, Maroc, Algérie) : préconise une unité fédérale immédiate, un « gouvernement continental » avec une armée, une monnaie, une citoyenneté communes (thèse de Nkrumah : \emph{United States of Africa}).
    \item Le \cle{groupe de Monrovia} (Nigéria, Liberia, Côte d'Ivoire, Sénégal, etc.) : privilégie une coopération progressive, respectueuse des souverainetés nationales, aboutissant à l'OUA (1963), puis à l'Union Africaine (UA, 2002).
\end{itemize}
Aujourd'hui, l'UA synthétise ces approches : elle respecte les frontières héritées (\emph{uti possidetis juris}) mais construit des institutions supranationales (Parlement panafricain, Cour de justice, Commission de l'UA, Conseil de paix et de sécurité \cle{CPS}). L'objectif ultime est une \cle{voix commune} dans les relations internationales : peser à l'ONU, dans les négociations climatiques (COP), commerciales (OMC) et géopolitiques face aux grandes puissances (USA, Chine, UE, Russie).

\begin{savaistu}[Nkrumah et le groupe de Casablanca]
Dès 1963, au sommet d'Addis-Abeba, Kwame Nkrumah propose un « gouvernement continental » doté d'un commandement militaire unifié, d'une monnaie unique et d'une politique étrangère unique. Ses partisans (Ghana, Guinée, Mali, Égypte, Maroc, Algérie) forment le « groupe de Casablanca », opposé au « groupe de Monrovia » plus modéré. L'OUA naîtra du compromis : charte respectant la souveraineté des États, mais engagement solennel de libération totale (décolonisation de l'Afrique australe) et non-ingérence. Nkrumah, renversé en 1966, reste le symbole de l'unité fédérale radicale.
\end{savaistu}

\courssub{L'objectif économique : intégration, industrialisation, transformation locale et souveraineté monétaire}

L'indépendance politique sans indépendance économique est une coquille vide. L'objectif économique structurel est de briser l'héritage colonial : l'Afrique productrice de matières premières (cacao, café, coton, minerais, pétrole) exportées brutes, importatrice de biens manufacturés. Trois leviers sont actionnés :

\begin{enumerate}
    \item \textbf{Intégration des marchés :} Créer un marché unique continental. La \cle{ZLECAf} (Zone de libre-échange continentale africaine), entrée en vigueur en 2021, vise à éliminer les droits de douane sur 90\% des lignes tarifaires, libérant la circulation des biens, services, capitaux et personnes. C'est le plus grand marché unique au monde par le nombre de pays (54 signataires, 47 ratifications en 2024).
    \item \textbf{Industrialisation et transformation locale :} Passer de l'exportation de fèves de cacao à l'exportation de chocolat, du bois brut aux meubles, du minerai de fer à l'acier. L'Agenda 2063 de l'UA fixe l'industrialisation comme priorité absolue (programme \cle{PIDA} pour les infrastructures, \cle{AIDA} pour l'industrialisation).
    \item \textbf{Souveraineté monétaire et bancaire :} Sortir de la dépendance au Franc CFA (zone CEMAC/UEMOA) ou aux devises étrangères. Projets : monnaie unique de la CEDEAO (l'\cle{Eco}, sans cesse reportée), Banque centrale africaine (prévue par le Traité d'Abuja de 1991), Banque africaine de développement (\cle{BAD}) comme levier de financement endogène.
\end{enumerate}

\begin{exemplebox}[Le cacao ivoiro-ghanéen : levier de transformation]
La Côte d'Ivoire et le Ghana produisent ensemble \~60\% du cacao mondial. En 2019, ils instaurent un \emph{Différentiel de Revenu Décent} (DRD) de 400 \$/tonne pour garantir un prix plancher aux planteurs. C'est une application concrète de la \cle{puissance de négociation collective} (objectif politique) au service de la \cle{redistribution de la valeur ajoutée} (objectif économique). Prochaine étape : transformer 50\% du cacao sur place d'ici 2030 (usines de broyage, chocolateries).
\end{exemplebox}

\courssub{L'objectif culturel : revalorisation, dignité et estime de soi}

Le colonialisme a été une entreprise de \cle{déshumanisation culturelle} : langues interdites à l'école, histoire niée (« l'Afrique n'a pas d'histoire », Hegel), religions traditionnelles diabolisées, normes esthétiques européennes imposées. Le panafricanisme culturel vise la \cle{restauration de la dignité} :
\begin{itemize}
    \item \textbf{Réhabilitation de l'histoire :} Cheikh Anta Diop (\emph{Nations nègres et culture}, 1954) démontre l'antériorité de la civilisation négro-africaine (Égypte pharaonique) et l'apport de l'Afrique à l'humanité.
    \item \textbf{Valorisation des langues africaines :} Promotion du swahili, haoussa, yoruba, lingala, bamoun, fulfulde, etc., comme langues de travail de l'UA et vecteurs d'intégration.
    \item \textbf{Philosophie africaine :} L'\cle{ubuntu} (« Je suis parce que nous sommes ») comme éthique du développement : la personne ne se réalise qu'en communauté. Cela inspire des politiques publiques inclusives (justice transitionnelle post-génocide au Rwanda \emph{Gacaca}, politique d'habitat groupé \emph{Imidugudu}).
    \item \textbf{Restitution du patrimoine :} Combat pour le retour des œuvres d'art spoliées (rapport Sarr-Savoy 2018, restitutions récentes du Bénin, du Sénégal, du Nigeria).
\end{itemize}

\begin{definition}[Développement endogène]
Le \textbf{développement endogène} est un développement \emph{pensé, décidé et conduit par les Africains eux-mêmes}, à partir de \emph{leurs propres ressources} (humaines, naturelles, culturelles, financières), en réponse à \emph{leurs besoins réels} (sécurité alimentaire, santé, éducation, emploi) et en conformité avec \emph{leurs valeurs} (solidarité, respect des aînés, rapport harmonieux à la nature). Il s'oppose au développement \emph{exogène} (imposé par des modèles occidentaux, conditionné par l'aide publique au développement, déconnecté des réalités locales).
\end{definition}

\courssub{L'unité comme condition \emph{sine qua non} du développement : l'argument de la masse critique}

Pourquoi l'unité n'est-elle pas un simple idéal romantique, mais une \cle{nécessité structurelle} ?
\begin{itemize}
    \item \textbf{Taille des marchés :} 54 États, dont 16 enclavés, avec des marchés nationaux trop étroits pour rentabiliser des usines (économies d'échelle). Le marché unique ZLECAf = 1,4 milliard de consommateurs, PIB cumulé \~3 400 milliards \$.
    \item \textbf{Pouvoir de négociation :} Seul, le Cameroun négocie un accord de partenariat économique (APE) avec l'UE en position de faiblesse. Uni (UA + groupes régionaux), le continent impose la réciprocité, les clauses de sauvegarde, le transfert de technologie.
    \item \textbf{Infrastructures transversales :} Barrage de Grand Inga (RDC, 40 GW potentiel) ne peut être rentable que s'il alimente un réseau panafricain (projet \emph{Cairo-Cape Town Highway}, \emph{Transsaharien}, interconnexion électrique WAPP/EAPP/CAPP).
    \item \textbf{Paix et sécurité :} L'architecture de paix et de sécurité de l'UA (CPS, Force africaine en attente) permet de gérer les conflits (Soudan, RDC, Sahel) sans attendre l'ONU ou l'ex-puissance coloniale.
\end{itemize}
Nkrumah résumait ainsi : « Nous faisons face ni à l'Est ni à l'Ouest ; nous faisons face à l'avant. » L'unité est la \cle{condition de possibilité} pour passer du statut d'objet de l'histoire à celui de sujet.

\begin{center}
\begin{popfigure}
\begin{tikzpicture}[scale=0.9, every node/.style={font=\small}]
    % Couleurs
    \colorlet{polcol}{popblue!80}
    \colorlet{ecocol}{popgreen!80}
    \colorlet{culcol}{poporange!80}
    \colorlet{centcol}{poppurple!90}
    % Cercle central : Développement
    \node[draw=centcol, thick, circle, minimum width=3.2cm, fill=centcol!15, align=center] (centre) at (0,0){\textbf{DÉVELOPPEMENT}\\\textbf{ DE L'AFRIQUE}\\(endogène, souverain,\\inclusif, durable)};
    % Objectif Politique
    \node[draw=polcol, thick, rounded corners=15pt, fill=polcol!10, text width=4cm, align=left] (pol) at (-5, 3) {
        \textbf{POLITIQUE}\\
        $\bullet$ Libération totale (néocolonialisme)\\
        $\bullet$ Unité des États (UA, Parlement, Armée)\\
        $\bullet$ Voix commune (ONU, Climat, Commerce)\\
        $\bullet$ Paix \& Sécurité (CPS, Force en attente)
    };
    % Objectif Économique
    \node[draw=ecocol, thick, rounded corners=15pt, fill=ecocol!10, text width=4cm, align=left] (eco) at (5, 3) {
        \textbf{ÉCONOMIQUE}\\
        $\bullet$ Intégration marchés (ZLECAf)\\
        $\bullet$ Industrialisation (AIDA, PIDA)\\
        $\bullet$ Transformation locale (matières 1ères)\\
        $\bullet$ Souveraineté monétaire (Eco, BAD, BC Africaine)
    };
    % Objectif Culturel
    \node[draw=culcol, thick, rounded corners=15pt, fill=culcol!10, text width=4cm, align=left] (cul) at (0, -3.5) {
        \textbf{CULTUREL}\\
        $\bullet$ Réhabilitation histoire (C.A. Diop)\\
        $\bullet$ Valorisation langues (Swahili, etc.)\\
        $\bullet$ Éthique Ubuntu / Maât\\
        $\bullet$ Restitution patrimoine / Dignité
    };
    % Flèches de convergence
    \draw[->, thick, polcol] (pol.south east) -- node[above right, sloped, xshift=2mm] {légitimité} (centre.north west);
    \draw[->, thick, ecocol] (eco.south west) -- node[above left, sloped, xshift=-2mm] {moyens} (centre.north east);
    \draw[->, thick, culcol] (cul.north) -- node[right] {sens \& valeurs} (centre.south);
    % Interactions entre piliers
    \draw[<->, dashed, popmuted] (pol.east) -- node[above] {stabilité} (eco.west);
    \draw[<->, dashed, popmuted] (pol.south) .. controls (0,1) .. node[left] {citoyenneté} (cul.north west);
    \draw[<->, dashed, popmuted] (eco.south) .. controls (0,-1) .. node[right] {capital humain} (cul.north east);
\end{tikzpicture}
\legende{Schéma des trois piliers convergents du panafricanisme : les objectifs politique, économique et culturel sont interdépendants et convergent vers un développement endogène de l'Afrique.}
\end{popfigure}
\end{center}

\courssub{Difficultés et limites du panafricanisme}

Malgré la clarté des objectifs, la réalisation bute sur des obstacles structurels et conjoncturels majeurs :
\begin{itemize}
    \item \textbf{Jalousie des souverainetés :} Réticence des États à céder des prérogatives (douane, monnaie, défense) à des instances supranationales (ex. : non-ratification du protocole sur la libre circulation des personnes par plusieurs pays).
    \item \textbf{Éclatement institutionnel (\emph{spaghetti bowl}) :} Chevauchement des 8 CER reconnues par l'UA (CEMAC, CEDEAO, SADC, EAC, COMESA, CEN-SAD, IGAD, UMA) avec des règles différentes, complexifiant l'intégration.
    \item \textbf{Déficit infrastructurel et financier :} Manque de routes, rails, énergie, ports ; dépendance au financement extérieur pour le budget de l'UA et les grands projets (PIDA).
    \item \textbf{Barrières non-tarifaires :} Corruption aux frontières, normes techniques divergentes, tracasseries administratives, freinent le commerce intra-africain plus que les droits de douane.
    \item \textbf{Instabilité politique et sécuritaire :} Coups d'État récurrents (Sahel, Afrique centrale), terrorisme, rébellions, déplacements de populations (ex. : retrait du Mali, Burkina Faso, Niger de la CEDEAO en 2024-2025 pour former l'AES).
    \item \textbf{Hétérogénéité économique :} Économies souvent concurrentes (productrices des mêmes matières premières) plutôt que complémentaires, rendant l'intégration commerciale difficile.
    \item \textbf{Dépendance extérieure persistante :} Accords de partenariat asymétriques (APE UE-ACP), endettement auprès de créanciers non-africains, conditionnalités FMI/BM.
\end{itemize}

\courssub{Applications concrètes : intégration régionale (CER) et réussites nationales}

Le panafricanisme ne vit pas seulement dans les textes de l'UA ; il s'incarne dans les \cle{Communautés économiques régionales (CER)}, piliers de la construction continentale.

\begin{center}
\begin{tabularx}{\linewidth}{|l|X|X|X|}
\hline
\rowcolor{popblueL}
\textbf{Communauté} & \textbf{Pays membres (ex.)} & \textbf{Institutions clés} & \textbf{Réalisations / Enjeux} \\
\hline
\cle{CEMAC} (Communauté Économique et Monétaire de l'Afrique Centrale) & Cameroun, Tchad, RCA, Gabon, Guinée équatoriale, Congo & \cle{BEAC} (Banque des États de l'Afrique Centrale), Parlement, Cour de justice & Monnaie commune (FCFA), marché commun (libre circulation personnes/biens \emph{théorique}), interconnectivité routière (Ndjamena-Douala), défis : convergence budgétaire, intégration réelle faible. \\
\hline
\cle{CEDEAO} (Afrique de l'Ouest) & 12 pays (Nigéria, Ghana, Sénégal, Côte d'Ivoire...)\newline*Note : Mali, Burkina, Niger retirés (AES 2024-25)* & Banque centrale (BCEAO), Parlement, Cour de justice, Force en attente & Libre circulation effective (passeport biométrique), projet monnaie unique \cle{Eco}, interventions militaires (Gambie, Guinée). Crise AES majeure. \\
\hline
\cle{SADC} (Afrique Australe) & 16 pays (Afrique du Sud, Angola, RDC, Zambie...) & Secrétariat, Tribunal, Parlement & Intégration énergétique (SAPP), corridor Nord-Sud, puissance industrielle (Afrique du Sud). \\
\hline
\cle{EAC} (Afrique de l'Est) & 8 pays (Kenya, Tanzanie, Ouganda, Rwanda, Burundi, RDC, Soudan du Sud, Somalie 2023) & Douane unique, Marché commun, Monnaie unique (projet Shilling) & Intégration la plus avancée : libre circulation travail/capitaux, passeport unique, infrastructure ferroviaire (SGR). \\
\hline
\cle{COMESA} & 21 pays (Égypte, RDC, Kenya, Zambie...) & Secrétariat, Cour de justice, Banque de développement & Zone de libre-échange, Union douanière (en cours), grand marché potentiel. \\
\hline
\cle{IGAD} & 8 pays (Djibouti, Érythrée, Éthiopie, Kenya, Somalie, Soudan, Soudan du Sud, Ouganda) & Secrétariat & Paix/sécurité (Soudan, Somalie), résilience climatique, développement bassin Nil. \\
\hline
\cle{CEN-SAD} & 28 pays (zone sahélo-saharienne) & Secrétariat (Tripoli) & Grande zone, mais intégration très faible, structure légère. \\
\hline
\cle{UMA} (Union du Maghreb Arabe) & 5 pays (Maroc, Algérie, Tunisie, Libye, Mauritanie) & Secrétariat (Rabat) & Quasi inactive (conflit Sahara occidental, tensions Algérie-Maroc). \\
\hline
\end{tabularx}
\end{center}

\textbf{Réussites nationales inspirantes (modèles de développement endogène) :}
\begin{itemize}
    \item \textbf{Maurice :} Passage économie sucrière mono-culturelle $\to$ services financiers, textile, TIC, tourisme. PIB/hab \~11 000 \$ (2023). Gouvernance stable, éducation gratuite, protection sociale.
    \item \textbf{Rwanda :} Post-génocide (1994), \emph{Vision 2020/2050} : « Singapour de l'Afrique ». Propreté (Kigali), e-gouvernance (Irembo), assurance santé universelle (Mutuelles de santé \~90\% couverture), attractivité investisseurs (Kigali International Financial Centre).
    \item \textbf{Botswana :} Gestion transparente des diamants (Debswana = partenariat 50/50 De Beers/Gouvernement), fonds souverain (Pula Fund), démocratie stable, éducation/santé gratuites. PIB/hab \~7 000 \$.
\end{itemize}
Ces pays partagent : \cle{leadership visionnaire}, \cle{institutions fortes}, \cle{investissement capital humain}, \cle{diversification économique}.

\begin{exemplebox}[Chiffre clé : le déficit d'intégration commerciale]
Le \textbf{commerce intra-africain} ne représente qu'environ \textbf{15 à 18\%} du commerce total du continent (données CNUCED/BAD 2022-2023), contre \textbf{plus de 60\% en Europe} (UE) et \textbf{environ 60\% en Asie} (ASEAN+Chine/Japon/Corée). L'Afrique commerce davantage \emph{avec le reste du monde} qu'entre elle-même. La ZLECAf vise à porter ce taux à 25\% d'ici 2030, puis 50\% à l'horizon 2040. C'est l'enjeu économique majeur du panafricanisme contemporain.
\end{exemplebox}

\begin{center}
\begin{popfigure}
\begin{tikzpicture}
\begin{axis}[
    ybar,
    bar width=18pt,
    width=\linewidth,
    height=9cm,
    enlarge x limits=0.35,
    symbolic x coords={Intra-Africain, Intra-Européen (UE), Intra-Asiatique (ASEAN+)},
    xtick=data,
    nodes near coords,
    nodes near coords align={vertical},
    ylabel={Part dans le commerce total (\%)},
    ymin=0, ymax=75,
    ytick={0,15,30,45,60,75},
    ymajorgrids=true,
    grid style=dashed,
    axis lines*=left,
    tick label style={font=\small, align=center},
    every node near coord/.append style={font=\small\bfseries, popdark},
    legend style={at={(0.5,-0.2)}, anchor=north, legend columns=-1, font=\small},
]
\addplot[fill=popblue!80, draw=popblue!60] coordinates {(Intra-Africain,16) (Intra-Européen (UE),62) (Intra-Asiatique (ASEAN+),60)};
\addplot[fill=poporange!80, draw=poporange!60] coordinates {(Intra-Africain,84) (Intra-Européen (UE),38) (Intra-Asiatique (ASEAN+),40)};
\legend{Commerce intra-régional, Commerce extra-régional}
\end{axis}
\end{tikzpicture}
\legende{Part du commerce intra-régional vs extra-régional (estimations CNUCED/BAD ~2022). L'Afrique est la région la moins intégrée commercialement : 84\% de ses échanges se font vers l'extérieur. Source : UNCTAD Handbook of Statistics, African Development Bank.}
\end{popfigure}
\end{center}

\courssub{TP guidé : construction d'une carte des organisations régionales et d'un diagramme des échanges}

\begin{methode}[TP : Cartographier l'intégration africaine et visualiser le déficit commercial]
\textbf{Objectif :} Appliquer les notions d'intégration régionale et de commerce intra-africain à des données réelles pour produire une analyse critique.

\textbf{Matériel :} Fond de carte Afrique (politique, A4), jeux de données (CNUCED, BAD, FMI) : PIB/hab, commerce intra/extra par CER, statuts ZLECAf (ratifié/non ratifié).

\textbf{Étape 1 — Carte des CER (8 CER reconnues par l'UA) :}
\begin{enumerate}
    \item Colorier chaque CER d'une couleur distincte (ex. : CEMAC=vert, CEDEAO=bleu, SADC=orange, EAC=rouge, COMESA=violet, CEN-SAD=jaune, IGAD=rose, UMA=gris).
    \item Superposer les pays membres de plusieurs CER (hachures croisées : ex. RDC $\in$ SADC, COMESA, EAC ; Congo $\in$ CEMAC, CEEAC).
    \item Localiser les sièges : Yaoundé (CEMAC), Abuja (CEDEAO), Gaborone (SADC), Arusha (EAC), Lusaka (COMESA), Tripoli (CEN-SAD), Djibouti (IGAD), Rabat (UMA).
    \item Tracer 3 corridors prioritaires PIDA : \emph{Transsaharien} (Alger-Lagos), \emph{Transafricain} (Dakar-Ndjamena-Djibouti), \emph{Nord-Sud} (Le Cap-Le Caire).
\end{enumerate}

\textbf{Étape 2 — Diagramme en barres groupées (pgfplots/Excel) :}
\begin{enumerate}
    \item Axe X : 8 CER. Axe Y : \% commerce intra-CER / commerce total de la CER.
    \item Données approx. 2022 : CEMAC \~5\%, CEDEAO \~12\%, SADC \~20\%, EAC \~18\%, COMESA \~8\%, UMA \~4\%, IGAD \~6\%, CEN-SAD \~3\%.
    \item Ajouter une barre « ZLECAf cible 2030 : 25\% » et « UE : 62\% » pour comparaison.
\end{enumerate}

\textbf{Étape 3 — Analyse rédigée (15 lignes) :}
\begin{itemize}
    \item Constater l'hétérogénéité (EAC/SADC leaders, CEMAC/UMA/CEN-SAD à la traîne).
    \item Expliquer : infrastructures, conflits, barrières non-tarifaires, monnaies non convertibles, économies similaires (concurrence vs complémentarité).
    \item Conclure sur l'urgence ZLECAf + infrastructures + libre circulation effective.
\end{itemize}
\end{methode}

\begin{attention}[Erreur fréquente à l'examen]
Ne citez \textbf{jamais} les objectifs du panafricanisme (unité, industrialisation, culture) \emph{sans les illustrer par des institutions ou des projets précis}.
\begin{itemize}
    \item \textbf{Mauvais :} « Le panafricanisme veut l'unité économique. »
    \item \textbf{Bon :} « L'objectif économique se concrétise par la ZLECAf (2021), le PIDA (infrastructures), l'AIDA (industrialisation), la future Banque centrale africaine et la monnaie Eco (CEDEAO). Au niveau sous-régional, la CEMAC dispose de la BEAC et d'un marché commun encore perfectible. »
\end{itemize}
L'examinateur attend la preuve que vous connaissez l'\emph{architecture institutionnelle} actuelle, pas seulement les slogans des années 1960.
\end{attention}

\courssub{Exercice résolu : Appliquer la notion de développement endogène}

\begin{exoresolu}[Analyse du lancement de la ZLECAf au prisme du développement endogène]
\textbf{Énoncé :} En quoi le lancement de la Zone de libre-échange continentale africaine (ZLECAf) en janvier 2021 illustre-t-il la mise en œuvre de la philosophie du développement endogène ? Justifiez en mobilisant la définition et trois arguments précis.

\tcblower
\textbf{Solution détaillée :}

\textbf{1. Rappel de la définition (critère de réussite) :}
Le développement endogène est un développement \emph{pensé, décidé et conduit par les Africains}, à partir de \emph{leurs ressources}, pour \emph{leurs besoins}, selon \emph{leurs valeurs}.

\textbf{2. Identification des éléments pertinents dans la situation (ZLECAf) :}
\begin{itemize}
    \item \textbf{Initiative africaine :} Négociée sous l'égide de l'UA (institution panafricaine), par des ministres africains du Commerce, sans tutelle occidentale (contrairement aux APE UE-ACP imposés).
    \item \textbf{Ressources mobilisées :} Marché de 1,4 milliard de consommateurs, main-d'œuvre jeune, matières premières, potentialités agricoles/industrielles.
    \item \textbf{Besoins ciblés :} Emploi (industrialisation), sécurité alimentaire (libre circulation produits vivriers), réduction pauvreté (revenus paysans via élimination droits de douane intra-africains).
    \item \textbf{Valeurs :} Solidarité (Traitement Spécial et Différencié \cle{TSD} pour PMA), paix (interdépendance économique), souveraineté (réduction dépendance exportations hors continent).
\end{itemize}

\textbf{3. Établissement des liens (argumentation) :}
\begin{itemize}
    \item \textbf{Argument 1 (Appropriation politique) :} La ZLECAf est un \emph{choix souverain} des 54 États signataires (ratifiée par 47 en 2024). Elle traduit la volonté politique de s'affranchir de la fragmentation héritée de Berlin (1885) et des accords de partenariat asymétriques (APE). C'est l'application de l'objectif politique (unité, voix commune) au service de l'économique.
    \item \textbf{Argument 2 (Mobilisation ressources internes) :} En supprimant les droits de douane sur 90\% des lignes tarifaires, la ZLECAf stimule la \emph{demande intra-africaine} pour les produits africains. Elle incite à transformer localement (ex. coton du Burkina $\to$ textile au Bénin $\to$ habillement au Maroc) plutôt qu'exporter brut vers l'Europe/Chine. C'est l'utilisation des \emph{ressources africaines} (travail, matières premières, marché) pour créer de la valeur \emph{en Afrique}.
    \item \textbf{Argument 3 (Réponse aux besoins structurels) :} L'Afrique importe \~40 milliards \$ de denrées alimentaires/an. La ZLECAf, couplée au Programme détaillé pour le développement de l'agriculture africaine (PDDAA), vise l'autosuffisance via le commerce intra-régional (riz de la Casamance $\to$ Sahel, maïs zambien $\to$ RDC). Elle répond au \emph{besoin vital} de sécurité alimentaire par la coopération, non par l'aide alimentaire.
\end{itemize}

\textbf{4. Conclusion :}
La ZLECAf est l'\emph{outil opérationnel majeur} du panafricanisme contemporain pour réaliser le développement endogène. Elle ne garantit pas le succès (défis : infrastructures, barrières non-tarifaires, volonté politique), mais elle en pose le \emph{cadre institutionnel africain}, pensé par les Africains pour les Africains.
\end{exoresolu}

\begin{aretenir}[Synthèse : Les trois piliers du panafricanisme comme philosophie du développement]
\begin{itemize}
    \item \textbf{Politique :} Libération totale $\to$ Unité (UA) $\to$ Voix commune $\to$ Paix/Sécurité (CPS). \emph{Condition sine qua non.}
    \item \textbf{Économique :} Intégration (ZLECAf) $\to$ Industrialisation (AIDA) $\to$ Transformation locale $\to$ Souveraineté monétaire/financière (Eco, BC Africaine). \emph{Moteur matériel.}
    \item \textbf{Culturel :} Réhabilitation histoire (C.A. Diop) $\to$ Valorisation langues $\to$ Éthique Ubuntu/Maât $\to$ Restitution patrimoine. \emph{Âme et direction.}
\end{itemize}
\textbf{Thèse centrale :} L'unité n'est pas une option sentimentale, c'est une \textbf{nécessité structurelle} (marché, négociation, infrastructures, sécurité) pour passer du sous-développement (extroversion, dépendance) au développement endogène (autonomie, dignité, prospérité partagée). Les réussites (Maurice, Rwanda, Botswana) prouvent que la \cle{gouvernance visionnaire} et l'\cle{investissement dans l'humain} sont les leviers internes, que le panafricanisme doit généraliser à l'échelle continentale.
\end{aretenir}

\coursec{Difficultés et limites du panafricanisme}

Après avoir analysé les objectifs ambitieux du panafricanisme — unité politique, intégration économique, solidarité culturelle et émancipation —, il est indispensable de confronter cet idéal à la réalité historique et contemporaine. Le chemin qui mène de la déclaration d'intention à la réalisation concrète est semé d'obstacles structurels, politiques et géostratégiques. Comprendre ces limites ne revient pas à nier la pertinence du projet, mais à en mesurer la complexité pour mieux penser les voies de son dépassement.

\courssub{L'héritage colonial : des frontières artificielles aux systèmes hétérogènes}

\paragraph{La Conférence de Berlin (1884-1885) : l'acte de naissance des États africains modernes.}
L'obstacle le plus fondamental est d'ordre territorial. Les frontières actuelles de l'Afrique ne sont pas le fruit d'une histoire longue de peuplement, de conquêtes ou de traités entre peuples africains. Elles ont été tracées à la règle et au compas, dans les salons de Berlin, par des puissances européennes qui ignoraient tout des réalités ethniques, linguistiques, culturelles et économiques du continent.

\begin{savaistu}[La Conférence de Berlin : un partage sans Africains]
La Conférence de Berlin (novembre 1884 – février 1885) a réuni 14 puissances européennes (Allemagne, France, Royaume-Uni, Portugal, Belgique, etc.) et les États-Unis comme observateur. \textbf{Aucun représentant africain n'était présent.} Le principe de l'\emph{effectivité} (occupation effective du terrain) a déclenché la « course au partage ». Résultat : environ \textbf{80\% des frontières africaines actuelles} sont des lignes droites (méridiens, parallèles) ou suivent des cours d'eau, coupant arbitrairement des espaces de vie séculaires.
\end{savaistu}

\begin{exemplebox}[Exemple concret : le peuple Bamiléké et la frontière Cameroun/Nigeria]
Les Bamiléké (Cameroun occidental) et les populations apparentées du côté nigérian (État d'Enugu, Cross River) partagent la même langue, les mêmes coutumes, les mêmes chefferies traditionnelles. La frontière tracée en 1919 (mandat Société des Nations) a séparé des familles, des chefferies et des marchés traditionnels. Aujourd'hui encore, les commerçants de Bafoussam ou Dschang doivent payer des droits de douane pour vendre sur le marché de la frontière, alors que leurs arrière-grands-parents y allaient librement.
\end{exemplebox}

\paragraph{Hétérogénéité des héritages administratifs et linguistiques.}
Cet héritage a produit trois grands blocs linguistiques et administratifs qui complètent l'émiettement territorial :
\begin{itemize}
    \item \textbf{Zone francophone :} Droit civil (Code Napoléon), administration centralisée, langue française, monnaie CFA (pour la plupart).
    \item \textbf{Zone anglophone :} Common Law, administration indirecte (Indirect Rule), langue anglaise, systèmes éducatifs et judiciaires distincts.
    \item \textbf{Zone lusophone :} Héritage portugais (Angola, Mozambique, Guinée-Bissau, Cap-Vert, São Tomé).
\end{itemize}
Cette triple fragmentation rend l'harmonisation des lois, des programmes scolaires, des normes techniques et des procédures douanières extraordinairement coûteuse et lente. L'intégration ne se décrète pas ; elle suppose une \emph{traduction} permanente entre des systèmes juridiques et administratifs incompatibles.

\begin{popfigure}
\begin{tikzpicture}[scale=0.75]
    % Contour très simplifié de l'Afrique (polygone approximatif)
    \draw[popdark, line width=1.5pt, fill=popcream] 
    (0,0) .. controls (2,1) and (4,2) .. (6,1) .. controls (8,0) and (9,-1) .. (8,-3) 
    .. controls (7,-4) and (5,-5) .. (3,-4) .. controls (1,-3) and (0,-2) .. (0,0);
    % Détail Corne de l'Est
    \draw[popdark, line width=1.5pt, fill=popcream] (6,1) .. controls (7,2) and (8,3) .. (9,2) .. controls (10,1) and (10,0) .. (9,-1) .. controls (8,0) and (7,1) .. (6,1);
    % Détail Afrique Australe
    \draw[popdark, line width=1.5pt, fill=popcream] (3,-4) .. controls (4,-5) and (5,-6) .. (4,-7) .. controls (2,-7) and (1,-6) .. (1,-5) .. controls (1,-4) and (2,-4) .. (3,-4);
    
    % Frontières rectilignes (exemples emblématiques)
    \draw[poporange, line width=1.2pt, dash pattern=on 4pt off 2pt] (1.5, -0.5) -- (4.5, -0.5); 
    \draw[poporange, line width=1.2pt, dash pattern=on 4pt off 2pt] (2.5, 0.5) -- (2.5, -2.5); 
    \draw[poporange, line width=1.2pt, dash pattern=on 4pt off 2pt] (5.5, 0.2) -- (5.5, -2.0); 
    \draw[poporange, line width=1.2pt, dash pattern=on 4pt off 2pt] (1.0, -2.0) -- (4.0, -2.0); 
    \draw[poporange, line width=1.2pt, dash pattern=on 4pt off 2pt] (4.0, -2.0) -- (4.0, -3.5); 
    \draw[poporange, line width=1.2pt, dash pattern=on 4pt off 2pt] (3.5, -4.5) -- (5.5, -4.5); 
    \draw[poporange, line width=1.2pt, dash pattern=on 4pt off 2pt] (5.5, -4.5) -- (5.5, -6.0); 
    
    % Légende visuelle
    \node[poporange, font=\bfseries\small] at (5.5, 1.5) {Frontières rectilignes (Berlin)};
    \node[popdark, font=\small] at (1, 1) {Méditerranée};
    \node[popdark, font=\small] at (8, -5) {Océan Indien};
    \node[popdark, font=\small] at (-1, -3) {Atlantique};
    
    % Quelques capitales pour repère
    \foreach \x/\y/\name in {2.5/0.2/Le Caire, 3.5/-1.5/Khartoum, 2.0/-2.5/Niamey, 4.5/-3.0/Windhoek} {
        \fill[popblue] (\x,\y) circle (2pt);
        \node[popblue, font=\tiny, anchor=west] at (\x+0.1,\y) {\name};
    }
    
    % Flèche Nord
    \draw[->, popdark, line width=1pt] (7.5, 2.5) -- (7.5, 3.2);
    \node[popdark, font=\small\bfseries] at (7.5, 3.5) {N};
\end{tikzpicture}
\legende{Carte schématique de l'Afrique mettant en évidence le tracé rectiligne (pointillés orange) de nombreuses frontières héritées de la Conférence de Berlin (1884-1885). Ces lignes droites (méridiens et parallèles) coupent arbitrairement des groupes ethniques, des bassins hydrographiques et des espaces économiques cohérents.}
\end{popfigure}

\courssub{Divergences idéologiques et égoïsmes nationaux : l'unité impossible ?}

\paragraph{De Casablanca à Monrovia : la fracture fondatrice.}
Dès l'accession aux indépendances (années 1960), le panafricanisme s'est heurté à une divergence stratégique majeure sur \emph{la méthode} et \emph{le rythme} de l'unité.

\begin{center}
\begin{tabularx}{\linewidth}{|l|X|X|}\hline
\rowcolor{popblueL}
\textbf{Dimension} & \textbf{Bloc de Casablanca (1961)} & \textbf{Bloc de Monrovia (1961)} \\ \hline
\textbf{Pays leaders} & Ghana (Nkrumah), Guinée (Touré), Mali (Keïta), Maroc, Algérie (GPRA), Égypte & Nigeria, Liberia, Sénégal, Côte d'Ivoire, Cameroun, Togo, etc. \\ \hline
\textbf{Vision} & \textbf{Union fédérale immédiate :} « \emph{United States of Africa} » maintenant. Gouvernement continental, armée unique, monnaie unique, citoyenneté unique. & \textbf{Intégration progressive :} Coopération fonctionnelle (économie, culture, technique) d'abord. Respect de la souveraineté nationale acquise. \\ \hline
\textbf{Philosophie} & \emph{« Cherchez d'abord le royaume politique, tout le reste vous sera donné par surcroît. »} (Nkrumah) & \emph{« L'unité ne se décrète pas, elle se construit. »} (Senghor, Houphouët-Boigny) \\ \hline
\textbf{Risque perçu} & Dilution des souverainetés naissantes, dictature continentale. & Balkanisation définitive, néocolonialisme durable. \\ \hline
\end{tabularx}
\end{center}

Cette opposition a abouti à un compromis : la création de l'\textbf{OUA (Organisation de l'Unité Africaine)} en 1963 à Addis-Abeba, sur la base de la Charte de Monrovia (souveraineté, non-ingérence, règlement pacifique des différends). L'idéal fédéral de Casablanca a été relégué au rang d'« objectif ultime ».

\begin{attention}[Piège à éviter : lire l'histoire comme une simple querelle de personnes]
Il ne s'agit pas seulement d'un conflit d'ego entre Nkrumah et Houphouët-Boigny. C'est un débat philosophique et stratégique profond : \textbf{la souveraineté nationale est-elle un obstacle à l'unité ou sa condition nécessaire ?} Une bonne dissertation doit montrer que ce dilemme structure encore les débats actuels au sein de l'Union Africaine (UA) : fédéralisme vs confédération, supranationalité vs intergouvernementalisme.
\end{attention}

\paragraph{Le principe de non-ingérence : bouclier des tyrans.}
L'article III de la Charte de l'OUA consacre le « respect de la souveraineté et de l'intégrité territoriale de chaque État » et la « non-ingérence dans les affaires intérieures ». Ce principe, né de la peur du néocolonialisme et du désir de protéger les jeunes États, est devenu un \textbf{paralysant majeur} :
\begin{itemize}
    \item Il a empêché l'OUA d'intervenir lors des génocides (Rwanda 1994), des guerres civiles (Liberia, Sierra Leone, Somalie) ou des dictatures sanglantes (Bokassa en Centrafrique, Amin Dada en Ouganda, Mobutu au Zaïre).
    \item Il a légitimé la « diplomatie du silence » entre chefs d'État : on ne critique pas son pair pour ne pas être critiqué à son tour.
\end{itemize}
Ce n'est qu'avec la transformation de l'OUA en \textbf{Union Africaine (UA) en 2002} (Acte constitutif de Lomé 2000) que le principe de « non-indifférence » (droit d'intervention en cas de crimes de guerre, génocide, crimes contre l'humanité) a été inscrit, mais son application reste sélective et politisée.

\courssub{Dépendances extérieures persistantes : le piège néocolonial}

\begin{definition}[Néocolonialisme]
Domination économique, politique, culturelle et militaire qu'exercent encore d'anciennes puissances coloniales ou de nouvelles grandes puissances sur des États africains formellement indépendants. Elle passe par le contrôle des matières premières, l'endettement, les accords de défense inégaux, l'influence monétaire (zone franc) et le conditionnement de l'aide au développement.
\end{definition}

\paragraph{Asymétrie commerciale et position de « fournisseur de matières premières ».}
L'intégration africaine bute sur la structure même des économies : elles sont \textbf{extraverties} (tournées vers l'exportation de matières premières brutes vers l'Europe, la Chine, les USA) et non \textbf{introverties} (échanges entre elles).
\begin{itemize}
    \item \textbf{Part du commerce intra-africain :} environ \textbf{15 à 18\%} du commerce total du continent (contre 60\% en Europe, 60\% en Asie, 45\% en Amérique du Nord).
    \item \textbf{Produits échangés :} mostly raw materials (pétrole, minerais, cacao, coton) ou produits semi-finis. Très peu de biens manufacturés à forte valeur ajoutée.
\end{itemize}
Tant que le Cameroun exporte du bois brut vers la Chine et importe des meubles de Chine, et que la Côte d'Ivoire exporte du cacao vers l'Europe et importe du chocolat d'Europe, l'intégration régionale reste un vœu pieux. La \textbf{ZLECAf (Zone de Libre-Échange Continentale Africaine)}, entrée en vigueur en 2021, vise précisément à inverser cette logique, mais les obstacles non-tarifaires (normes, corruption aux frontières, infrastructures) sont immenses.

\paragraph{Endettement et conditionnalités.}
La dette publique africaine (environ 650 milliards USD en 2023) contraint les États à négocier avec le FMI et la Banque Mondiale. Les programmes d'ajustement structurel (PAS) des années 80-90 ont démantelé les services publics et les industries naissantes. Aujourd'hui, la dette auprès de la Chine (infrastructures contre ressources) crée de nouvelles dépendances. Un État endetté n'est pas un État souverain capable de mener une politique panafricaine hardie.

\courssub{Conflits, instabilité et déficit d'infrastructures : le corps malade de l'unité}

\paragraph{L'instabilité politique : un frein structurel.}
Coups d'État (Mali, Guinée, Burkina Faso, Niger, Gabon récemment), guerres civiles (Soudan, RDC, Somalie), terrorisme (Sahel, Lac Tchad, Mozambique) : l'insécurité détruit le capital humain, physique et la confiance nécessaire à l'intégration.
\begin{itemize}
    \item Les flux de réfugiés déstabilisent les pays voisins.
    \item Les budgets militaires absorbent les ressources de l'éducation et de la santé.
    \item L'État de droit recule, décourageant l'investissement privé intra-africain.
\end{itemize}

\paragraph{Le déficit infrastructurel : l'Afrique mal connectée à elle-même.}
C'est l'obstacle le plus tangible pour les populations et les entrepreneurs.

\begin{savaistu}[Le paradoxe des liaisons aériennes]
Il est souvent \textbf{plus rapide et moins cher} de voyager de \textbf{Douala (Cameroun) à Paris (France)} (vol direct ~6h30, liaisons quotidiennes) que de Douala à \textbf{N'Djamena (Tchad)} ou \textbf{Bangui (RCA)} (pas de vol direct, correspondance via Paris, Istanbul, Addis-Abeba ou Nairobi ; durée 12h à 24h).
\textbf{Cause :} L'héritage colonial a orienté les infrastructures (rails, routes, ports, aérien) vers les métropoles (extraction), pas vers les voisins (intégration). Le ciel africain reste fragmenté malgré la décision de Yamoussoukro (1999) pour un ciel unique africain.
\end{savaistu}

\begin{experience}[Enquête de terrain : le corridor Douala-Bangui-N'Djamena]
\textbf{Problématique :} Pourquoi le transport de marchandises entre la zone CEMAC (Communauté Économique et Monétaire de l'Afrique Centrale) coûte-t-il 3 à 4 fois plus cher qu'entre l'Europe et l'Afrique ?
\textbf{Observations (simulées) :}
\begin{enumerate}
    \item \textbf{Route :} Piste dégradée (Ketté - Batouri - Garoua-Boulaï - Bangui), barrages routiers multiples (police, douane, gendarmerie, eaux et forêts), rackets quotidiens (~50 000 à 100 000 FCFA par camion par trajet).
    \item \textbf{Rail :} Pas de liaison ferroviaire continue (Camrail s'arrête à Ngaoundéré ; pas de rail en RCA ni au Tchad).
    \item \textbf{Fleuve :} Oubangui non navigable en saison sèche ; port de Bangui vétuste.
    \item \textbf{Aérien :} Fret inexistant ou hors de prix.
\end{enumerate}
\textbf{Interprétation :} L'intégration physique est la condition \emph{sine qua non} de l'intégration économique. Sans routes, rails, énergie et numérique transfrontaliers, la ZLECAf reste un texte juridique sans effet réel.
\end{experience}

\courssub{Bilan nuancé : échecs réels, avancées concrètes, horizon pertinent}

Il serait intellectuellement malhonnête et pédagogiquement contre-productif de dresser un bilan purement négatif. La dialectique hégélienne (Thèse - Antithèse - Synthèse) s'impose ici.

\begin{methode}[Construire un bilan argumenté (Thèse - Antithèse - Synthèse)]
\textbf{Étape 1 : La Thèse (Les acquis / Réussites).} Lister les institutions créées, les normes adoptées, les avancées réelles.
\textbf{Étape 2 : L'Antithèse (Les échecs / Limites).} Lister les obstacles structurels, les promesses non tenues, les souffrances des peuples.
\textbf{Étape 3 : La Synthèse (Analyse nuancée et personnelle).} Expliquer \emph{pourquoi} les acquis n'ont pas résolu les problèmes (gap ratification/application, manque de volonté politique, contexte géopolitique). Conclure sur la pertinence \emph{aujourd'hui} du panafricanisme : mythe mobilisateur ou projet politique opératoire ?
\end{methode}

\begin{popfigure}
\begin{tikzpicture}
    % Tableau deux colonnes : Acquis / Difficultés
    \node[anchor=north west, inner sep=0, align=center] (table) at (0,0){
    \begin{tabularx}{\linewidth}{|>{\bfseries}X|X|}
    \hline
    \rowcolor{popgreenL}
    \multicolumn{2}{|c|}{\textbf{BILAN DU PANAFRICANISME : ACQUIS ET DIFFICULTÉS PERSISTANTES}} \\
    \hline
    \rowcolor{popgreen!20}
    \textbf{ACQUIS MAJEURS (Thèse)} & \textbf{DIFFICULTÉS PERSISTANTES (Antithèse)} \\
    \hline
    \textbf{1. Existence institutionnelle :} OUA (1963) $\rightarrow$ UA (2002). Siège à Addis-Abeba. Parlement panafricain (Midrand). Cour africaine des droits de l'homme (Arusha). & \textbf{1. Faible supranationalité :} Décisions non contraignantes (sauf Cour, peu saisie). Budget UA dépendant à \textasciitilde60\% de partenaires extérieurs (UE, Chine). \\
    \hline
    \textbf{2. Cadre juridique :} Traité d'Abuja (1991) : 6 étapes vers la Communauté Économique Africaine. Protocole sur la libre circulation (2018). ZLECAf (2021) : plus grande ZLE au monde (1,3 Md hab, 3,4 \$ PIB). & \textbf{2. Gap ratification/application :} 54 signataires ZLECAf, mais ratifications lentes. Protocole libre circulation : 4 ratifications seulement (sur 15 nécessaires). Barrières non-tarifaires massives. \\
    \hline
    \textbf{3. Paix et sécurité :} UAPS (Architecture de paix et sécurité). Forces en attente (brigades régionales). Interventions : AMISOM (Somalie), MISCA (RCA), G5 Sahel (soutenu par UA). & \textbf{3. Réactivité limitée :} Dépendance logistique/financière extérieures. Non-ingérence encore frein. Coups d'État récurrents (belt sahélien). Terrorisme en expansion. \\
    \hline
    \textbf{4. Identité et culture :} Kiswahili langue de travail UA. Année de la culture. Passeport africain (lancé 2016). Journée de l'Afrique (25 mai). & \textbf{4. Citoyenneté virtuelle :} Passeport non reconnu par la plupart des États. Visa encore requis entre la majorité des pays africains. Xénophobie (Afrique du Sud, Gabon, Guinée Équatoriale). \\
    \hline
    \textbf{5. Développement :} NEPAD (Nouveau Partenariat pour le Développement de l'Afrique). Agenda 2063 (« L'Afrique que nous voulons »). PIDA (Programme Infrastructures). & \textbf{5. Développement inégal :} Croissance souvent sans emploi, sans transformation structurelle. Pauvreté extrême concentrée en Afrique. Changement climatique non causé par l'Afrique mais subi. \\
    \hline
    \end{tabularx}
    };
\end{tikzpicture}
\legende{Tableau de synthèse dialectique : les acquis institutionnels et juridiques indéniables (colonne gauche) se heurtent à un fossé d'application, de financement propre et de volonté politique réelle (colonne droite). Source : Rapports UA, CEA, AfDB, analyses croisées.}
\end{popfigure}

\courssub{Synthèse : L'unité africaine, horizon indépassable ou chimère ?}

Le panafricanisme n'a pas « échoué » ; il est \textbf{inachevé}. Il a réussi l'étape politique (indépendance, existence d'une organisation continentale), mais bute sur l'étape économique et fédérale.

\begin{aretenir}[Synthèse pour le Probatoire / Baccalauréat]
\begin{enumerate}
    \item \textbf{Les causes structurelles :} Frontières de Berlin (émiettement), économies extraverties (néocolonialisme), hétérogénéité juridico-linguistique.
    \item \textbf{Les causes politiques :} Égoïsmes nationaux, non-ingérence (OUA), instabilité/coups d'État, manque de leadership visionnaire partagé.
    \item \textbf{Les causes physiques :} Infrastructures coloniales (radiales vers l'extérieur), coût des transports intra-africains prohibitif.
    \item \textbf{Le bilan dialectique :} Des institutions existent (UA, ZLECAf, Cour, Parlement) mais manquent de souveraineté effective (budget, pouvoir contraignant, armée unique).
    \item \textbf{Pertinence actuelle :} Le panafricanisme reste l'\textbf{unique horizon stratégique} pour peser dans la mondialisation (démographie, ressources, climat). La ZLECAf est le test décisif : si le commerce intra-africain décolle, l'intégration politique suivra. Sinon, le risque de « balkanisation » ou de recolonisation rampante persiste.
\end{enumerate}
\end{aretenir}

\begin{exoresolu}[Sujet type Bac : « Le panafricanisme est-il encore une philosophie pertinente pour le développement de l'Afrique ? »]
\textbf{Analyse du sujet :} « Encore » $\rightarrow$ histoire, bilan. « Philosophie pertinente » $\rightarrow$ capacité à penser/guider l'action pour le développement. « Développement » $\rightarrow$ pas seulement croissance PIB, mais autonomie, bien-être, souveraineté.
\textbf{Plan dialectique :}
\textbf{I. Thèse : Une philosophie fondatrice et encore mobilisatrice (Pertinence idéelle).}
\begin{itemize}
    \item A. Héritage de la décolonisation : a donné sens à l'indépendance (Nkrumah, Nyerere, Sankara).
    \item B. Réponse à la mondialisation : Seule l'échelle continentale permet de négocier (OMC, Climat, Dette), de mutualiser les marchés (ZLECAf), de gérer les bassins fluviaux (Niger, Congo, Nil).
    \item C. Agenda 2063 : Vision long terme endogène (transformation structurelle, intégration, paix).
\end{itemize}
\textbf{II. Antithèse : Des limites structurelles qui en font une « philosophie sans puissance » (Pertinence opératoire faible).}
\begin{itemize}
    \item A. Primauté de l'État-nation postcolonial : les élites captent la rente, craignent la supranationalité.
    \item B. Dépendance externe persistante : Aide, dette, sécurité, technologie $\rightarrow$ perte d'autonomie décisionnelle.
    \item C. Faiblesse de la société civile panafricaine : Peuples peu consultés, frontières fermées, xénophobie.
\end{itemize}
\textbf{III. Synthèse : Le panafricanisme comme \emph{projet régulateur} (Kant) et chantier permanent.}
\begin{itemize}
    \item Il n'est pas une idéologie close (comme le marxisme-léninisme d'État), mais un \emph{idéal régulateur} : on ne l'atteint jamais tout à fait, mais il oriente l'action.
    \item La ZLECAf est le test de vérité : passer du « panafricanisme des chefs d'État » au « panafricanisme des peuples et des entrepreneurs ».
    \item Condition de succès : \textbf{réappropriation par le bas} (éducation panafricaine, médias, mobilité étudiante, paiement panafricain PAPSS) + \textbf{réforme institutionnelle} (budget UA 100\% africain, suffrage universel pour Parlement, Cour effective).
\end{itemize}
\textbf{Conclusion :} Oui, pertinent, \emph{à condition de cesser d'être une rhétorique de sommet pour devenir une ingénierie quotidienne de l'intégration par le bas.}
\end{exoresolu}

\begin{exercice}[Entraînement : Commentaire de texte / Dissertation]
\textbf{Sujet :} « L'Afrique n'a pas besoin de charity, elle a besoin de justice. » (Discours d'Emmanuel Macron à Ouagadougou, 2017, reprenant une idée panafricaine classique).
\textbf{Consignes :} En vous appuyant sur la leçon « Panafricanisme et développement », expliquez en quoi cette phrase résume le passage du néocolonialisme (aide/charité) à la souveraineté panafricaine (justice/commerce équitable/intégration). Montrez les obstacles actuels à cette justice.
\end{exercice}

\begin{corrige}[Exercice n°1]
\textbf{Introduction :} Accroche sur la relation Afrique-France/Monde. Définition : Néocolonialisme = domination déguisée en aide. Panafricanisme = quête d'autonomie. Problématique : Le passage de la charité (relation verticale, humiliante) à la justice (relation horizontale, contractuelle) est-il le cœur du projet panafricain de développement ?
\textbf{Développement :}
\begin{itemize}
    \item \textbf{1. La charité comme outil néocolonial :} Aide conditionnée (gouvernance, ouverture marchés), dette odieuse, accords de défense inégaux, franc CFA (garantie Trésor français). Crée dépendance, pas développement.
    \item \textbf{2. La justice comme horizon panafricain :} Justice commerciale (ZLECAf, fin barrières), justice monétaire (monnaie unique africaine ?), justice climatique (pollueur-payeur, fonds pertes/dommages), justice mémorielle (restitution biens culturels).
    \item \textbf{3. Obstacles à cette justice :} Intérêts divergents des 54 États, faiblesse négociatrice face aux blocs (UE, Chine, USA), corruption élites, manque d'industrialisation (pas de levier).
\end{itemize}
\textbf{Conclusion :} La phrase de Macron (paradoxe) illustre la prise de conscience. Mais la justice ne se mendie pas, elle se conquiert par l'unité effective (ZLECAf, PIDA, UA forte). Le panafricanisme est la philosophie de cette conquête.
\end{corrige}

\coursec{Synthèse : l'unité africaine est-elle la voie du développement ?}

La section précédente a montré que le panafricanisme rencontre de réelles difficultés : rivalités entre États, faiblesse des institutions continentales, dépendance économique extérieure, diversité linguistique et politique. Faut-il en conclure que le projet panafricain est un échec ? Ce serait aller trop vite. Les difficultés d'un projet ne prouvent pas que son idéal est faux ; elles montrent seulement que sa réalisation exige des conditions précises. Il est temps de rassembler les acquis de la leçon et de répondre à la question philosophique qui la traverse : \emph{le panafricanisme peut-il être une philosophie du développement de l'Afrique ?}

\courssub{Rappel de la problématique}

\begin{definition}[Philosophie du développement]
Une \cle{philosophie du développement} est une réflexion d'ensemble sur les \cle{finalités}, les \cle{moyens} et les \cle{conditions} d'un développement authentique de l'Afrique. Elle ne se contente pas de mesurer la richesse produite : elle demande \emph{pour qui}, \emph{pour quoi} et \emph{comment} un peuple se développe.
\end{definition}

Poser la question « l'unité africaine est-elle la voie du développement ? », c'est donc demander si le panafricanisme réussit à articuler trois choses : une fin (la dignité et la prospérité des Africains), un moyen (l'union politique et économique du continent) et des conditions (gouvernance, éducation, paix). La dissertation exige de discuter chacun de ces trois termes, et non de réciter l'histoire de l'OUA et de l'Union africaine.

\courssub{Argument 1 : oui, l'unité renforce le poids économique et politique du continent}

L'intuition est simple : un seul doigt se casse facilement, cinq doigts fermés forment un poing. Séparés, les États africains pèsent peu dans les négociations internationales ; unis, ils représentent un marché de plus d'un milliard de consommateurs et un ensemble de ressources considérables.

\begin{exemplebox}[La ZLECAf : l'unité en actes]
La \cle{Zone de libre-échange continentale africaine} (ZLECAf), entrée en vigueur en 2021, crée un marché unique pour les biens et services. Selon la Banque mondiale, elle pourrait augmenter les revenus du continent de plusieurs centaines de milliards de dollars d'ici 2035 et sortir des dizaines de millions de personnes de l'extrême pauvreté. Les négociations collectives (sur les prix des matières premières, les accords commerciaux, le climat) donnent à l'Afrique une voix qu'aucun État isolé ne pourrait avoir.
\end{exemplebox}

Philosophiquement, cet argument prolonge la pensée de Kwame \cle{Nkrumah} : « L'Afrique doit s'unir ou périr », car aucun pays africain ne peut être indépendant et prospère seul. L'unité apparaît comme la \emph{condition de possibilité} de la souveraineté économique réelle.

\courssub{Argument 2 : l'unité seule ne suffit pas}

Mais l'unité n'est pas une baguette magique. Un continent uni mais mal gouverné resterait pauvre. L'expérience montre que le développement exige aussi :

\begin{itemize}
\item la \cle{bonne gouvernance} : lutte contre la corruption, justice équitable, gestion transparente des ressources publiques ;
\item l'\cle{éducation} : former les jeunes aux sciences, aux techniques et à l'esprit d'entreprise, car aucun développement durable ne repose sur une population sans qualifications ;
\item la \cle{transformation locale} des ressources : vendre du chocolat plutôt que des fèves de cacao, des meubles plutôt que des grumes de bois.
\end{itemize}

Autrement dit, l'unité est une condition \emph{nécessaire} mais non \emph{suffisante} : elle multiplie les forces, mais elle ne remplace ni le travail, ni la compétence, ni l'honnêteté.

\courssub{Argument 3 : le développement véritable est d'abord un choix de société}

Au fond, la question dépasse l'économie. Amartya \cle{Sen} définit le développement comme un \emph{processus d'élargissement des libertés réelles} des personnes : vivre longtemps, être instruit, participer à la vie publique. Sous cet angle, le développement n'est pas seulement une affaire de revenus, mais de \cle{capacités} et de \cle{dignité}.

C'est ici que le panafricanisme retrouve toute sa profondeur philosophique. Pour Nkrumah, l'unité doit rendre à l'Africain la maîtrise de son destin ; pour Léopold Sédar \cle{Senghor}, la renaissance africaine suppose de puiser dans les valeurs propres du continent — solidarité, communauté, dialogue des cultures — plutôt que de copier servilement des modèles étrangers. Le développement authentique est donc \cle{endogène} : il part des ressources, des besoins et des valeurs de la société africaine elle-même.

\begin{aretenir}[La thèse de synthèse]
Le panafricanisme est bien une philosophie du développement de l'Afrique, à condition de ne pas réduire l'unité à un slogan : elle doit s'incarner dans des institutions efficaces (ZLECAf, unions régionales), s'accompagner de bonne gouvernance et d'éducation, et viser un développement endogène fondé sur la dignité et les libertés réelles des Africains.
\end{aretenir}

\courssub{Application aux situations concrètes camerounaises}

\begin{exemplebox}[Le Cameroun et l'intégration africaine]
\begin{itemize}
\item \textbf{Transformation du cacao} : le Cameroun, cinquième producteur mondial, exporte encore l'essentiel de ses fèves brutes. Des unités locales de transformation (chocolateries à Yaoundé et Douala) créent plus d'emplois et de valeur ajoutée : c'est le développement endogène en pratique.
\item \textbf{Corridors CEMAC} : les corridors Douala--N'Djamena et Douala--Bangui relient le port de Douala à l'hinterland ; l'intégration régionale fait du Cameroun la porte de l'Afrique centrale.
\item \textbf{Entrepreneuriat des jeunes} : les start-up numériques camerounaises (paiement mobile, agriculture, santé) montrent qu'une jeunesse formée crée elle-même ses emplois.
\item \textbf{Diaspora} : les transferts d'argent et de compétences des Camerounais de l'étranger financent écoles, cliniques et entreprises : la diaspora est un prolongement vivant du panafricanisme.
\end{itemize}
\end{exemplebox}

\begin{savaistu}[Le Cameroun, carrefour de l'intégration africaine]
Le Cameroun a accueilli en 2021 des rencontres continentales sur l'intégration africaine, et Yaoundé abrite le siège de la Banque des États de l'Afrique centrale (BEAC). La \cle{CEMAC} regroupe six États d'Afrique centrale (Cameroun, Congo, Gabon, Guinée équatoriale, Tchad, Centrafrique) autour d'une monnaie commune, le franc CFA, et d'institutions communes : un exemple concret, quoique perfectible, d'unité régionale.
\end{savaistu}

\courssub{Méthode de la dissertation}

\begin{methode}[Rédiger une introduction]
\begin{enumerate}
\item \textbf{Accroche} : un fait ou une citation qui met en scène le problème (ex. : la création de la ZLECAf en 2021).
\item \textbf{Définitions} : définir les termes du sujet (panafricanisme, développement).
\item \textbf{Problématique} : formuler la question philosophique sous forme interrogative.
\item \textbf{Annonce du plan} : présenter les deux ou trois parties du développement.
\end{enumerate}
\end{methode}

\begin{methode}[Rédiger et justifier une conclusion]
\begin{enumerate}
\item \textbf{Répondre explicitement} à la problématique posée dans l'introduction (oui, non, ou réponse nuancée et justifiée).
\item \textbf{Synthétiser} les arguments du développement en deux ou trois phrases, sans les répéter mot pour mot.
\item \textbf{Ouvrir} sur une question voisine sans hors-sujet (ex. : « Reste à savoir si les jeunes générations sauront transformer cet idéal en réalité vécue »).
\end{enumerate}
\end{methode}

\begin{attention}[Erreur fréquente]
Réciter l'histoire du panafricanisme (congrès, dates, personnalités) \emph{sans répondre à la question philosophique posée}. L'histoire n'est qu'un matériau : elle doit servir une argumentation. Chaque fait cité doit être suivi de sa portée philosophique (« cela montre que... »).
\end{attention}

\begin{popfigure}
\begin{tikzpicture}[node distance=0.55cm, every node/.style={font=\small}]
\node[draw=popblue, fill=popblueL, rounded corners, minimum width=4.2cm, minimum height=0.8cm] (a) {Accroche};
\node[draw=popblue, fill=popblueL, rounded corners, minimum width=4.2cm, minimum height=0.8cm, below=of a] (b) {Définitions des termes};
\node[draw=popblue, fill=popblueL, rounded corners, minimum width=4.2cm, minimum height=0.8cm, below=of b] (c) {Problématique};
\node[draw=popblue, fill=popblueL, rounded corners, minimum width=4.2cm, minimum height=0.8cm, below=of c] (d) {Annonce du plan};
\node[draw=poporange, fill=poporangeL, rounded corners, minimum width=4.2cm, minimum height=0.8cm, below=of d] (e) {Arguments développés};
\node[draw=popgreen, fill=popgreenL, rounded corners, minimum width=4.2cm, minimum height=0.8cm, below=of e] (f) {Conclusion justifiée};
\draw[-{Stealth}, thick, popdark] (a) -- (b);
\draw[-{Stealth}, thick, popdark] (b) -- (c);
\draw[-{Stealth}, thick, popdark] (c) -- (d);
\draw[-{Stealth}, thick, popdark] (d) -- (e);
\draw[-{Stealth}, thick, popdark] (e) -- (f);
\node[right=0.4cm of a, text=popmuted, font=\footnotesize\itshape] {Introduction};
\node[right=0.4cm of e, text=popmuted, font=\footnotesize\itshape] {Développement};
\end{tikzpicture}
\legende{Schéma de la démarche de dissertation : de l'accroche à la conclusion, chaque étape prépare la suivante.}
\end{popfigure}

\begin{popfigure}
\begin{tikzpicture}
\node[font=\small, align=center]{
\begin{tabular}{|l|l|l|}
\hline
\rowcolor{popblueL} \textbf{Concept / Institution} & \textbf{Auteur ou date} & \textbf{Idée clé} \\
\hline
Développement & A. Sen & Élargissement des libertés réelles \\
\hline
Sous-développement & --- & Pauvreté, dépendance, inégalités \\
\hline
Panafricanisme & Nkrumah, Senghor & Unité et dignité de l'Afrique \\
\hline
OUA & 1963 & Première organisation continentale \\
\hline
Union africaine & 2002 & Relance de l'intégration \\
\hline
ZLECAf & 2021 & Marché unique africain \\
\hline
CEMAC & 1994 & Intégration régionale, franc CFA \\
\hline
\end{tabular}};
\end{tikzpicture}
\legende{Tableau récapitulatif de la leçon : concepts clés, auteurs, dates et institutions du panafricanisme.}
\end{popfigure}

\courssub{Entraînement : rédaction guidée}

\begin{exoresolu}[Introduction et conclusion rédigées]
Sujet : \emph{Le panafricanisme est-il la voie du développement de l'Afrique ?} Rédiger l'introduction et la conclusion de la dissertation.
\tcblower
\textbf{Introduction.} En 2021, l'entrée en vigueur de la ZLECAf a relancé un vieux rêve : celui d'une Afrique unie et prospère. Le \cle{panafricanisme} désigne le mouvement politique et culturel qui promeut l'unité des peuples africains ; le \cle{développement} désigne le processus par lequel une société améliore durablement les conditions de vie et les libertés de ses membres. Mais l'unité politique et économique suffit-elle à produire un développement authentique ? Le panafricanisme est-il la voie du développement de l'Afrique ? Nous montrerons d'abord que l'unité renforce le poids du continent, puis qu'elle ne suffit pas sans gouvernance et éducation, enfin que le développement véritable est un choix de société fondé sur la dignité.

\textbf{Conclusion.} Le panafricanisme est bien une voie du développement de l'Afrique, car l'unité donne au continent la force économique et politique qui manque à chacun de ses États isolés. Mais cette voie n'est pas un chemin tout tracé : elle exige la bonne gouvernance, l'éducation des jeunes et la transformation locale des ressources, et elle n'a de sens que si elle sert la dignité et les libertés réelles des Africains. Reste à savoir si les générations actuelles sauront transformer cet idéal philosophique en réalité vécue.
\end{exoresolu}

\begin{exercice}[Pour s'entraîner]
Sujet : \emph{Faut-il être unis pour se développer ?} 1) Définis les deux termes du sujet. 2) Formule une problématique. 3) Propose un plan en deux parties avec un argument et un exemple camerounais pour chaque partie.
\end{exercice}

\begin{corrige}[Exercice]
1) « Unis » : solidaires, organisés ensemble au niveau régional ou continental. « Se développer » : améliorer durablement les conditions de vie et les libertés. 2) Problématique : l'unité est-elle une condition nécessaire et suffisante du développement ? 3) Plan possible : I. L'unité est nécessaire (argument : poids économique ; exemple : ZLECAf et corridors CEMAC au départ de Douala). II. Mais elle n'est pas suffisante (argument : sans gouvernance ni éducation, l'unité reste vide ; exemple : transformation locale du cacao comme complément indispensable).
\end{corrige}

\newpage

\coursec{Activité d'intégration}

\coursec{Leçon 09 : Panafricanisme et développement de l'Afrique}

\courssub{Objectifs de l'activité}
\begin{itemize}
    \item Mobiliser les concepts de \cle{développement}, de \cle{sous-développement} et d'\cle{unité africaine} pour analyser des documents authentiques.
    \item Construire une représentation cartographique synthétique des grands ensembles régionaux africains et des flux commerciaux.
    \item Exploiter des données statistiques pour comparer le commerce intra-africain et le commerce extra-africain.
    \item Rédiger un paragraphe argumenté structuré répondant à une question philosophique, en croisant des données factuelles et des références théoriques du cours.
\end{itemize}

\courssub{Situation}
En cette année 2025, le Cameroun accueille le sommet extraordinaire de la \cle{CEMAC} (Communauté Économique et Monétaire de l'Afrique Centrale) à Yaoundé. Le thème retenu est : « \emph{Accélération de la ZLECAf : levier pour l'industrialisation et la résilience économique de l'Afrique centrale} ». En marge de ce sommet, votre classe de Terminale Technique est invitée à participer à un « \emph{Forum des Jeunes Penseurs de l'Intégration} ». 

Votre mission, en tant qu'experts juniors, est de produire une \textbf{note de synthèse illustrée} destinée à alimenter le débat public. Cette note doit répondre à la question centrale : \textbf{« L'unité africaine est-elle la condition \emph{sine qua non} du développement du Cameroun et du continent ? »}

Pour ce faire, le comité d'organisation met à votre disposition un \textbf{dossier documentaire} composé des pièces suivantes :

\begin{enumerate}
    \item \textbf{Carte muette de l'Afrique} (format A3) avec les frontières des 54 États reconnus par l'UA.
    \item \textbf{Tableau 1 : Données socio-économiques de 10 pays africains (Années 2023-2024)}.
    \item \textbf{Tableau 2 : Structure des échanges commerciaux (Exportations + Importations) en milliards USD (2023)}.
    \item \textbf{Extrait 1 : Discours de Kwame Nkrumah, Conférence des Chefs d'État africains, Addis-Abeba, mai 1963}. 
    \item \textbf{Extrait 2 : Préambule de la Charte de l'Organisation de l'Unité Africaine (OUA), 25 mai 1963}.
    \item \textbf{Photographies} : (a) Autoroute Yaoundé-Douala (tronçon financé par la BEAC/BADEA) ; (b) Marché transfrontalier de Kousseri (frontière Cameroun-Tchad) ; (c) Port autonome de Kribi (phase 1, exploitation chinoise) ; (d) Siège de la Commission de l'UA à Addis-Abeba.
\end{enumerate}

\textbf{Tableau 1 : Indicateurs de développement (Sources : Banque Mondiale, PNUD, BCEAO/BEAC 2024)}

\begin{center}
\begin{tabularx}{\linewidth}{|l|X|X|X|X|X|}\hline
\textbf{Pays} & \textbf{PIB/hab (PPA \$)} & \textbf{IDH (2022)} & \textbf{Taux de pauvreté (\% < \$2,15/j)} & \textbf{Part industrie (\% PIB)} & \textbf{Principaux partenaires commerciaux} \\ \hline
Cameroun & 4\,100 & 0,576 & 23,0 & 26,5 & UE, Chine, Nigéria, pays CEMAC \\ \hline
Nigéria & 5\,400 & 0,548 & 30,9 & 30,8 & Inde, Espagne, Pays-Bas, CEDEAO \\ \hline
Côte d'Ivoire & 6\,200 & 0,550 & 18,2 & 21,4 & UE, Chine, CEDEAO \\ \hline
Sénégal & 4\,300 & 0,517 & 30,5 & 23,1 & UE, Chine, Mali, CEDEAO \\ \hline
Gabon & 16\,500 & 0,706 & 3,4 & 45,2 & Chine, UE, pays CEMAC \\ \hline
Tchad & 1\,500 & 0,394 & 42,3 & 14,8 & UE, Chine, Cameroun, pays CEMAC \\ \hline
RDC & 1\,300 & 0,481 & 62,0 & 41,5 & Chine, UE, Zambie, SADC \\ \hline
Kenya & 6\,500 & 0,601 & 27,0 & 17,6 & UE, Ouganda, Tanzanie, EAC \\ \hline
Afrique du Sud & 13\,400 & 0,713 & 18,9 & 25,2 & Chine, UE, USA, SADC \\ \hline
Maroc & 9\,100 & 0,683 & 1,7 & 25,8 & UE, Espagne, France, CEDEAO (observateur) \\ \hline
\end{tabularx}
\end{center}

\textbf{Tableau 2 : Commerce total (Export + Import) en Milliards USD (Source : CNUCED, Rapport 2024)}

\begin{center}
\begin{tabularx}{\linewidth}{|l|X|X|X|}\hline
\textbf{Zone / Pays} & \textbf{Commerce Intra-Africain} & \textbf{Commerce Extra-Africain} & \textbf{Taux d'intra-africain (\%)} \\ \hline
Cameroun & 1,8 & 10,2 & 15,0 \\ \hline
Nigéria & 5,5 & 95,0 & 5,5 \\ \hline
Côte d'Ivoire & 3,2 & 18,8 & 14,5 \\ \hline
CEDEAO (15 pays) & 45,0 & 255,0 & 15,0 \\ \hline
CEMAC (6 pays) & 4,5 & 35,5 & 11,2 \\ \hline
SADC (16 pays) & 65,0 & 235,0 & 21,7 \\ \hline
\textbf{Total Continent (54 pays)} & \textbf{190,0} & \textbf{1\,010,0} & \textbf{15,8} \\ \hline
\end{tabularx}
\end{center}

\textbf{Extrait 1 (Nkrumah, 1963) :} \emph{« Nous devons nous unir ou périr... Nous avons besoin d'une union économique pour éviter d'être les victimes du néocolonialisme. L'indépendance politique ne suffit pas si elle n'est pas adossée à une indépendance économique. Un marché commun africain, une monnaie unique, une banque centrale, une politique de défense commune, une citoyenneté africaine : voilà les piliers de l'État-Uni d'Afrique. »}

\textbf{Extrait 2 (Charte OUA, 1963) :} \emph{« Convaincus que la Charte des Nations Unies et la Déclaration universelle des droits de l'homme constituent le fondement d'une coopération pacifique et positive entre les États... Animés de la volonté de promouvoir l'unité et la solidarité des États africains... Déterminés à sauvegarder et à consolider l'indépendance, la souveraineté et l'intégrité territoriale des États africains... »}

\courssub{Consignes}
Travaillant en groupes de 4 à 5 élèves, vous réaliserez les trois tâches suivantes. Chaque groupe rendra une production unique (carte, graphique, paragraphe) qui sera présentée au tableau lors de la restitution.

\begin{enumerate}
    \item \textbf{Cartographie raisonnée (30 min) :} Sur la carte muette fournie :
    \begin{itemize}
        \item Localisez et colorez distinctement les 4 grands ensembles régionaux suivants : \cle{CEDEAO} (Afrique de l'Ouest), \cle{CEMAC} (Afrique Centrale), \cle{SADC} (Afrique Australe), \cle{UEMOA} (Zone Franc Ouest). Indiquez la capitale de chaque organisation.
        \item Tracez des \textbf{flèches proportionnelles} (épaisseur = volume) représentant les 3 principaux flux commerciaux \emph{extra-africains} (vers UE, Chine, USA/Inde) et les 2 principaux flux \emph{intra-africains} identifiables depuis le Cameroun (vers Nigéria, vers Tchad/CEMAC).
        \item Légendez votre carte (Titre, Auteur, Date, Échelle indicative, Orientation, Signes conventionnels).
    \end{itemize}
    
    \item \textbf{Analyse statistique et graphique (20 min) :}
    \begin{itemize}
        \item À partir du \textbf{Tableau 2}, construisez un \textbf{diagramme en barres groupées} comparant, pour le Cameroun, la CEDEAO, la CEMAC, la SADC et le Total Continental, les volumes du commerce Intra-Africain vs Extra-Africain.
        \item Calculez le \textbf{taux de couverture du commerce intra-africain par le commerce extra-africain} pour le Cameroun (Rapport Extra / Intra). Interprétez ce ratio en une phrase.
    \end{itemize}
    
    \item \textbf{Rédaction argumentée (40 min) :} Rédigez un \textbf{paragraphe unique d'environ 250 mots} répondant à la question : \emph{« L'unité africaine est-elle la condition du développement du Cameroun et du continent ? »}
    \begin{itemize}
        \item Votre réponse doit comporter une \textbf{thèse claire} (oui, non, ou nuancée), au moins \textbf{deux arguments distincts} appuyés sur des \textbf{données chiffrées} (Tableaux 1 ou 2) et sur \textbf{au moins un extrait de texte} (Nkrumah ou Charte OUA).
        \item Vous devez mobiliser \textbf{explicitement la définition du cours} du \cle{développement} (processus d'amélioration durable du bien-être) ou du \cle{sous-développement} (structure de dépendance, extraversion).
        \item La conclusion du paragraphe doit ouvrir sur une perspective (ZLECAf, fédéralisme, coopération technique).
    \end{itemize}
\end{enumerate}

\courssub{Ressources à mobiliser}
\begin{definition}[Développement (rappel de cours)]
Processus global et durable d'amélioration des conditions de vie d'une population (économiques, sociales, culturelles, politiques), mesuré par des indicateurs quantitatifs (PIB/hab, IDH, taux d'industrialisation) et qualitatifs (libertés, équité, autonomie). Le développement suppose une \cle{transformation structurelle} de l'économie (passage du primaire au secondaire/tertiaire) et une \cle{maîtrise endogène} des choix stratégiques.
\end{definition}

\begin{definition}[Sous-développement (rappel de cours)]
Situation structurelle caractérisée par la \cle{dépendance} (extériorisation de l'économie : exportation de matières premières, importation de produits finis), la \cle{paupérisation} relative, la \cle{faible industrialisation} et l'\cle{extraversion} des échanges. Ce n'est pas un « stade » antérieur au développement, mais une \cle{relation de domination} entretenue par le partage international du travail (thèse de Samir Amin, Walter Rodney).
\end{definition}

\begin{propriete}[Objectifs du Panafricanisme (Philosophie du développement)]
\begin{enumerate}
    \item \textbf{Unité politique :} Création d'États-Unis d'Afrique (Nkrumah) ou Union progressive (Monrovia/OUA) pour peser sur la scène internationale.
    \item \textbf{Intégration économique :} Marché commun, monnaie unique, infrastructures transfrontalières, ZLECAf (Zone de Libre-Échange Continentale Africaine, 2018/2021).
    \item \textbf{Décolonisation mentale/culturelle :} Renaissance africaine, valorisation des langues, histoire, sciences endogènes (Cheikh Anta Diop, Ki-Zerbo).
    \item \textbf{Sécurité collective :} Défense commune, prévention des conflits (PSD de l'UA).
\end{enumerate}
\end{propriete}

\begin{methode}[Rédiger un paragraphe argumenté philosophique]
\begin{enumerate}
    \item \textbf{Accroche / Définition :} Définir le concept clé (unité, développement) ou poser le problème.
    \item \textbf{Thèse :} Affirmer sa position (ex: « L'unité est une condition \emph{nécessaire mais non suffisante}... »).
    \item \textbf{Argument 1 (Économique/Structurel) :} Donnée chiffrée + Analyse (lien avec sous-développement/extraversion).
    \item \textbf{Argument 2 (Politique/Historique) :} Citation auteur + Analyse (volonté politique, souveraineté).
    \item \textbf{Nuance / Limite :} Difficultés (hétérogénéité, souveraineté nationale, gouvernance).
    \item \textbf{Conclusion ouverte :} Perspective concrète (ZLECAf, industrialisation régionale).
\end{enumerate}
\end{methode}

\courssub{Démarche et corrigé commenté}
\begin{exoresolu}[Corrigé commenté de l'activité d'intégration]
\textbf{Consignes rappelées :} Réaliser une carte raisonnée, un diagramme en barres groupées avec calcul de ratio, et un paragraphe argumenté d'environ 250 mots à partir du dossier documentaire.
\tcblower
\textbf{Étape 1 : La carte raisonnée (Attendus)}
\begin{itemize}
    \item \textbf{Zones coloriées :} 
    \begin{itemize}
        \item \cle{CEDEAO} (15 pays, Afrique de l'Ouest) : Siège de la Commission à \textbf{Abuja} (Nigéria). Couleur unique.
        \item \cle{UEMOA} (8 pays, Zone Franc Ouest, subset CEDEAO) : Siège de la Commission à \textbf{Ouagadougou} (Burkina Faso). Hachures distinctes superposées à la CEDEAO.
        \item \cle{CEMAC} (6 pays, Afrique Centrale) : Siège de la Commission à \textbf{Bangui} (RCA) ; Secrétariat permanent/Conférence des Chefs d'État à \textbf{Yaoundé} (Cameroun). Couleur distincte.
        \item \cle{SADC} (16 pays, Afrique Australe) : Siège du Secrétariat à \textbf{Gaborone} (Botswana). Couleur distincte.
    \end{itemize}
    \item \textbf{Flèches Cameroun (proportionnelles aux volumes Tableau 2) :} 
    \begin{itemize}
        \item \emph{Extra-africaines (Grosses flèches sortantes) :} Vers \textbf{Union Européenne} (France, Pays-Bas, Italie, Espagne), Vers \textbf{Chine}, Vers \textbf{Inde/USA}.
        \item \emph{Intra-africaines (Flèches moyennes/fines) :} Vers \textbf{Nigéria} (Pétrole, ciment, commerce informel important), Vers \textbf{Tchad/RCA/Congo} (Produits manufacturés, denrées alimentaires, ciment Cimencam/Dangote).
    \end{itemize}
    \item \textbf{Légende complète :} Titre : « Cameroun : Insertion régionale et extraversion commerciale (2023) ». Auteur : Nom/Groupe. Date. Échelle indicative. Orientation (Nord). Signes conventionnels : Zones de couleur/hachures, Flèches proportionnelles (ex: 1 mm = 1 Md \$).
\end{itemize}

\textbf{Étape 2 : Diagramme en barres et Calcul (Attendus)}
\begin{itemize}
    \item \textbf{Graphique :} Axe des abscisses : 5 catégories (Cameroun, CEDEAO, CEMAC, SADC, Total Continent). Axe des ordonnées : Volume (Mds USD). Deux barres groupées par catégorie : Bleu « Intra-africain », Orange « Extra-africain ». Écart visuel saisissant : barre orange $\gg$ barre bleue pour toutes les catégories.
    \item \textbf{Calcul Cameroun :} Taux de couverture (Extra / Intra) = 10,2 / 1,8 = \textbf{5,67}.
    \item \textbf{Interprétation :} « Le Cameroun échange près de \textbf{5,7 fois plus} avec le reste du monde qu'avec l'Afrique, confirmant une structure d'\cle{extraversion} massive de son économie, caractéristique du \cle{sous-développement} au sens de la théorie de la dépendance. »
\end{itemize}

\textbf{Étape 3 : Paragraphe argumenté (Proposition de corrigé type)}

L'unité africaine apparaît comme une \textbf{condition nécessaire mais non suffisante} du développement du Cameroun et du continent. Au sens du cours, le \cle{développement} exige une transformation structurelle de l'économie pour sortir de la \cle{dépendance} et de l'\cle{extraversion} qui définissent le \cle{sous-développement}. 

Premier argument, \textbf{structurel et économique} : les données du Tableau 2 révèlent que le commerce intra-africain ne représente que \textbf{15,8\%} du total continental et seulement \textbf{15\%} pour le Cameroun (1,8 Md \$ sur 12 Md \$). Le ratio Extra/Intra de 5,67 pour le Cameroun prouve que nos économies sont tournées vers l'extérieur (UE, Chine) pour écouler des matières premières (bois, cacao, pétrole, minerais) et importer des biens finis. Comme l'anticipait \textbf{Nkrumah en 1963} (Extrait 1), sans « union économique », « marché commun » et « monnaie unique », l'indépendance politique reste vide : la ZLECAf, en supprimant les droits de douane sur 90\% des lignes tarifaires, est l'outil technique pour réorienter ces flux, créer des chaînes de valeur régionales (ex: transformation du cacao ivoirien/camerounais au Nigeria) et favoriser l'industrialisation, condition \emph{sine qua non} du développement endogène.

Deuxième argument, \textbf{politique et stratégique} : la \textbf{Charte de l'OUA (1963)} posait le principe de « l'unité et la solidarité » pour « sauvegarder la souveraineté ». Aujourd'hui, face aux grands ensembles mondiaux (UE, USA, Chine, BRICS), aucun État africain seul — pas même le Nigéria ou l'Afrique du Sud — ne pèse assez pour négocier des termes de l'échange équitables ou imposer des normes environnementales et sociales. L'unité (via l'UA, la CEMAC, la CEDEAO) permet une \textbf{mutualisation de la souveraineté} (défense commune, diplomatie unique, infrastructures transfrontalières comme le corridor Douala-Ndjamena-Dakar) qui renforce le pouvoir de négociation du continent.

Toutefois, l'unité ne saurait être une \emph{baguette magique}. L'hétérogénéité des niveaux de développement (IDH 0,394 au Tchad vs 0,713 en Afrique du Sud, Tableau 1), les conflits armés, la faiblesse des institutions et la persistance des barrières non tarifaires (corruption aux frontières, normes techniques divergentes) freinent l'intégration. 

En conclusion, l'unité africaine est le \textbf{cadre politique indispensable} pour inverser la logique extravertie du sous-développement, mais elle doit se concrétiser par une \textbf{intégration productive effective} (ZLECAf, infrastructures énergétiques/numériques, politique industrielle commune) pour devenir le levier réel du développement camerounais et africain.
\end{exoresolu}

\courssub{Bilan de l'activité}
Cette activité d'intégration a permis de :
\begin{itemize}
    \item \textbf{Ancrer la philosophie dans le réel :} Les élèves ont manipulé des données officielles (Banque Mondiale, CNUCED, BEAC) concernant leur pays (Cameroun) et leur région (CEMAC/CEDEAO), transformant des concepts abstraits (extraversion, sous-développement) en réalités chiffrées observables.
    \item \textbf{Articuler les savoir-faire techniques et philosophiques :} La cartographie (spatialisation), la statistique (quantification) et la dissertation (argumentation) ne sont plus des exercices disjoints mais des outils convergents pour penser la même question : le destin de l'Afrique.
    \item \textbf{Mettre en évidence la tension centrale du panafricanisme :} Le corrigé montre que l'unité est un \emph{impératif catégorique} (Nkrumah, Charte OUA) pour briser la dépendance, mais que sa réalisation bute sur des \emph{obstacles structurels} (hétérogénéité, gouvernance, extraversion persistante mesurée à 85\% du commerce).
    \item \textbf{Préparer l'épreuve du Baccalauréat :} L'exercice type « paragraphe argumenté contraint » (thèse, 2 arguments documentés, définition cours, nuance, conclusion ouverte) est un entraînement direct à la dissertation philosophique (Sujet type : « L'unité africaine est-elle la condition du développement ? » ou « Le sous-développement est-il une fatalité ? »).
    \item \textbf{Valoriser l'agentivité :} En conclusion, l'accent mis sur la ZLECAf et l'industrialisation régionale replace les élèves (futurs techniciens, ingénieurs, cadres) comme acteurs du développement, non simples spectateurs de l'histoire.
\end{itemize}

\newpage

\coursec{Méthodes et exercices résolus}

\coursec{Leçon 09 : Panafricanisme et développement de l'Afrique}

\courssub{Le développement : définition et indicateurs}

\begin{definition}[Développement]
Le \cle{développement} est un processus global, durable et multidimensionnel par lequel une société améliore le bien-être de l'ensemble de sa population. Il ne se réduit pas à la croissance économique (augmentation du PIB) : il intègre le progrès social (santé, éducation, protection sociale), culturel (valorisation des identités), politique (démocratie, droits de l'homme) et humain (libertés réelles, dignité). Au sens d'Amartya Sen, le développement est l'\emph{expansion des libertés substantielles} dont jouissent les êtres humains.
\end{definition}

\begin{aretenir}[Indicateurs principaux du développement]
\begin{itemize}
\item \textbf{Économiques :} PIB par habitant (PPA), taux d'industrialisation, part de l'agriculture/industrie/services, taux d'épargne et d'investissement, balance commerciale.
\item \textbf{Sociaux et humains (IDH - PNUD) :} Espérance de vie à la naissance, durée moyenne de scolarisation / durée attendue, revenu national brut par habitant.
\item \textbf{Autres :} Accès à l'eau potable et à l'électricité, taux d'alphabétisation, mortalité infantile/maternelle, indice de développement humain (IDH), indice d'inégalité de genre (IIG), empreinte écologique.
\end{itemize}
\end{aretenir}

\begin{exemplebox}[Le Cameroun face aux indicateurs]
Le Cameroun, pays à revenu intermédiaire de la tranche inférieure (Banque mondiale), affiche un IDH de 0,576 (rang 151/193 en 2022). Malgré une croissance moyenne de 3-4\%/an, les inégalités persistent (indice de Gini $\approx$ 0,44), l'accès à l'électricité reste limité en zone rurale ($\approx$ 20\%) et le taux de transformation locale du cacao, du coton ou du bois est faible ($<$ 30\%). Cela illustre la distinction entre \emph{croissance} (production de richesse) et \emph{développement} (partage et transformation structurelle).
\end{exemplebox}

\courssub{Le sous-développement : réalité et causes}

\begin{definition}[Sous-développement]
Le \cle{sous-développement} n'est pas un « stade » antérieur au développement, mais une \emph{structure} : situation d'un pays ou d'une région caractérisée par la pauvreté de masse, la faiblesse des infrastructures, la dépendance économique (exportation de matières premières, importation de produits finis), le retard technique, la précarité sociale et une insertion subordonnée dans l'économie mondiale. C'est un rapport de dépendance (théorie de la dépendance : A. G. Frank, S. Amin).
\end{definition}

\begin{aretenir}[Causes du sous-développement : double lecture]
\begin{itemize}
\item \textbf{Causes internes (endogènes) :} Mauvaise gouvernance (corruption, népotisme), instabilité politique, faible taux d'épargne nationale, insuffisance de l'investissement productif, croissance démographique non accompagnée, faible capital humain (éducation, santé), absence de politique industrielle volontariste.
\item \textbf{Causes externes (exogènes) :} Héritage colonial (économie d'enclave, infrastructures tournées vers l'export), \cle{échanges inégaux} (détérioration des termes de l'échange, prix des matières premières fixés à l'extérieur), endettement extérieur et conditionnalités des institutions de Bretton Woods, fuite des capitaux (évasion fiscale, transferts illicites), néocolonialisme (contrôle des multinationales, accords de partenariat déséquilibrés).
\end{itemize}
\end{aretenir}

\begin{attention}[Piège conceptuel]
Ne pas confondre \emph{sous-développement} et \emph{non-développement} (terme désuet). Ne pas attribuer le sous-développement à une cause unique (ex. « la colonisation » ou « la corruption »). L'analyse rigoureuse articule \textbf{responsabilités internes} et \textbf{contraintes structurelles externes}.
\end{attention}

\courssub{Historique du panafricanisme}

\begin{aretenir}[Grandes étapes du panafricanisme]
\begin{enumerate}
\item \textbf{Origines diasporiques (fin XIX\textsuperscript{e} -- 1945) :} Réaction à l'esclavage, à la traite et au racisme. Figures : \textbf{W. E. B. Du Bois} (Congrès panafricains de 1900, 1919, 1921, 1923, 1927), \textbf{Marcus Garvey} (UNIA, « L'Afrique aux Africains », retour symbolique). Revendication : dignité, droits civiques, conscience noire.
\item \textbf{Tournant politique : Congrès de Manchester (1945) :} Passage du culturel au politique. Résolution : indépendance immédiate des colonies. Présence de \textbf{Kwame Nkrumah}, \textbf{Jomo Kenyatta}, \textbf{George Padmore}.
\item \textbf{Ère des indépendances et rivalités doctrinales (1957--1963) :}
      \begin{itemize}
      \item Indépendance du \textbf{Ghana} (1957) : Nkrumah proclame « L'indépendance du Ghana n'a de sens que si elle est liée à la libération totale de l'Afrique ».
      \item Deux blocs : \emph{Groupe de Casablanca} (Nkrumah, Sékou Touré, Modibo Keïta : fédéralisme politique immédiat, armée commune) vs \emph{Groupe de Monrovia} (Senghor, Houphouët-Boigny, Tubman : coopération fonctionnelle, souveraineté nationale, gradualisme).
      \end{itemize}
\item \textbf{Institutionnalisation : OUA (Addis-Abeba, 25 mai 1963) :} Charte de l'OUA : souveraineté, non-ingérence, libération des territoires colonisés (Afrique du Sud, Rhodésie, Angola, Mozambique), coopération économique. Compromis monrovien (pas d'État fédéral).
\item \textbf{Relance économique et culturelle (années 1980--2000) :} Plan d'action de Lagos (1980), Traité d'Abuja (1991) créant la \textbf{CAE} (Communauté économique africaine), \textbf{NEPAD} (2001).
\item \textbf{Union africaine (Durban, 2002) :} Remplace l'OUA. Objectifs : intégration accélérée, paix et sécurité (PSC), gouvernance (APRM), citoyenneté africaine, \textbf{ZLECAf} (Zone de libre-échange continentale africaine, entrée en vigueur 2019, démarrage commercial 2021).
\item \textbf{Courant culturel continu :} \emph{Négritude} (Senghor, Césaire, Damas) : revalorisation des valeurs noires. \textbf{Cheikh Anta Diop} : africanité de l'Égypte antique, unité culturelle, langues africaines comme vecteurs de développement.
\end{enumerate}
\end{aretenir}

\begin{savaistu}[Le panafricanisme au Cameroun]
Le Cameroun, membre fondateur de l'OUA (1963), adhère à l'UA (2002). Il appartient à la \textbf{CEMAC} (Communauté économique et monétaire de l'Afrique centrale, 1994) : monnaie commune (FCFA), libre circulation, marché commun imparfait. Le Cameroun accueille le Secrétariat exécutif de la \textbf{COMIFAC} (forêts du bassin du Congo) et participe activement à la ZLECAf. Des figures comme \textbf{Ruben Um Nyobé} (UPC) ou \textbf{Mongo Beti} (écrivain militant) ont porté une vision panafricaniste de libération et de justice sociale.
\end{savaistu}

\courssub{Objectifs du panafricanisme comme philosophie du développement}

\begin{aretenir}[Trois piliers indissociables]
\begin{itemize}
\item \textbf{Objectif politique :} Souveraineté réelle, non-alignement, paix collective (Mécanisme de prévention des conflits), défense des intérêts africains sur la scène internationale (ONU, OMC, G20), citoyenneté africaine (passeport africain, libre circulation).
\item \textbf{Objectif économique :} \cle{Autocentrage} (développement endogène), transformation locale des matières premières (industrialisation), intégration régionale et continentale (ZLECAf : marché unique 1,3 milliard d'habitants, PIB $\approx$ 3 400 milliards \$), pouvoir de négociation face aux blocs (UE, Chine, USA), infrastructures transafricaines (routes, rail, énergie, numérique).
\item \textbf{Objectif culturel et humain :} Réhabilitation de l'histoire et des civilisations africaines (Cheikh Anta Diop), promotion des langues africaines dans l'éducation et l'administration, valorisation des savoirs endogènes, éducation panafricaine, rôle de la diaspora (6\textsuperscript{e} région de l'UA).
\end{itemize}
\end{aretenir}

\courssub{Difficultés et limites du panafricanisme}

\begin{aretenir}[Obstacles majeurs à l'unité africaine]
\begin{itemize}
\item \textbf{Héritage colonial :} Frontières artificielles (Conférence de Berlin 1884-85) coupant les peuples, langues officielles étrangères (français, anglais, portugais) créant des barrières, élites formées hors sol.
\item \textbf{Égoïsmes nationaux et souveraineté jalousée :} Réticence à céder des prérogatives à des institutions supranationales (Parlement panafricain consultatif seulement), multiplicité des \textbf{Communautés économiques régionales (CER)} chevauchantes (CEMAC, CEEAC, UEMOA, CEDEAO, SADC, COMESA, etc.) : « bol de spaghettis » institutionnel.
\item \textbf{Faiblesse économique structurelle :} Faible commerce intra-africain ($\approx$ 15-18\% du total vs 60\% en Europe, 60\% en Asie), dépendance aux matières premières, déficit d'infrastructures transfrontalières, financement de l'UA dépendant de l'extérieur ($\approx$ 60\% du budget).
\item \textbf{Instabilités et gouvernance :} Conflits armés (Sahel, RDC, Soudan, Corne de l'Afrique), coups d'État récurrents, déficit démocratique, corruption, fuite des cerveaux.
\item \textbf{Diversité culturelle et religieuse :} Gestion du pluralisme (laïcité vs droit coutumier/religieux), tensions identitaires instrumentalisées.
\end{itemize}
\end{aretenir}

\courssub{Synthèse : l'unité africaine est-elle la voie du développement ?}

\begin{aretenir}[Réponse nuancée attendue au Bac]
L'unité africaine (\cle{panafricanisme}) est une \textbf{condition nécessaire} du développement : aucun État africain isolé n'a la taille critique (marché, ressources, poids diplomatique) pour s'industrialiser et négocier équitablement (thèse de Nkrumah confirmée par l'économie d'échelle). Mais elle n'est \textbf{pas suffisante} : l'unité politique sans bonne gouvernance, sans capital humain, sans transformation structurelle des économies, risque de produire une « coquille vide ». Le panafricanisme du XXI\textsuperscript{e} siècle est \emph{pragmatique} : intégration économique progressive (ZLECAf), paix et sécurité, gouvernance partagée (APRM), appropriation par la jeunesse. C'est un \emph{horizon régulateur} (Kant) plus qu'un État fédéral immédiat.
\end{aretenir}

% ============================================================
% MÉTHODES ET EXERCICES CORRIGÉS (Bloc collègue revu et corrigé)
% ============================================================

\courssub{Méthode 1 : Appliquer les notions de développement et de sous-développement à une situation concrète}

\begin{methode}[Analyser une situation concrète avec les notions de l'unité]
Pour appliquer les notions de \cle{développement} et de \cle{sous-développement} à une situation de la vie courante (un village, une entreprise, un pays), suivre une démarche en cinq étapes :
\begin{enumerate}
\item \textbf{Lire attentivement la situation} et repérer les faits : niveau de vie, infrastructures, éducation, santé, activités économiques.
\item \textbf{Dégager la notion pertinente} : la situation illustre-t-elle le développement (progrès économique, social et humain durable) ou le sous-développement (pauvreté, dépendance, retard technique) ?
\item \textbf{Identifier les indicateurs} présents dans le texte : revenu par habitant, accès à l'eau potable, scolarisation, espérance de vie, industrialisation.
\item \textbf{Expliquer les causes} : internes (mauvaise gouvernance, faible épargne, corruption) ou externes (dépendance économique, échanges inégaux, endettement).
\item \textbf{Conclure en proposant une voie de solution}, en lien avec le panafricanisme si le sujet s'y prête (unité, solidarité, intégration africaine).
\end{enumerate}
\textbf{Réflexe} : toujours passer du \emph{fait concret} à la \emph{notion philosophique}, puis revenir au fait pour vérifier que l'analyse colle à la réalité.
\end{methode}

\begin{exoresolu}[Analyse d'une situation de vie courante]
\textbf{Énoncé.} Dans un village de la région de l'Extrême-Nord du Cameroun, les habitants cultivent le mil avec des outils traditionnels. Il n'y a ni électricité, ni centre de santé ; les enfants parcourent huit kilomètres pour aller à l'école. Les récoltes sont vendues à bas prix à des commerçants étrangers qui les transforment et les revendent cher. Un jeune diplômé affirme : « Notre village est sous-développé, mais l'union avec les autres villages africains peut nous sauver. »
Appliquez les notions de l'unité pour analyser cette situation et justifier l'affirmation du jeune diplômé.
\tcblower
\textbf{Solution détaillée.}
\begin{enumerate}
\item \textbf{Repérage des faits.} La situation décrit : une agriculture traditionnelle à faible rendement, l'absence d'infrastructures (électricité, santé), l'éloignement de l'école, la vente de matières premières brutes à bas prix et la dépendance vis-à-vis de commerçants étrangers.
\item \textbf{Notion pertinente.} Il s'agit du \cle{sous-développement} : état d'un pays ou d'une communauté caractérisé par la pauvreté, le retard technique, la faiblesse des infrastructures et la dépendance économique.
\item \textbf{Indicateurs mobilisés.} Faible productivité agricole (outils traditionnels), absence d'accès aux services de base (santé, électricité), difficultés de scolarisation, exportation de produits bruts sans valeur ajoutée : autant d'indices qui montrent que le développement ne se réduit pas à la richesse, mais concerne aussi le \emph{développement humain} (éducation, santé, dignité).
\item \textbf{Causes.} \emph{Cause interne} : faiblesse des moyens techniques et de l'organisation locale. \emph{Cause externe} : le système d'échange inégal — le village vend brut et achète cher les produits transformés, ce qui l'enferme dans la dépendance.
\item \textbf{Justification de l'affirmation du diplômé.} Le jeune reprend l'idée centrale du \cle{panafricanisme} : seule, une petite communauté (ou un petit État africain) manque de moyens pour négocier les prix, transformer ses produits et financer ses infrastructures. L'union (coopératives, intégration régionale, marché commun africain) permettrait de mutualiser les ressources, de créer des unités de transformation locales et de renforcer le pouvoir de négociation. C'est la thèse défendue par Kwame Nkrumah : « L'Afrique doit s'unir ou périr. »
\end{enumerate}
\textbf{Conclusion.} La situation illustre le sous-développement comme réalité multidimensionnelle (économique, sociale, humaine) ; l'affirmation du diplômé est justifiée car l'unité africaine, objectif du panafricanisme, apparaît comme une condition du développement véritable.
\end{exoresolu}

\courssub{Méthode 2 : Rédiger une introduction de dissertation sur le panafricanisme}

\begin{methode}[Construire une introduction problématisée]
Une bonne introduction au Baccalauréat technique comporte quatre temps :
\begin{enumerate}
\item \textbf{L'amorce (accroche)} : un fait, une constatation historique ou une citation liée au sujet (ex. l'indépendance du Ghana en 1957, la création de l'OUA en 1963).
\item \textbf{La définition des termes clés} : définir brièvement \cle{panafricanisme} (mouvement visant l'unité et la solidarité des peuples africains et de la diaspora) et \cle{développement} (progrès global, économique, social et humain).
\item \textbf{La problématique} : formuler le problème sous forme de question, en montrant la tension du sujet (l'unité est-elle une condition suffisante du développement ?).
\item \textbf{L'annonce du plan} : présenter les deux ou trois moments de la réflexion, sans les détailler.
\end{enumerate}
\textbf{Réflexe} : l'introduction tient en un seul paragraphe de 8 à 12 lignes ; jamais de plan détaillé ni d'exemples dans l'introduction.
\end{methode}

\begin{exoresolu}[Rédaction d'une introduction]
\textbf{Énoncé.} Sujet de dissertation : « Le panafricanisme peut-il conduire l'Afrique au développement ? » Rédigez une introduction complète.
\tcblower
\textbf{Solution détaillée.}
\begin{enumerate}
\item \textbf{Amorce.} Depuis les indépendances des années 1960, l'Afrique reste, malgré l'immensité de ses richesses naturelles, le continent où la pauvreté est la plus répandue. Face à ce paradoxe, des penseurs et des hommes d'État africains ont fait de l'unité du continent la condition de son émancipation.
\item \textbf{Définitions.} Le \cle{panafricanisme} est le mouvement de pensée et d'action politique qui promeut la solidarité et l'unité de tous les Africains, du continent comme de la diaspora, en vue de leur libération et de leur progrès. Le \cle{développement} désigne quant à lui le progrès durable d'une société dans les domaines économique, social, culturel et humain.
\item \textbf{Problématique.} Mais l'unité politique et culturelle des peuples africains suffit-elle à transformer les structures économiques héritées de la colonisation ? Le panafricanisme est-il une voie réaliste vers le développement, ou un idéal généreux mais irréalisable ?
\item \textbf{Annonce du plan.} Nous verrons d'abord que le panafricanisme constitue une philosophie du développement fondée sur l'unité et la solidarité ; nous examinerons ensuite les difficultés qui limitent sa réalisation ; nous montrerons enfin que l'unité africaine, repensée de manière concrète et progressive, demeure une condition nécessaire du développement du continent.
\end{enumerate}
\textbf{Conclusion de la méthode.} L'introduction ainsi rédigée respecte les quatre temps exigés : elle accroche, définit, problématise et annonce le plan.
\end{exoresolu}

\courssub{Méthode 3 : Construire un paragraphe argumenté avec un auteur}

\begin{methode}[Rédiger un paragraphe de développement]
Chaque paragraphe de la dissertation suit le schéma en cinq étapes :
\begin{enumerate}
\item \textbf{Idée directrice} : une phrase qui énonce clairement la thèse du paragraphe.
\item \textbf{Explication} : développer l'idée en deux ou trois phrases.
\item \textbf{Référence à un auteur} : citer brièvement un penseur du panafricanisme (Nkrumah, Senghor, Cheikh Anta Diop, Du Bois) et sa thèse.
\item \textbf{Exemple concret} : illustrer par un fait camerounais ou africain (création de la CEMAC, de la ZLECAf, panafricanisme culturel).
\item \textbf{Phrase de transition} : conclure le paragraphe et annoncer le suivant.
\end{enumerate}
\textbf{Formule à retenir} : \emph{Idée + Explication + Auteur + Exemple + Transition} (schéma « IEAET »).
\end{methode}

\begin{exoresolu}[Rédaction d'un paragraphe argumenté]
\textbf{Énoncé.} Dans une dissertation sur le thème « L'unité africaine, condition du développement », rédigez un paragraphe montrant que la division de l'Afrique explique sa faiblesse économique.
\tcblower
\textbf{Solution détaillée.}
\begin{enumerate}
\item \textbf{Idée directrice.} La fragmentation politique de l'Afrique en une cinquantaine d'États souvent de petite taille constitue un obstacle majeur à son développement économique.
\item \textbf{Explication.} Divisés, les États africains disposent de marchés trop étroits, négocient isolément le prix de leurs matières premières et restent dépendants des puissances étrangères. La concurrence qu'ils se livrent affaiblit leur position dans les échanges internationaux.
\item \textbf{Référence à un auteur.} Kwame Nkrumah, père de l'indépendance du Ghana et figure du panafricanisme, soutenait dans \emph{L'Afrique doit s'unir} (1963) que seule l'union politique et économique du continent pouvait briser le néocolonialisme et permettre une industrialisation véritable : « L'Afrique doit s'unir ou périr. »
\item \textbf{Exemple concret.} L'expérience des organisations régionales comme la CEMAC, à laquelle appartient le Cameroun, montre que la mise en commun des ressources (monnaie commune, libre circulation) renforce les économies ; à l'échelle continentale, la Zone de libre-échange continentale africaine (ZLECAf, 2019) vise précisément à créer un marché unique de plus d'un milliard de consommateurs.
\item \textbf{Transition.} Si l'unité apparaît ainsi comme une condition du développement, encore faut-il examiner les obstacles qui en freinent la réalisation.
\end{enumerate}
\textbf{Conclusion.} Le paragraphe respecte le schéma IEAET : il affirme une idée, l'explique, l'appuie sur Nkrumah, l'illustre par des exemples africains et prépare la suite du raisonnement.
\end{exoresolu}

\courssub{Méthode 4 : Justifier une réponse ou une conclusion personnelle}

\begin{methode}[Rédiger et justifier une conclusion]
Le savoir-faire « rédiger et justifier une conclusion » exige une démarche rigoureuse :
\begin{enumerate}
\item \textbf{Bilan des thèses} : rappeler en une phrase chacune des positions examinées (thèse et antithèse).
\item \textbf{Prise de position argumentée} : dire clairement ce que l'on pense, en s'appuyant sur les arguments les plus solides du développement — jamais sur une simple opinion.
\item \textbf{Dépassement} : montrer que la réponse nuance ou dépasse l'opposition initiale (ex. l'unité est nécessaire mais non suffisante).
\item \textbf{Ouverture} : poser une question nouvelle liée au sujet (ex. le rôle de la jeunesse africaine, l'intégration économique).
\end{enumerate}
\textbf{Attention} : une conclusion ne contient ni nouvel exemple, ni nouvel auteur ; elle justifie une position déjà construite.
\end{methode}

\begin{exoresolu}[Rédaction d'une conclusion justifiée]
\textbf{Énoncé.} Sujet : « Le panafricanisme peut-il conduire l'Afrique au développement ? » Rédigez la conclusion de la dissertation.
\tcblower
\textbf{Solution détaillée.}
\begin{enumerate}
\item \textbf{Bilan.} Nous avons vu, d'une part, que le panafricanisme constitue une philosophie du développement fondée sur l'unité, la solidarité et la dignité retrouvée des peuples africains ; d'autre part, qu'il se heurte à des difficultés réelles : divisions héritées de la colonisation, égoïsmes nationaux, dépendance économique, faiblesse des institutions panafricaines.
\item \textbf{Prise de position justifiée.} Nous soutenons que le panafricanisme demeure une condition \emph{nécessaire} du développement de l'Afrique, car aucun État africain isolé ne dispose du poids économique suffisant pour transformer ses matières premières et imposer des échanges équitables — l'argument de Nkrumah conserve toute sa force. Mais il n'est pas une condition \emph{suffisante} : l'unité doit s'accompagner de bonne gouvernance, d'éducation et d'industrialisation.
\item \textbf{Dépassement.} Le vrai panafricanisme du XXI\textsuperscript{e} siècle n'est donc pas seulement politique : il est d'abord économique et culturel, comme le montre la ZLECAf, et il se construit progressivement, par étapes régionales.
\item \textbf{Ouverture.} Reste à savoir si la jeunesse africaine, majoritaire sur le continent, saura s'approprier cet idéal pour en faire un projet concret de développement.
\end{enumerate}
\textbf{Conclusion de la méthode.} Cette conclusion fait le bilan, justifie une position nuancée, dépasse l'opposition initiale et ouvre sur une question nouvelle : elle répond pleinement au savoir-faire exigé.
\end{exoresolu}

\courssub{Méthode 5 : Traiter un commentaire de texte sur le panafricanisme}

\begin{methode}[Analyser un texte philosophique]
Pour réussir le commentaire d'un texte de Nkrumah, Senghor ou Cheikh Anta Diop :
\begin{enumerate}
\item \textbf{Présenter le texte} : auteur, œuvre, contexte (ex. luttes pour les indépendances, création de l'OUA en 1963).
\item \textbf{Dégager la thèse} : en une phrase, que soutient l'auteur ? (ex. l'unité africaine est la condition du développement).
\item \textbf{Relever la structure argumentative} : repérer les étapes du raisonnement (constat, argument, conséquence).
\item \textbf{Expliquer les concepts} : définir les termes importants du texte (unité, néocolonialisme, développement, conscience noire).
\item \textbf{Discuter} : apprécier la pertinence de la thèse, en mobilisant les difficultés du panafricanisme et des exemples actuels.
\end{enumerate}
\textbf{Réflexe} : toujours citer une courte phrase du texte pour appuyer chaque explication.
\end{methode}

\begin{exoresolu}[Commentaire guidé d'un extrait]
\textbf{Énoncé.} Texte : « L'indépendance politique n'a de sens que si elle s'accompagne de l'indépendance économique. Aucun de nos États, pris isolément, n'est assez fort pour résister aux pressions extérieures. Notre salut est dans l'union. » (d'après K. Nkrumah). Dégagez la thèse de l'auteur, expliquez son argumentation et discutez sa pertinence pour l'Afrique d'aujourd'hui.
\tcblower
\textbf{Solution détaillée.}
\begin{enumerate}
\item \textbf{Présentation.} Cet extrait est inspiré de la pensée de Kwame Nkrumah (1909-1972), premier président du Ghana indépendant (1957) et principal théoricien du panafricanisme politique, notamment dans \emph{L'Afrique doit s'unir} (1963), l'année de la création de l'OUA.
\item \textbf{Thèse.} L'auteur soutient que l'indépendance politique des États africains demeure illusoire sans unité, car seule l'union garantit l'indépendance économique réelle.
\item \textbf{Structure argumentative.} Le raisonnement comporte trois temps : (a) un \emph{constat} — la souveraineté politique ne suffit pas ; (b) un \emph{argument} — l'isolement rend chaque État vulnérable aux « pressions extérieures » (ce que Nkrumah appelle le néocolonialisme) ; (c) une \emph{conclusion} — « notre salut est dans l'union ».
\item \textbf{Explication des concepts.} L'\emph{indépendance économique} désigne la capacité d'un peuple à maîtriser ses ressources, à transformer ses matières premières et à fixer librement sa politique ; l'\emph{union} renvoie à l'objectif panafricaniste d'intégration politique et économique du continent.
\item \textbf{Discussion.} La thèse garde une forte pertinence : les États africains exportent encore massivement des produits bruts (cacao camerounais, pétrole, cobalt) dont les prix sont fixés ailleurs. La ZLECAf traduit aujourd'hui cette volonté d'union économique. Cependant, l'expérience montre les limites du projet : rivalités entre États, frontières héritées de la colonisation, gouvernances parfois défaillantes. L'unité reste donc un horizon nécessaire, mais qui exige un travail intérieur préalable dans chaque État.
\end{enumerate}
\textbf{Conclusion.} Nkrumah établit rigoureusement que l'unité conditionne le développement ; l'histoire récente confirme sa pertinence tout en révélant les obstacles que le panafricanisme doit encore surmonter.
\end{exoresolu}

\courssub{Méthode 6 : Mobiliser l'historique du panafricanisme dans une copie}

\begin{methode}[Situer les grandes étapes du panafricanisme]
Pour utiliser correctement l'historique du panafricanisme dans un devoir :
\begin{enumerate}
\item \textbf{Retenir les origines} : le mouvement naît dans la diaspora noire (Amérique, Caraïbe) à la fin du XIX\textsuperscript{e} siècle, contre l'esclavage, la colonisation et le racisme (W. E. B. Du Bois, Marcus Garvey).
\item \textbf{Mémoriser les dates clés} : congrès panafricains (Manchester, 1945) ; indépendance du Ghana (1957) ; création de l'OUA à Addis-Abeba (1963) ; transformation en Union africaine (2002).
\item \textbf{Distinguer les courants} : panafricanisme politique (Nkrumah : unité des États), culturel (Senghor, Cheikh Anta Diop : revalorisation des civilisations négro-africaines), économique (intégration régionale, ZLECAf).
\item \textbf{Relier chaque étape à l'objectif de développement} : montrer comment l'unité recherchée vise la libération puis le progrès du continent.
\end{enumerate}
\textbf{Réflexe} : ne citer une date que si elle sert l'argumentation ; l'historique n'est pas un récit mais un outil de preuve.
\end{methode}

\begin{exoresolu}[Utilisation de l'historique dans une réponse argumentée]
\textbf{Énoncé.} Question : « Montrez, en vous appuyant sur l'histoire du panafricanisme, que ce mouvement constitue une philosophie du développement de l'Afrique. » (Réponse rédigée d'une quinzaine de lignes.)
\tcblower
\textbf{Solution détaillée.}
\begin{enumerate}
\item \textbf{Origines.} Né à la fin du XIX\textsuperscript{e} siècle dans la diaspora noire, avec des figures comme W. E. B. Du Bois et Marcus Garvey, le panafricanisme est d'abord un mouvement de résistance : il affirme la dignité des peuples noirs face à l'esclavage, à la colonisation et au racisme. Or la dignité est le premier fondement du développement : un peuple ne peut progresser s'il se méprise lui-même.
\item \textbf{Étape décisive.} Le congrès de Manchester (1945) marque le passage de la revendication culturelle au projet politique : les indépendances deviennent l'objectif explicite. Avec l'indépendance du Ghana en 1957, Nkrumah fait du panafricanisme un programme de développement : unir le continent pour industrialiser, maîtriser les ressources et échapper au néocolonialisme.
\item \textbf{Institutionnalisation.} La création de l'OUA à Addis-Abeba en 1963, puis sa transformation en Union africaine en 2002, traduit la volonté de doter l'Afrique d'institutions communes capables de coordonner son développement. Les organisations régionales (CEMAC pour le Cameroun) et la ZLECAf (2019) prolongent cette dynamique économique.
\item \textbf{Courant culturel.} Senghor et Cheikh Anta Diop ajoutent que le développement suppose la revalorisation des civilisations africaines : développer l'Afrique, c'est aussi lui rendre confiance en son génie propre.
\end{enumerate}
\textbf{Conclusion.} De la diaspora à l'Union africaine, le panafricanisme apparaît ainsi comme une philosophie du développement : il lie inséparablement libération politique, unité économique et fierté culturelle comme conditions du progrès du continent.
\end{exoresolu}

\begin{attention}[Les 5 erreurs qui coûtent des points au Bac]
\begin{enumerate}
\item \textbf{Confondre croissance et développement.} La croissance est une simple augmentation de la richesse produite ; le développement inclut le progrès social et humain (santé, éducation, dignité). Une définition réductrice du développement au seul revenu fait perdre des points dès l'introduction.
\item \textbf{Transformer la dissertation en récit historique.} Aligner des dates (1945, 1957, 1963, 2002) sans les relier à la problématique n'est pas philosopher : chaque fait historique doit servir un argument.
\item \textbf{Citer un auteur sans expliquer sa thèse.} Écrire « Nkrumah a dit que l'Afrique doit s'unir » sans expliquer \emph{pourquoi} (faiblesse des États isolés, néocolonialisme) ni illustrer par un exemple rend la référence inutile.
\item \textbf{Attribuer le sous-développement à une seule cause.} Oublier la distinction entre causes internes (gouvernance, faible industrialisation) et causes externes (échanges inégaux, dépendance, endettement) appauvrit l'analyse ; le correcteur attend les deux dimensions.
\item \textbf{Conclure par une opinion non justifiée.} Affirmer « le panafricanisme est impossible » ou « c'est la seule solution » sans argument, ou introduire un nouvel exemple dans la conclusion, ruine la démarche : la conclusion doit découler des arguments du développement et proposer une position nuancée (condition nécessaire mais non suffisante).
\end{enumerate}
\end{attention}

\newpage

\coursec{Exercices}

\courssub{Je vérifie mes connaissances}

\begin{exercice}[QCM : Concepts fondamentaux]
Parmi les propositions suivantes, une seule est exacte. Recopiez la lettre correspondante et justifiez votre choix en une phrase.

\begin{enumerate}
\item L'Indice de Développement Humain (IDH) combine :
      \begin{itemize}
      \item A. PIB par habitant, taux d'alphabétisation, espérance de vie.
      \item B. Taux de croissance, inflation, chômage.
      \item C. Exportations, importations, balance commerciale.
      \item D. Population active, taux de natalité, taux de mortalité.
      \end{itemize}
\item Le concept de « sous-développement » désigne, pour les théoriciens de la dépendance (A. G. Frank, S. Amin) :
      \begin{itemize}
      \item A. Un retard historique rattrapable par l'industrialisation.
      \item B. Une situation structurelle produite par le système capitaliste mondial.
      \item C. L'absence de ressources naturelles.
      \item D. Une faible densité de population.
      \end{itemize}
\item La Conférence de Bandung (1955) a posé les bases du :
      \begin{itemize}
      \item A. Panafricanisme politique.
      \item B. Non-alignement et solidarité afro-asiatique.
      \item C. Marché commun africain.
      \item D. Organisation de l'Unité Africaine (OUA).
      \end{itemize}
\item Le panafricanisme, comme philosophie du développement, vise prioritairement :
      \begin{itemize}
      \item A. L'unification politique totale de l'Afrique en un seul État fédéral.
      \item B. L'intégration économique, politique et culturelle pour un développement endogène.
      \item C. Le renforcement des liens avec les anciennes puissances coloniales.
      \item D. La création d'une monnaie unique pour tout le continent.
      \end{itemize}
\end{enumerate}
\end{exercice}

\begin{exercice}[Vrai / Faux justifié]
Indiquez si chaque affirmation est vraie ou fausse, puis justifiez brièvement votre réponse en mobilisant une notion du cours.

\begin{enumerate}
\item La croissance économique (augmentation du PIB) est une condition suffisante pour garantir le développement humain.
\item Le panafricanisme est né uniquement après les indépendances des années 1960, sous l'impulsion de Kwame Nkrumah.
\item La Zone de Libre-Échange Continentale Africaine (ZLECAf) est un instrument d'intégration économique panafricaine.
\item La « fuite des cerveaux » (exode des cadres qualifiés) illustre une cause interne du sous-développement sans lien avec le contexte international.
\end{enumerate}
\end{exercice}

\begin{exercice}[Définitions précises]
Définissez rigoureusement les concepts suivants en deux ou trois lignes chacun, en précisant leur portée philosophique dans la leçon.

\begin{enumerate}
\item Développement (au sens d'Amartya Sen / PNUD).
\item Sous-développement (au sens de la théorie de la dépendance).
\item Panafricanisme (comme philosophie du développement).
\item Intégration régionale (ex. : CEMAC, CEDEAO, ZLECAf).
\end{enumerate}
\end{exercice}

\begin{exercice}[Calculs et indicateurs directs]
Le tableau ci-dessous donne des données pour trois pays africains (année de référence).

\begin{center}
\begin{tabularx}{\linewidth}{|l|X|X|X|}\hline
Pays & PIB/habitant (PPA \$US) & Espérance de vie (ans) & Années de scolarité moyennes \\ \hline
Cameroun & 4\,200 & 60 & 6,2 \\
Botswana & 18\,000 & 69 & 9,3 \\
Niger & 1\,200 & 62 & 2,1 \\ \hline
\end{tabularx}
\end{center}

\begin{enumerate}
\item Classez ces trois pays par niveau de développement humain probable (IDH) du plus élevé au plus faible. Justifiez en citant les trois composantes de l'IDH.
\item Le Cameroun a un PIB/habitant supérieur au Niger mais une espérance de vie proche. Que cela suggère-t-il sur la conversion de la richesse en bien-être ?
\end{enumerate}
\end{exercice}

\courssub{Je m'entraîne}

\begin{exercice}[Application : Notions et situations camerounaises]
Pour chacune des trois situations ci-dessous, identifiez \textbf{deux} notions du cours mobilisées et justifiez leur pertinence.

\begin{enumerate}
\item \textbf{Situation 1 :} Le Cameroun exporte 90\% de sa production de cacao sous forme de fèves brutes vers l'Europe, où il est transformé en chocolat à haute valeur ajoutée.
\item \textbf{Situation 2 :} Entrée en vigueur de la ZLECAf (Zone de Libre-Échange Continentale Africaine) : suppression progressive des droits de douane sur 90\% des lignes tarifaires entre États africains.
\item \textbf{Situation 3 :} Chaque année, plus de 3\,000 médecins et ingénieurs formés au Cameroun émigrent vers l'Europe ou l'Amérique du Nord pour y exercer.
\end{enumerate}
\end{exercice}

\begin{exercice}[Analyse de documents statistiques : Croissance vs Développement]
À partir du tableau suivant (données fictives simplifiées, année 2023), rédigez un paragraphe argumenté de 10 lignes comparant croissance économique et développement humain pour ces deux pays.

\begin{center}
\begin{tabularx}{\linewidth}{|l|X|X|X|X|}\hline
Pays & Taux de croissance du PIB (\%) & PIB/hab (PPA \$US) & IDH (0-1) & Part du quintile le plus pauvre dans le revenu national (\%) \\ \hline
Pays A (Pétrolier) & 4,5 & 12\,000 & 0,55 & 4 \\
Pays B (Diversifié) & 2,8 & 5\,500 & 0,68 & 8 \\ \hline
\end{tabularx}
\end{center}

\textit{Consignes :} Utilisez les termes « croissance », « développement », « IDH », « inégalités », « conversion de la richesse ». Concluez sur la pertinence du PIB seul comme indicateur.
\end{exercice}

\begin{exercice}[Débat argumenté : Causes du sous-développement]
Sujet : « Le sous-développement de l'Afrique est-il principalement dû à des causes externes ? »

Construisez un plan dialectique structuré (Thèse / Antithèse / Synthèse) en respectant les consignes :
\begin{itemize}
\item \textbf{Thèse :} Deux arguments majeurs en faveur des causes externes (historiques, structurels), avec un exemple précis par argument.
\item \textbf{Antithèse :} Deux arguments majeurs en faveur des causes internes (gouvernance, démographie, choix de politiques), avec un exemple précis par argument.
\item \textbf{Synthèse :} Une proposition articulée dépassant l'opposition, intégrant la responsabilité partagée et le rôle du panafricanisme comme levier d'autonomie.
\end{itemize}
\end{exercice}

\begin{exercice}[Commentaire guidé : Extrait de Kwame Nkrumah]
Lisez l'extrait suivant (adapté de \emph{Africa Must Unite}, 1963) et répondez aux questions.

\begin{quote}
« L'unité de l'Afrique est une nécessité vitale. Nous avons les ressources, les hommes, le potentiel pour bâtir une Afrique prospère. Mais tant que nous resterons divisés en petits États économiquement non viables, nous serons la proie du néocolonialisme. L'indépendance politique sans indépendance économique est illusoire. L'Organisation de l'Unité Africaine doit devenir un gouvernement continental capable de planifier notre développement industriel et agricole. »
\end{quote}

\begin{enumerate}
\item Dégagez la thèse principale du texte en une phrase.
\item Identifiez et expliquez trois arguments utilisés par Nkrumah pour fonder cette thèse.
\item Quelle conception du panafricanisme (politique, économique, ou les deux ?) ce texte illustre-t-il ? Justifiez.
\item En quoi cette vision répond-elle aux difficultés du sous-développement évoquées dans le cours ?
\end{enumerate}
\end{exercice}

\courssub{Je me prépare au Bac}

\begin{exercice}[Sujet de dissertation type Baccalauréat Technique]
\textbf{Sujet :} « Le panafricanisme est-il la voie du développement de l'Afrique ? »

Vous traiterez ce sujet en respectant la méthode de la dissertation philosophique (introduction avec accroche, définition des termes, problématique, annonce du plan ; développement en trois parties équilibrées ; conclusion avec ouverture). Vous mobiliserez obligatoirement :
\begin{itemize}
\item Les distinctions croissance / développement (Sen, IDH).
\item L'analyse des causes du sous-développement (internes / externes, théorie de la dépendance).
\item L'historique, les objectifs (intégration, autonomie, marché commun, paix) et les difficultés du panafricanisme (souverainetés nationales, disparités économiques, néocolonialisme, gouvernance).
\item Des exemples concrets camerounais et africains (ZLECAf, CEMAC,-UA, industrialisation, formation).
\end{itemize}
\end{exercice}

\begin{exercice}[Commentaire de texte type Baccalauréat Technique]
\textbf{Texte :} Extrait de \emph{Consciencisme} (1964) de Kwame Nkrumah (adapté).

\begin{quote}
« La société africaine traditionnelle était égalitaire et communautaire. Le colonialisme a imposé une structure de classe et une économie d'extraction. Le socialisme africain n'est pas une importation marxiste ; c'est la réaffirmation moderne de nos valeurs communautaires traditionnelles, adaptée aux techniques industrielles. Le panafricanisme est le véhicule politique de ce socialisme : sans unité continentale, aucun État africain ne peut briser le cercle de la dépendance économique. L'indépendance politique n'est que la première étape ; la révolution économique et sociale en est la finalité. »
\end{quote}

\begin{enumerate}
\item \textbf{Dégagez la thèse} de l'auteur en une phrase claire.
\item \textbf{Identifiez et expliquez} les trois arguments principaux qui la soutiennent.
\item \textbf{Problématisez :} Quelle tension le texte révèle-t-il entre héritage culturel et modernité technique ?
\item \textbf{Discutez la portée philosophique :} En quoi cette conception du panafricanisme comme « véhicule » du socialisme africain répond-elle (ou non) aux défis du développement endogène ? Apportez des nuances critiques (ex. : diversité africaine, expériences socialistes passées, rôle du marché).
\end{enumerate}
\end{exercice}

\begin{exercice}[Question à réponse construite / Étude de documents composite]
\textbf{Document 1 :} Extrait du rapport PNUD 2023 — « L'Afrique subsaharienne reste la région au plus faible IDH moyen (0,55), malgré une croissance moyenne du PIB de 3,5\% sur la décennie 2010-2020. Les inégalités de revenus (indice de Gini moyen 0,42) freinent la conversion de la croissance en développement. »

\textbf{Document 2 :} Carte schématique de la ZLECAf (Zone de Libre-Échange Continentale Africaine) — 54 pays signataires, 1,3 milliard de consommateurs potentiels, PIB cumulé 3\,400 milliards \$US. Objectif : +52\% de commerce intra-africain d'ici 2035.

\textbf{Document 3 :} Témoignage d'un entrepreneur camerounais (2024) — « Je produis des jus de fruits locaux à Douala. Mes bouteilles coûtent 30\% plus cher que les jus importés d'Europe (subventionnés, transport maritime peu taxé). La ZLECAf ne m'aide pas encore : les barrières non tarifaires (normes, corruption aux frontières, manque d'infrastructures routières) bloquent mon export vers le Gabon ou le Nigeria. »

\textbf{Questions :}
\begin{enumerate}
\item À partir du Document 1, expliquez pourquoi la croissance économique ne suffit pas à garantir le développement humain en Afrique. Mobilisez les concepts d'IDH, d'inégalités et de conversion de la richesse. (6 pts)
\item Le Document 2 présente la ZLECAf comme un levier majeur. Identifiez l'objectif panafricain visé et expliquez en quoi l'élargissement du marché (échelle continentale) peut favoriser l'industrialisation et la transformation locale (ex. : transformation du cacao, jus de fruits). (6 pts)
\item Le Document 3 révèle des obstacles persistants. En croisant les trois documents, analysez pourquoi l'unité panafricaine (ZLECAf) ne produit pas automatiquement du développement pour les PME locales. Distinguez \textbf{intégration de droit} (traité) et \textbf{intégration de fait} (infrastructures, normes, gouvernance, lutte contre le néocolonialisme commercial). Proposez deux mesures concrètes au niveau national et continental pour lever ces freins. (8 pts)
\end{enumerate}
\end{exercice}

\newpage

\coursec{Corrigés des exercices}

\begin{corrige}[Exercice 1 — QCM : Concepts fondamentaux]
\begin{enumerate}
\item \textbf{Réponse A.} L'IDH, construit par le PNUD (inspiré des travaux d'Amartya Sen), combine trois dimensions : le niveau de vie (PIB par habitant en PPA), l'éducation (scolarisation et alphabétisation) et la santé (espérance de vie à la naissance).
\item \textbf{Réponse B.} Pour André Gunder Frank et Samir Amin (théorie de la dépendance), le sous-développement n'est pas un simple retard mais le produit structurel de l'insertion inégale des périphéries dans le système capitaliste mondial dominé par les « centres ».
\item \textbf{Réponse B.} Bandung (1955) réunit 29 pays d'Asie et d'Afrique et fonde la solidarité afro-asiatique et le non-alignement face aux deux blocs de la Guerre froide ; l'OUA ne sera créée qu'en 1963 à Addis-Abeba.
\item \textbf{Réponse B.} Le panafricanisme vise l'intégration économique, politique et culturelle du continent afin de rendre possible un développement endogène, fondé sur les forces propres de l'Afrique, et non une fusion étatique immédiate ni une simple monnaie unique.
\end{enumerate}
\end{corrige}

\begin{corrige}[Exercice 2 — Vrai / Faux justifié]
\begin{enumerate}
\item \textbf{Faux.} La croissance est une condition \emph{nécessaire} mais non \emph{suffisante} du développement : une hausse du PIB peut coexister avec de fortes inégalités et un IDH faible si la richesse n'est pas convertie en santé, éducation et libertés réelles (Sen).
\item \textbf{Faux.} Le panafricanisme naît dans la diaspora noire dès la fin du XIX\textsuperscript{e} siècle (W. E. B. Du Bois, Marcus Garvey, congrès panafricains de 1900 à 1945) ; Nkrumah et les indépendances en marquent l'étape politique, non la naissance.
\item \textbf{Vrai.} La ZLECAf (entrée en vigueur en 2021) traduit concrètement l'objectif panafricain d'intégration économique en créant un marché continental de 1,3 milliard de consommateurs.
\item \textbf{Faux.} La fuite des cerveaux a bien des causes internes (salaires, conditions de travail), mais elle est inséparable du contexte international : les pays riches attirent les cadres formés aux frais des pays pauvres, ce qui illustre la dépendance et l'échange inégal des ressources humaines.
\end{enumerate}
\end{corrige}

\begin{corrige}[Exercice 3 — Définitions précises]
\begin{enumerate}
\item \textbf{Développement (Sen / PNUD) :} processus d'élargissement des libertés réelles et des capabilités des personnes (vivre longtemps, être instruit, participer à la vie collective), mesuré notamment par l'IDH. Portée philosophique : le développement est une fin humaine, non une simple accumulation de richesses.
\item \textbf{Sous-développement (théorie de la dépendance) :} situation structurelle d'un pays périphérique dont l'économie, spécialisée dans l'exportation de matières premières brutes, est façonnée par et subordonnée aux centres capitalistes. Portée : le sous-développement est produit par le système mondial, non un simple retard naturel.
\item \textbf{Panafricanisme :} doctrine et mouvement visant l'unité politique, économique et culturelle de l'Afrique et de sa diaspora, conçu comme philosophie du développement : seule l'union rendrait possible un développement endogène et l'émancipation vis-à-vis du néocolonialisme.
\item \textbf{Intégration régionale :} processus par lequel des États voisins (CEMAC, CEDEAO, ZLECAf) associent leurs économies (libre-échange, monnaie, politiques communes) pour élargir les marchés, mutualiser les infrastructures et peser davantage dans les échanges mondiaux.
\end{enumerate}
\end{corrige}

\begin{corrige}[Exercice 4 — Calculs et indicateurs directs]
\begin{enumerate}
\item \textbf{Classement probable : Botswana $>$ Cameroun $>$ Niger.} Justification par les trois composantes de l'IDH :
\begin{itemize}
\item \emph{Niveau de vie :} 18\,000 \$ (Botswana) $>$ 4\,200 \$ (Cameroun) $>$ 1\,200 \$ (Niger) ;
\item \emph{Éducation :} 9,3 années de scolarité $>$ 6,2 $>$ 2,1 ;
\item \emph{Santé :} 69 ans $>$ 62 ans $>$ 60 ans.
\end{itemize}
Le Botswana domine sur les trois dimensions ; le Niger est dernier sur le revenu et l'éducation malgré une espérance de vie correcte.
\item Le Cameroun gagne 3,5 fois plus par habitant que le Niger (4\,200 contre 1\,200 \$) pour une espérance de vie seulement \emph{légèrement supérieure} (62 contre 60 ans). Cela suggère que la richesse produite n'est pas efficacement \textbf{convertie en bien-être} : inégalités de revenus, faiblesse des systèmes de santé et d'éducation, mauvaise répartition. Le PIB seul ne mesure donc pas le développement humain, ce qui justifie le recours à l'IDH.
\end{enumerate}
\end{corrige}

\begin{corrige}[Exercice 5 — Application : Notions et situations camerounaises]
\begin{enumerate}
\item \textbf{Situation 1 (cacao brut) :} \emph{Échange inégal} — la valeur ajoutée (transformation en chocolat) est captée à l'étranger, le Cameroun ne vendant que la matière première ; \emph{théorie de la dépendance} — cette spécialisation dans l'export de produits bruts reproduit la position périphérique du pays dans le système mondial.
\item \textbf{Situation 2 (ZLECAf) :} \emph{Intégration régionale} — suppression des droits de douane pour créer un marché continental ; \emph{panafricanisme économique} — réalisation concrète de l'objectif d'unité économique visant un développement endogène (industrialisation, commerce intra-africain).
\item \textbf{Situation 3 (fuite des cerveaux) :} \emph{Cause interne du sous-développement} — perte de capitaux humains formés aux frais de l'État, qui prive le pays de médecins et d'ingénieurs ; \emph{dépendance / échange inégal} — les pays riches bénéficient gratuitement d'une main-d'œuvre qualifiée financée par les pays pauvres, ce qui relie la cause interne au contexte international.
\end{enumerate}
\end{corrige}

\begin{corrige}[Exercice 6 — Analyse de documents : Croissance vs Développement]
\textbf{Paragraphe argumenté attendu (corrigé-type) :}

Le Pays A affiche une croissance forte (4{,}5\,\%) et un PIB par habitant élevé (12\,000 \$US), supérieur au double de celui du Pays B (5\,500 \$US). Pourtant, son IDH (0{,}55) est nettement inférieur à celui du Pays B (0{,}68). Ce paradoxe s'explique par les inégalités : dans le Pays A, le quintile le plus pauvre ne capte que 4\,\% du revenu national, contre 8\,\% dans le Pays B. La richesse pétrolière du Pays A, concentrée entre quelques mains, ne se convertit pas en santé, en éducation ni en libertés réelles pour la majorité : c'est le problème de la \emph{conversion de la richesse} en développement humain, théorisé par Amartya Sen. Le Pays B, à croissance plus modeste (2{,}8\,\%) mais diversifiée et mieux répartie, offre à sa population de meilleures capabilités. On en conclut que le PIB seul est un indicateur insuffisant : il mesure la production de richesses, non leur répartition ni leur traduction en bien-être. Seul un indicateur composite comme l'IDH rend compte du développement entendu comme élargissement des libertés humaines.

\textbf{Points de vigilance :} exiger la citation chiffrée des deux pays, l'usage des cinq termes imposés, et une conclusion explicite sur la pertinence limitée du PIB.
\end{corrige}

\begin{corrige}[Exercice 7 — Débat argumenté : Causes du sous-développement]
\textbf{Thèse — le sous-développement vient principalement de causes externes :}
\begin{itemize}
\item \emph{Argument historique :} la colonisation et la traite ont structuré des économies d'extraction tournées vers l'extérieur. Exemple : les infrastructures camerounaises coloniales relient les mines et plantations aux ports (Douala), non les régions entre elles.
\item \emph{Argument structurel :} l'échange inégal et la baisse des cours des matières premières maintiennent la dépendance (A. G. Frank, S. Amin). Exemple : le Cameroun exporte son cacao brut et importe du chocolat transformé à forte valeur ajoutée.
\end{itemize}
\textbf{Antithèse — le sous-développement vient principalement de causes internes :}
\begin{itemize}
\item \emph{Argument de gouvernance :} corruption, mauvaise gestion et instabilité détournent les ressources du développement. Exemple : des rentes pétrolières africaines mal redistribuées n'élèvent pas l'IDH.
\item \emph{Argument des choix de politiques :} absence de stratégies d'industrialisation et de transformation locale, démographie non maîtrisée. Exemple : faible transformation locale du cacao ou du coton malgré des décennies d'indépendance.
\end{itemize}
\textbf{Synthèse :} les causes externes et internes s'articulent : le système mondial inégal (héritage colonial, néocolonialisme) fixe un cadre contraignant, mais les élites africaines ont une responsabilité propre dans la gouvernance et les choix économiques. Le panafricanisme apparaît alors comme levier d'autonomie : l'intégration continentale (ZLECAf, marchés communs) peut réduire la dépendance externe, tandis que l'exigence de bonne gouvernance répond aux causes internes. Le développement suppose donc une responsabilité partagée et une stratégie endogène collective.
\end{corrige}

\begin{corrige}[Exercice 8 — Commentaire guidé : Kwame Nkrumah, \emph{Africa Must Unite}]
\begin{enumerate}
\item \textbf{Thèse :} seule l'unité politique et économique du continent africain peut rendre l'indépendance réelle et permettre un développement viable, car la division expose l'Afrique au néocolonialisme.
\item \textbf{Trois arguments :}
\begin{itemize}
\item \emph{Argument des ressources :} l'Afrique possède les ressources, les hommes et le potentiel d'une prospérité autonome — l'unité permettrait de les mobiliser ensemble.
\item \emph{Argument de la non-viabilité :} divisés en petits États économiquement non viables, les pays africains restent « la proie du néocolonialisme » — la fragmentation est une faiblesse structurelle.
\item \emph{Argument de l'indépendance économique :} l'indépendance politique sans indépendance économique est « illusoire » — la souveraineté formelle ne suffit pas sans maîtrise de la production.
\end{itemize}
\item Ce texte illustre \textbf{les deux dimensions} du panafricanisme : politique (un « gouvernement continental », une OUA renforcée) et économique (planification du « développement industriel et agricole »). Pour Nkrumah, l'unité politique est l'instrument de la libération économique.
\item Cette vision répond directement au diagnostic du cours sur le sous-développement : si celui-ci résulte de la dépendance et de l'échange inégal, alors l'unité continentale — marché élargi, industrialisation planifiée, transformation locale — est la condition d'un développement endogène brisant le cercle de la dépendance.
\end{enumerate}
\end{corrige}

\begin{corrige}[Exercice 9 — Sujet de dissertation type Baccalauréat Technique]
\textbf{Corrigé-type (plan détaillé et éléments de rédaction).}

\textbf{Introduction.} Accroche : soixante ans après les indépendances, l'Afrique reste la région au plus faible IDH malgré d'immenses ressources. Définitions : le \emph{développement} (Sen) est l'élargissement des libertés réelles, mesuré par l'IDH, et non la seule croissance du PIB ; le \emph{panafricanisme} est la doctrine de l'unité politique, économique et culturelle du continent. Problématique : l'unité africaine est-elle la condition suffisante du développement, ou seulement une condition parmi d'autres ? Annonce du plan : le panafricanisme comme réponse nécessaire (I) ; ses difficultés internes et externes (II) ; une voie nécessaire mais à compléter par la gouvernance et l'endogénéité (III).

\textbf{I. Le panafricanisme, réponse nécessaire au sous-développement.} Le sous-développement n'est pas un retard mais une dépendance structurelle (Frank, Amin) : échange inégal, exportation de matières brutes (cacao camerounais transformé en Europe). La fragmentation en petits États non viables expose au néocolonialisme (Nkrumah, \emph{Africa Must Unite}). L'intégration élargit les marchés, permet l'industrialisation et la transformation locale : la ZLECAf (54 États, 1,3 milliard de consommateurs) en est l'instrument actuel ; la CEMAC en est une étape régionale.

\textbf{II. Les difficultés et limites du panafricanisme.} Attachement des États à leurs souverainetés nationales ; disparités économiques entre pays ; pesanteur du néocolonialisme (accords commerciaux inégaux, concurrence des produits subventionnés) ; faiblesses de gouvernance internes (corruption, fuite des cerveaux, infrastructures insuffisantes). L'échec relatif de l'OUA et la lenteur de l'Union africaine montrent que l'unité proclamée ne produit pas automatiquement le développement.

\textbf{III. Synthèse : une voie nécessaire mais non suffisante.} L'unité est une condition du développement, non sa garantie : elle doit s'accompagner de bonne gouvernance, d'investissement dans l'éducation et la santé (conversion de la richesse en capabilités, Sen), et d'une volonté de transformation locale. Le panafricanisme demeure ainsi moins un État fédéral à construire d'un coup qu'une philosophie pratique : avancer par intégrations régionales concrètes vers un développement endogène.

\textbf{Conclusion.} Le panafricanisme est bien la voie du développement de l'Afrique, à condition d'être compris comme processus exigeant et non comme slogan. Ouverture : la jeunesse africaine, majoritaire, sera-t-elle l'actrice de cette unité ?

\textbf{Barème indicatif (sur 20) :} introduction (accroche, définitions, problématique, plan) : 4 pts ; I : 4 pts ; II : 4 pts ; III : 4 pts ; conclusion et ouverture : 2 pts ; langue, cohérence, exemples camerounais : 2 pts.

\textbf{Points de vigilance :} distinguer explicitement croissance et développement ; citer au moins un auteur (Sen, Nkrumah, Amin) ; mobiliser au moins deux exemples africains ou camerounais ; éviter la paraphrase du cours sans argumentation personnelle.
\end{corrige}

\begin{corrige}[Exercice 10 — Commentaire de texte type Baccalauréat Technique : \emph{Consciencisme}]
\begin{enumerate}
\item \textbf{Thèse :} le socialisme africain, réaffirmation moderne des valeurs communautaires traditionnelles, ne peut triompher de la dépendance économique que par le panafricanisme, l'unité continentale étant la condition de la révolution économique et sociale.
\item \textbf{Trois arguments :}
\begin{itemize}
\item \emph{Argument anthropologique :} la société africaine traditionnelle était « égalitaire et communautaire » ; le socialisme africain en est la continuation, non une importation marxiste — il est donc légitime et endogène.
\item \emph{Argument historique :} le colonialisme a imposé « une structure de classe et une économie d'extraction » ; la dépendance actuelle est un produit historique qu'il faut défaire.
\item \emph{Argument politique :} « sans unité continentale, aucun État africain ne peut briser le cercle de la dépendance » ; l'indépendance politique n'est que « la première étape », la révolution économique en est « la finalité ».
\end{itemize}
\item \textbf{Problématisation :} le texte révèle une tension entre \emph{héritage culturel} (valeurs communautaires traditionnelles) et \emph{modernité technique} (techniques industrielles) : comment moderniser sans imiter l'Occident, et traditionnaliser sans folkloriser ? Nkrumah prétend les articuler, mais la compatibilité du communautarisme villageois avec l'industrie planifiée n'est pas démontrée.
\item \textbf{Discussion critique :} cette conception répond au défi du développement endogène en refusant les modèles importés et en fondant le développement sur les ressources propres du continent. Cependant, des nuances s'imposent : (a) l'Afrique est diverse — parler d'\emph{une} société traditionnelle égalitaire est une simplification ; (b) les expériences socialistes africaines (Ghana de Nkrumah lui-même, Tanzanie) ont souvent échoué économiquement, faute d'efficacité productive et à cause de l'autoritarisme ; (c) le rôle du marché et de l'initiative privée, que Nkrumah néglige, s'est révélé nécessaire à la création de richesses. Le panafricanisme comme « véhicule » reste donc pertinent comme horizon d'unité, mais son contenu économique doit être repensé au-delà du seul socialisme d'État.
\end{enumerate}
\textbf{Barème indicatif (sur 20) :} thèse : 3 pts ; arguments : 6 pts ; problématisation : 4 pts ; discussion nuancée : 5 pts ; langue et méthode : 2 pts.

\textbf{Points de vigilance :} ne pas confondre socialisme africain et marxisme (le texte les distingue explicitement) ; la discussion doit contenir au moins deux nuances critiques argumentées ; citer le texte à l'appui de chaque argument.
\end{corrige}

\begin{corrige}[Exercice 11 — Question à réponse construite / Étude de documents composite]
\begin{enumerate}
\item \textbf{Croissance sans développement (6 pts).} Le Document 1 montre un décalage : une croissance moyenne de 3{,}5\,\% du PIB sur dix ans coexiste avec l'IDH moyen le plus faible du monde (0{,}55). L'IDH mesure le niveau de vie, l'éducation et la santé : si la richesse produite n'améliore pas ces trois dimensions, il n'y a pas développement au sens de Sen. La clé est l'\emph{inégalité} (indice de Gini moyen de 0{,}42) : quand la croissance profite à une minorité, elle ne se \emph{convertit} pas en bien-être pour la majorité. La croissance est donc une condition nécessaire mais non suffisante du développement humain.
\item \textbf{La ZLECAf comme levier (6 pts).} L'objectif panafricain visé est l'\emph{intégration économique continentale} : créer un marché unique de 1,3 milliard de consommateurs (3\,400 milliards \$US de PIB cumulé) et augmenter de 52\,\% le commerce intra-africain. Un marché élargi rend rentables des industries qui ne le seraient pas à l'échelle d'un seul petit État (argument de Nkrumah sur la non-viabilité) : il permet la \emph{transformation locale} (chocolat à partir du cacao camerounais, jus de fruits industrialisés) au lieu de l'exportation brute, captant ainsi la valeur ajoutée sur le continent — c'est la voie d'un développement endogène brisant l'échange inégal.
\item \textbf{Pourquoi l'unité ne suffit pas (8 pts).} Le Document 3 montre que le traité ne suffit pas : il faut distinguer l'\emph{intégration de droit} (la ZLECAf existe juridiquement, les droits de douane baissent) de l'\emph{intégration de fait} (routes, normes harmonisées, fluidité des frontières, compétitivité réelle). L'entrepreneur de Douala subit des barrières non tarifaires (corruption, normes, infrastructures défaillantes) et la concurrence de produits européens subventionnés — forme commerciale du néocolonialisme. Croisé avec le Document 1, cela confirme que ni la croissance ni l'unité formelle ne produisent automatiquement le développement : il faut une conversion effective en capacités productives et humaines. \emph{Mesures possibles} (deux attendues, exemples) : au niveau national, investir dans les infrastructures (corridor Douala–N'Djamena) et soutenir les PME de transformation (financement, énergie) ; au niveau continental, harmoniser les normes et lutter contre la corruption aux frontières (guichets uniques), voire protéger temporairement les industries naissantes africaines face aux produits subventionnés.
\end{enumerate}
\textbf{Barème indicatif :} Q1 : 6 pts (IDH 2, inégalités 2, conversion 2) ; Q2 : 6 pts (objectif panafricain 2, effet de taille de marché 2, exemple de transformation 2) ; Q3 : 8 pts (distinction droit/fait 3, analyse croisée 3, deux mesures 2).

\textbf{Points de vigilance :} exiger l'usage explicite des trois documents dans la Q3 ; la distinction intégration de droit / intégration de fait doit être nommée ; les mesures proposées doivent être concrètes et réparties entre niveau national et continental.
\end{corrige}

\newpage

\coursec{Fiche bilan}

\courssub{Carte mentale de la leçon}
\begin{popfigure}
\begin{tikzpicture}[
  mindmap,
  concept color=popblue,
  level 1 concept/.append style={level distance=4.5cm,sibling angle=60, font=\small\bfseries},
  level 2 concept/.append style={level distance=3cm, font=\scriptsize},
  every node/.style={concept, execute at begin node=\hskip0pt},
  text=white,
  grow cyclic
]
\node[concept, scale=1.1] {Panafricanisme \& Développement}
  child[concept color=popgreen] { node[concept] {1. Développement}
    child { node {Définition : progrès global} }
    child { node {Indicateurs : PIB, IDH, santé, éducation} }
    child { node {Endogène vs Exogène} }
  }
  child[concept color=poporange] { node[concept] {2. Sous-développement}
    child { node {Réalité : pauvreté, dépendance} }
    child { node {Causes internes : gouvernance, corruption} }
    child { node {Causes externes : néocolonialisme, termes échange} }
  }
  child[concept color=poppurple] { node[concept] {3. Historique}
    child { node {Précurseurs : Du Bois, Garvey, Nkrumah} }
    child { node {Conférences : Manchester 1945, Addis-Abeba 1963} }
    child { node {OUA (1963) $\to$ UA (2002)} }
  }
  child[concept color=poppink] { node[concept] {4. Objectifs}
    child { node {Unité politique \& économique} }
    child { node {Libération totale (indépendance 2)} }
    child { node {Solidarité, paix, développement endogène} }
    child { node {Renaissance culturelle africaine} }
  }
  child[concept color=popgold] { node[concept] {5. Difficultés}
    child { node {Souverainetés nationales jalousées} }
    child { node {Hétérogénéité (langues, économies)} }
    child { node {Dépendance financière ext.} }
    child { node {Conflits, instabilité, coups d'État} }
  }
  child[concept color=popdark] { node[concept] {6. Synthèse}
    child { node {Unité = condition nécessaire} }
    child { node {Pas suffisante sans bonne gouvernance} }
    child { node {ZLECAf, Agenda 2063 : leviers actuels} }
  };
\end{tikzpicture}
\legende{Carte mentale récapitulative : concepts, histoire, objectifs, obstacles et enjeux du panafricanisme comme philosophie du développement.}
\end{popfigure}

\courssub{Bilan des savoirs essentiels}

\textbf{1. Définitions clés}
\begin{itemize}
  \item \cle{Développement} : processus global (économique, social, culturel, politique) visant l'amélioration durable du bien-être des populations. Il ne se réduit pas à la \emph{croissance} (augmentation du PIB) : il suppose des transformations structurelles (industrialisation, éducation, santé, justice sociale). Il est mesuré par l'\cle{IDH} (Indice de Développement Humain : santé, éducation, niveau de vie) et non plus seulement par le PIB par habitant.
  \item \cle{Sous-développement} : situation structurelle (et non simple « retard » ou stade transitoire) caractérisée par la pauvreté de masse, la dépendance économique, la faiblesse des infrastructures et du capital humain. Distinction capitale : le \emph{sous-développement} résulte d'une relation de domination (il est \emph{produit} par l'échange inégal), tandis que le \emph{non-développement} désignerait un retard encore rattrapable en dehors de toute domination.
  \item \cle{Panafricanisme} : idéologie et mouvement politique prônant l'unité, la solidarité et la coopération entre les peuples africains et la diaspora, en vue de leur libération, de leur dignité et de leur développement autonome.
\end{itemize}

\begin{attention}[Piège fréquent à l'examen]
Ne confonds pas \emph{croissance économique} et \emph{développement}. Un pays peut connaître une forte croissance (exportation de pétrole ou de cacao) sans développement réel, si les richesses ne profitent pas aux populations (chômage, écoles et hôpitaux insuffisants). Le correcteur attend cette nuance dès l'introduction de ta dissertation.
\end{attention}

\textbf{2. Repères historiques et auteurs majeurs (à citer)}
\begin{center}
\begin{tabularx}{\linewidth}{|l|X|X|}
\hline
\rowcolor{popblueL}
\textbf{Période / Étape} & \textbf{Acteurs / Événements clés} & \textbf{Apport conceptuel} \\ \hline
Années 1900-1945 & W.E.B. Du Bois, Marcus Garvey ; congrès panafricains (Paris 1919, Londres 1921, New York 1927, Manchester 1945) & Conscience raciale, « l'Afrique aux Africains », droit des peuples à l'autodétermination. \\ \hline
Années 1958-1963 & Kwame Nkrumah (Ghana), Sékou Touré (Guinée), Modibo Keïta (Mali) — Groupe de Casablanca ; Félix Houphouët-Boigny (Côte d'Ivoire), Léopold Sédar Senghor (Sénégal) — Groupe de Monrovia. & Débat fondateur : \textbf{unité fédérale immédiate} (Casablanca) contre \textbf{unité progressive par la coopération} (Monrovia). \\ \hline
25 mai 1963 & Création de l'\cle{OUA} (Organisation de l'Unité Africaine) à Addis-Abeba (Éthiopie). Charte : souveraineté des États, non-ingérence, libération du continent. & Compromis : unité politique différée, priorité à la décolonisation (Comité de libération). \\ \hline
1991 / 2002 & Traité d'Abuja (Communauté Économique Africaine) $\to$ \cle{UA} (Union Africaine), lancée à Durban (Afrique du Sud). & Passage de la « non-ingérence » à la « non-indifférence » ; objectifs : intégration économique, paix, bonne gouvernance. \\ \hline
2018 - aujourd'hui & \cle{ZLECAf} (Zone de Libre-Échange Continentale Africaine, accord signé à Kigali en 2018), Agenda 2063. & Outil économique concret de l'intégration : marché unique d'environ 1,3 milliard d'habitants. \\ \hline
\end{tabularx}
\end{center}

\begin{aretenir}[Repères à mémoriser]
\begin{itemize}
  \item \textbf{1945} : congrès de Manchester, tournant vers l'exigence d'indépendance.
  \item \textbf{1963} : OUA à Addis-Abeba — principe de \emph{non-ingérence}.
  \item \textbf{2002} : UA à Durban — principe de \emph{non-indifférence} (l'Union peut intervenir en cas de crimes graves).
  \item \textbf{2018} : ZLECAf — le volet économique du panafricanisme contemporain.
\end{itemize}
\end{aretenir}

\textbf{3. Objectifs du panafricanisme comme philosophie du développement}
\begin{enumerate}
  \item \textbf{Libération totale} : indépendance politique (acquise) \emph{plus} indépendance économique (souveraineté monétaire, industrielle, alimentaire) — c'est ce que Nkrumah appelait l'achèvement de la décolonisation.
  \item \textbf{Unité africaine} : marché commun, libre circulation des personnes, défense commune, institutions continentales (Parlement panafricain, Cour africaine de justice).
  \item \textbf{Développement endogène} : valorisation des ressources et des savoirs locaux, priorité à l'humain, dans l'esprit de la philosophie \emph{Ubuntu} (« Je suis parce que nous sommes »).
  \item \textbf{Renaissance culturelle} : réhabilitation de l'histoire, des langues et des valeurs africaines contre l'aliénation culturelle (Cheikh Anta Diop : l'unité culturelle comme fondement de l'unité politique).
  \item \textbf{Paix et sécurité} : mécanismes africains de prévention et de règlement des conflits (UA, CEMAC, CEDEAO).
\end{enumerate}

\textbf{4. Difficultés et limites (obstacles à l'unité)}
\begin{itemize}
  \item \textbf{Politiques} : jalousie des souverainetés nationales, instabilité (coups d'État au Sahel), absence de leadership continental fort, non-application des protocoles (libre circulation).
  \item \textbf{Économiques} : économies extraverties (exportation de matières premières, importation de produits finis), faible commerce intra-africain (environ 15\% des échanges du continent), dépendance à l'aide et à la dette extérieures, infrastructures transfrontalières insuffisantes.
  \item \textbf{Sociales et culturelles} : diversité linguistique (plus de 2000 langues), frontières artificielles héritées du partage colonial (conférence de Berlin, 1884-1885), népotisme, tribalisme, fuite des cerveaux.
  \item \textbf{Institutionnelles} : UA sous-financée et dépendante de contributeurs extérieurs, lourdeurs bureaucratiques, décisions peu contraignantes pour les États membres.
\end{itemize}

\begin{exemplebox}[Illustrations camerounaises]
\begin{itemize}
  \item Le Cameroun est membre fondateur de l'OUA (1963) et a ratifié la \cle{ZLECAf}, participant ainsi au marché unique africain.
  \item Le Cameroun joue un rôle d'intégration \emph{sous-régionale} au sein de la \cle{CEMAC} (Communauté Économique et Monétaire de l'Afrique Centrale), dont la Banque des États de l'Afrique Centrale (BEAC) a son siège à Yaoundé.
  \item Les difficultés internes (crise dans les régions du Nord-Ouest et du Sud-Ouest, insécurité dans l'Extrême-Nord) montrent qu'aucune intégration continentale ne peut réussir sans paix et gouvernance à l'échelle nationale.
\end{itemize}
\end{exemplebox}

\textbf{5. Méthode express : dissertation et commentaire sur ce thème}
\begin{methode}[Structurer une réponse « Panafricanisme et développement »]
\begin{enumerate}
  \item \textbf{Analyser le sujet} : identifie la tension (unité contre souveraineté, idéal contre réel, politique contre économique). Définis les termes clés (développement \emph{endogène} contre \emph{exogène}).
  \item \textbf{Problématiser} : par exemple « L'unité africaine est-elle une condition \emph{suffisante} du développement ou seulement \emph{nécessaire} ? » ou « Le panafricanisme est-il encore opératoire face aux nationalismes ? »
  \item \textbf{Construire un plan dialectique} :
    \begin{itemize}
      \item \textbf{Thèse} : le panafricanisme fonde la puissance collective (Nkrumah : « Divisés, nous sommes faibles ; unis, l'Afrique devient une force »). L'unité est le cadre \emph{sine qua non} : marché élargi, diplomatie commune, sécurité collective.
      \item \textbf{Antithèse} : l'unité formelle masque des divergences réelles. Le développement dépend d'abord de la gouvernance \emph{interne} (État de droit, éducation, lutte contre la corruption). Des États comme le Botswana ou le Rwanda progressent sans unité fédérale préalable ; l'UA peine à sanctionner les putschistes.
      \item \textbf{Synthèse} : l'unité est un \emph{levier structurel} (ZLECAf, infrastructures, voix commune à l'ONU), mais son efficacité exige une \emph{refondation de l'État} (contrat social, reddition des comptes). Il faut articuler un panafricanisme \emph{par le bas} (peuples, jeunesse, culture) et \emph{par le haut} (institutions).
    \end{itemize}
  \item \textbf{Illustrer avec le Cameroun} : membre fondateur de l'OUA, rôle dans la CEMAC, ratification de la ZLECAf, mais aussi défis internes (crise anglophone) qui freinent l'intégration.
  \item \textbf{Conclure} : ouverture sur l'Agenda 2063, « L'Afrique que nous voulons » — et sur la responsabilité de la jeunesse technique dans la transformation du continent.
\end{enumerate}
\end{methode}

\courssub{Checklist avant le Bac}
\begin{aretenir}[Checklist avant le Bac — Leçon 9]
\begin{itemize}
  \item $\square$ Je sais définir le \textbf{développement} (processus global, mesuré par l'IDH) et le distinguer de la simple croissance (PIB).
  \item $\square$ Je sais expliquer le \textbf{sous-développement} comme relation de dépendance (causes internes et externes) et non comme un simple « retard ».
  \item $\square$ Je connais les \textbf{3 grandes phases} du panafricanisme : diaspora et précurseurs (Du Bois, Garvey) $\to$ indépendances et débat Casablanca contre Monrovia (Nkrumah contre Houphouët-Boigny) $\to$ OUA, UA, ZLECAf.
  \item $\square$ Je sais citer \textbf{3 objectifs majeurs} : unité politique et économique, développement endogène, renaissance culturelle.
  \item $\square$ Je peux lister \textbf{4 obstacles majeurs} : souverainetés nationales, économies extraverties, instabilités politiques, faiblesse institutionnelle de l'UA.
  \item $\square$ Je maîtrise la \textbf{nuance thèse/antithèse} : l'unité est une condition nécessaire mais non suffisante (la gouvernance locale reste déterminante).
  \item $\square$ Je possède \textbf{2 exemples camerounais concrets} : membre fondateur de l'OUA, siège de la BEAC à Yaoundé, ratification de la ZLECAf, crise du Nord-Ouest et du Sud-Ouest comme frein à l'intégration.
  \item $\square$ Je sais mobiliser \textbf{une référence d'auteur} : Nkrumah (\emph{Africa Must Unite}), Cheikh Anta Diop (fondement culturel de l'unité), ou l'Agenda 2063 de l'UA.
  \item $\square$ Je sais rédiger une \textbf{problématique} pertinente (ex. : « L'unité africaine est-elle le remède au sous-développement ? »).
  \item $\square$ Je connais le \textbf{plan type en 3 parties} (thèse : nécessité de l'unité / antithèse : primauté du national / synthèse : articulation des échelles locale et continentale).
  \item $\square$ Je distingue \cle{OUA} (1963, non-ingérence, libération) et \cle{UA} (2002, non-indifférence, intégration, droits de l'homme).
  \item $\square$ Je sais expliquer la \cle{ZLECAf} comme le bras économique actuel du panafricanisme (marché unique, transformation structurelle).
\end{itemize}
\end{aretenir}

\findecours

\end{document}
